    \chapter{Set Theory and Logic}

    \relatedcourses{(MATH202) Set Theory}

    \relatedbooks{Introduction to Set Theory (Karel Hrbacek and Thomas Jech)}

    \begin{definition}[Set and Element]
        A set is a collection of objects, called its elements. We write
        $x\in A$ if $x$ is an element of $A$, and $x\notin A$ otherwise.
        The relation $A\subseteq B$ means
        \[
            \forall x,\qquad x\in A\Rightarrow x\in B.
        \]
    \end{definition}

    \begin{formula}[Basic Set Operations]
        For sets $A$ and $B$ inside an ambient set $U$, we use
        \[
            A\cup B=\{x\in U\mid (x\in A)\lor(x\in B)\},
        \]
        \[
            A\cap B=\{x\in U\mid (x\in A)\land(x\in B)\},
        \]
        and
        \[
            A\setminus B=\{x\in A\mid x\notin B\}.
        \]
    \end{formula}

    \begin{formula}[Set Algebra]
        For subsets $A,B,C\subseteq U$, the following identities are often useful:
        \[
            A\cup\varnothing=A,
            \qquad
            A\cap U=A,
            \qquad
            A\cup A=A,
            \qquad
            A\cap A=A.
        \]
        The distributive laws are
        \[
            A\cap(B\cup C)=(A\cap B)\cup(A\cap C),
        \]
        \[
            A\cup(B\cap C)=(A\cup B)\cap(A\cup C).
        \]
        The absorption laws are
        \[
            A\cup(A\cap B)=A,
            \qquad
            A\cap(A\cup B)=A.
        \]
    \end{formula}

    \begin{formula}[Basic Logical Equivalences]
        The most useful logical equivalences are
        \[
            (P\Rightarrow Q)\Longleftrightarrow(\neg Q\Rightarrow\neg P),
        \]
        \[
            \neg(P\Rightarrow Q)\Longleftrightarrow P\land\neg Q,
        \]
        \[
            \neg(P\land Q)\Longleftrightarrow(\neg P)\lor(\neg Q),
            \qquad
            \neg(P\lor Q)\Longleftrightarrow(\neg P)\land(\neg Q),
        \]
        and
        \[
            \neg(\forall x\in A,\ P(x))
            \Longleftrightarrow
            \exists x\in A\ \text{s.t.}\ \neg P(x),
        \]
        \[
            \neg(\exists x\in A\ \text{s.t.}\ P(x))
            \Longleftrightarrow
            \forall x\in A,\ \neg P(x).
        \]
    \end{formula}

    \begin{strategy}
        To prove $P\Rightarrow Q$, it is often enough to prove the contrapositive
        $\neg Q\Rightarrow\neg P$. To disprove a universal statement, find a
        counterexample.
    \end{strategy}

    \begin{formula}[De Morgan's Laws]
        For subsets $A,B\subseteq U$, we have
        \[
            (A\cup B)^c=A^c\cap B^c,
            \qquad
            (A\cap B)^c=A^c\cup B^c.
        \]
        More generally,
        \[
            \left(\bigcup_{i\in I}A_i\right)^c
            =
            \bigcap_{i\in I}A_i^c,
            \qquad
            \left(\bigcap_{i\in I}A_i\right)^c
            =
            \bigcup_{i\in I}A_i^c.
        \]
    \end{formula}

    \begin{definition}[Power Set]
        The power set of a set $A$ is
        \[
            \Pow(A)=\{B\mid B\subseteq A\}.
        \]
        If $|A|=n<\infty$, then
        \[
            |\Pow(A)|=2^n.
        \]
    \end{definition}

    \begin{definition}[Function]
        A function $f\colon X\to Y$ assigns to each element $x\in X$ a unique element
        $f(x)\in Y$. The set $X$ is the domain, and $Y$ is the codomain. The
        notation $x\mapsto f(x)$ describes the rule on elements.
    \end{definition}

    \begin{definition}[Injection, Surjection, and Bijection]
        A function $f\colon X\to Y$ is injective if distinct elements of $X$ have
        distinct images. Equivalently,
        \[
            \forall x_1,x_2\in X,
            \qquad
            f(x_1)=f(x_2)\Rightarrow x_1=x_2.
        \]
        Once injectivity has been established, one may write
        $f\colon X\hookrightarrow Y$.

        The function is surjective if every element of $Y$ has a preimage:
        \[
            \forall y\in Y,
            \qquad
            \exists x\in X\ \text{s.t.}\ f(x)=y.
        \]
        Once surjectivity has been established, one may write
        $f\colon X\twoheadrightarrow Y$. The function is bijective if it is both
        injective and surjective.
    \end{definition}

    \begin{proposition}[Finite Injective-Surjective Criterion]
        If $X$ and $Y$ are finite sets with $|X|=|Y|$, then for a function
        $f\colon X\to Y$, the following are equivalent:
        \[
            f\text{ is injective}
            \quad\Longleftrightarrow\quad
            f\text{ is surjective}
            \quad\Longleftrightarrow\quad
            f\text{ is bijective}.
        \]
    \end{proposition}

    \begin{proposition}[Inverse Image Preserves Set Operations]
        If $f\colon X\to Y$ and $A,B\subseteq Y$, then
        \[
            f^{-1}(A\cup B)=f^{-1}(A)\cup f^{-1}(B),
        \]
        \[
            f^{-1}(A\cap B)=f^{-1}(A)\cap f^{-1}(B),
        \]
        and
        \[
            f^{-1}(Y\setminus A)=X\setminus f^{-1}(A).
        \]
    \end{proposition}

    \begin{definition}[Relation and Equivalence Relation]
        A binary relation on a set $A$ is a subset $R\subseteq A\times A$.
        An equivalence relation $\sim$ is reflexive, symmetric, and transitive;
        explicitly,
        \[
            \forall x\in A,
            \qquad x\sim x,
        \]
        \[
            \forall x,y\in A,
            \qquad x\sim y\Rightarrow y\sim x,
        \]
        and
        \[
            \forall x,y,z\in A,
            \qquad
            (x\sim y\land y\sim z)\Rightarrow x\sim z.
        \]
        The equivalence class of $a\in A$ is
        \[
            [a]=\{x\in A\mid x\sim a\}.
        \]
    \end{definition}

    \begin{proposition}[Equivalence Relations and Partitions]
        Equivalence relations on a set $A$ are the same structure as partitions
        of $A$: each equivalence relation partitions $A$ into equivalence
        classes, and each partition of $A$ determines an equivalence relation.
    \end{proposition}

    \begin{definition}[Cardinality and Equinumerosity]
        Two sets $A$ and $B$ have the same cardinality if there exists a
        bijection between them. We write
        \[
            |A|=|B|
            \qquad\text{or}\qquad
            A\xrightarrow{\sim}B.
        \]
        Thus cardinality compares the sizes of finite and infinite sets through
        one-to-one correspondences rather than through geometric containment.
    \end{definition}

    \begin{definition}[Dedekind-Infinite Set]
        A set $A$ is Dedekind-infinite if there exists a proper subset
        $B\subsetneq A$ s.t.
        \[
            |A|=|B|.
        \]
        In ZFC, a set is Dedekind-infinite if and only if it is infinite. Without
        the Axiom of Choice, this equivalence cannot be proved in full
        generality.
    \end{definition}

    \begin{definition}[Countability]
        A set $A$ is countable if it is finite or if there exists a bijection
        $f\colon A\to\N$. In the countably infinite case,
        \[
            \exists f\colon A\xrightarrow{\sim}\N,
            \qquad
            |A|=|\N|=\aleph_0.
        \]
        A set is uncountable if it is not countable.
    \end{definition}

    \begin{theorem}[Basic Cardinalities]
        The standard countably infinite sets satisfy
        \[
            |\N|=|2\N|=|\Z|=|\Q|=\aleph_0,
        \]
        whereas
        \[
            |(0,1)|=|\R|=\mathfrak{c}=2^{\aleph_0}>\aleph_0.
        \]
        In particular, a proper subset of an infinite set may have the same
        cardinality as the whole set, while the real numbers cannot be listed by
        the natural numbers.
    \end{theorem}


    \begin{theorem}[Cantor's Theorem]
        For every set $A$, there is no surjection from $A$ onto $\Pow(A)$.
        Equivalently,
        \[
            |A|<|\Pow(A)|.
        \]
    \end{theorem}

    \begin{definition}[Aleph Numbers and the Continuum]
        The symbol $\aleph_0$ denotes the cardinality of a countably infinite
        set, and $\aleph_1$ denotes the least uncountable cardinal. The
        cardinality of the continuum is
        \[
            \mathfrak{c}=|\R|=2^{\aleph_0}.
        \]
        The Continuum Hypothesis is the assertion
        \[
            \mathfrak{c}=\aleph_1.
        \]
    \end{definition}

    \begin{theorem}[Schr\"oder--Bernstein Theorem; Cantor--Bernstein Theorem]
        If there exists an injection $f\colon A\hookrightarrow B$ and an injection
        $g\colon B\hookrightarrow A$, then there exists a bijection
        \[
            h\colon A\xrightarrow{\sim}B.
        \]
        Thus two sets have the same cardinality once each can be embedded into
        the other.
    \end{theorem}

    \begin{algorithm}[Mathematical Induction]
        To prove a statement $P(n)$ for all integers $n\geq n_0$:
        \begin{bookitems}
            \item Prove the base case $P(n_0)$.
            \item Prove the induction step: for every $k\geq n_0$, show that
                  $P(k)$ implies $P(k+1)$.
        \end{bookitems}
        Then $P(n)$ holds for all integers $n\geq n_0$.
    \end{algorithm}

    \begin{theorem}[Well-Ordering Principle]
        Every nonempty subset of $\N$ has a least element.
    \end{theorem}

    \begin{fact}[Russell's Paradox]
        The expression ``the set of all sets that do not contain themselves''
        leads to a contradiction in naive set theory. This is one reason modern
        set theory is formulated axiomatically.
    \end{fact}

    \begin{fact}[ZF and ZFC]
        ZF denotes Zermelo--Fraenkel set theory, and ZFC denotes ZF together
        with the Axiom of Choice. The axioms restrict set formation sufficiently
        to avoid Russell's paradox while supporting the standard foundations of
        modern mathematics.
    \end{fact}

    \begin{fact}[Axiom of Choice and Zorn's Lemma]
        The Axiom of Choice, Zorn's Lemma, and the Well-Ordering Theorem are
        equivalent over the usual Zermelo--Fraenkel axioms. They are standard
        foundational tools in advanced algebra, analysis, and topology.
    \end{fact}

    \chapter{Calculus and Analysis}

    \relatedcourses{(MATH101) Calculus I, (MATH102) Calculus II,
        (MATH210) Applied Complex Variables, (MATH311) Analysis I,
        (MATH312) Analysis II}

    \relatedbooks{Calculus With Applications (Peter D. Lax and Maria Shea Terrell),
        Multivariable Calculus with Applications (Peter D. Lax and Maria Shea Terrell),
        Principles of Mathematical Analysis (Walter Rudin),
        Complex Variables and Applications (James Ward Brown and Ruel V. Churchill)}

    \section*{One-Variable Calculus}

    \begin{definition}[Sequence, Series, and Partial Sum]
        A sequence is a function $a\colon \N\to\R$ or $a\colon \N\to\C$, usually written
        $(a_n)$. A series is the formal sum
        \[
            \sum_{n=1}^{\infty}a_n,
        \]
        whose $N$th partial sum is
        \[
            S_N=\sum_{n=1}^{N}a_n.
        \]
        The series converges to $S$ exactly when $S_N\to S$ as $N\to\infty$;
        otherwise it diverges. It converges absolutely if
        \[
            \sum_{n=1}^{\infty}|a_n|
        \]
        converges, and it converges conditionally if it converges but does not
        converge absolutely. Thus every convergence test for a series is a
        criterion for convergence of its sequence of partial sums.
    \end{definition}

    \begin{formula}[Arithmetic and Geometric Progressions]
        For an arithmetic progression with first term $a$ and common difference
        $d$, we have
        \[
            a_n=a+(n-1)d,
            \qquad
            \sum_{k=1}^{n}a_k=\frac{n}{2}(a_1+a_n).
        \]
        For a geometric progression with first term $a$ and common ratio $r$, we
        have
        \[
            a_n=ar^{n-1},
            \qquad
            \sum_{k=0}^{n-1}ar^k=a\frac{1-r^n}{1-r}\quad(r\neq 1).
        \]
        If $|r|<1$, then
        \[
            \sum_{k=0}^{\infty}ar^k=\frac{a}{1-r}.
        \]
    \end{formula}

    \begin{formula}[Sums of Powers; Faulhaber's Formulas]
        The first power sums are
        \[
            \sum_{k=1}^{n}k
            =
            \frac{n(n+1)}{2},
        \]
        \[
            \sum_{k=1}^{n}k^2
            =
            \frac{n(n+1)(2n+1)}{6},
            \qquad
            \sum_{k=1}^{n}k^3
            =
            \left(\frac{n(n+1)}{2}\right)^2,
        \]
        and
        \[
            \sum_{k=1}^{n}k^4
            =
            \frac{n(n+1)(2n+1)(3n^2+3n-1)}{30}.
        \]
        More generally, for each fixed nonnegative integer $m$,
        \[
            \sum_{k=1}^{n}k^m
        \]
        is a polynomial in $n$ of degree $m+1$ with leading coefficient
        $1/(m+1)$.
    \end{formula}

    \begin{formula}[Finite Differences; Newton's Forward Difference Formula]
        For a sequence $(a_n)$, define the forward difference by
        \[
            \Delta a_n=a_{n+1}-a_n.
        \]
        If $a_n=P(n)$ for a polynomial $P$ of degree $d$ with leading
        coefficient $c$, then
        \[
            \Delta^d a_n=d!c,
            \qquad
            \Delta^{d+1}a_n=0.
        \]
        Conversely, a sequence with constant $d$th differences is polynomial
        in $n$ of degree at most $d$. For every polynomial $P$ of degree at most
        $d$,
        \[
            P(n)
            =
            \sum_{k=0}^{d}\binom{n}{k}\Delta^kP(0).
        \]
    \end{formula}

    \begin{strategy}
        A table of successive differences is one of the most efficient methods to detect a
        hidden polynomial sequence. Constant first, second, and third differences
        indicate linear, quadratic, and cubic behavior, respectively.
    \end{strategy}

    \begin{formula}[Linear Recurrences with Constant Coefficients]
        Consider
        \[
            a_n=c_1a_{n-1}+c_2a_{n-2}+\cdots+c_ma_{n-m}.
        \]
        Its characteristic polynomial is
        \[
            r^m-c_1r^{m-1}-c_2r^{m-2}-\cdots-c_m.
        \]
        If the roots $r_1,\dots,r_m$ are distinct, then
        \[
            a_n=A_1r_1^n+\cdots+A_mr_m^n.
        \]
        A root $r$ of multiplicity $q$ contributes terms
        \[
            r^n,\ nr^n,\ \dots,\ n^{q-1}r^n.
        \]
        The constants are determined from the initial values.
    \end{formula}

    \begin{formula}[First-Order Linear Recurrence]
        If
        \[
            a_n=ra_{n-1}+b_n
            \qquad(n\geq 1),
        \]
        then iteration gives
        \[
            a_n
            =
            r^na_0
            +
            \sum_{j=1}^{n}r^{\,n-j}b_j.
        \]
        In particular, if $b_n=b$ and $r\neq 1$, then
        \[
            a_n
            =
            r^na_0+b\frac{r^n-1}{r-1}.
        \]
    \end{formula}

    \begin{formula}[Nonhomogeneous Linear Recurrences]
        For a linear recurrence with constant coefficients, the general solution
        is
        \[
            a_n=a_n^{(h)}+a_n^{(p)},
        \]
        where $a_n^{(h)}$ solves the associated homogeneous recurrence and
        $a_n^{(p)}$ is one particular solution. If the forcing term has the form
        \[
            P(n)\alpha^n,
        \]
        try a particular solution of the form
        \[
            n^sQ(n)\alpha^n,
        \]
        where $\deg Q=\deg P$ and $s$ is the multiplicity of $\alpha$ as a root
        of the characteristic polynomial.
    \end{formula}

    \begin{strategy}
        For recurrences, first identify whether the quickest language is
        characteristic roots, generating functions, a transfer matrix, or
        telescoping after a suitable normalization.
    \end{strategy}

    \begin{definition}[Limit of a Function; Epsilon--Delta Definition]
        Let $D\subseteq\R$, let $a$ be an accumulation point of $D$, and let
        $f\colon D\to\R$. The statement
        \[
            \lim\limits_{x\to a}f(x)=L
        \]
        means that $f(x)$ can be made arbitrarily close to $L$ by taking $x$
        sufficiently close to, but different from, $a$. Formally,
        \[
            \forall\varepsilon>0,
            \quad
            \exists\delta>0,
            \quad
            \forall x\in D,
            \qquad
            0<|x-a|<\delta
            \Rightarrow
            |f(x)-L|<\varepsilon.
        \]
    \end{definition}

    \begin{definition}[One-Sided Limits and Limits at Infinity]
        The right-hand limit $\lim\limits_{x\to a^+}f(x)=L$ is defined by
        \[
            \forall\varepsilon>0,
            \quad
            \exists\delta>0,
            \quad
            \forall x\in D,
            \qquad
            0<x-a<\delta
            \Rightarrow
            |f(x)-L|<\varepsilon.
        \]
        The left-hand definition is analogous, with $0<a-x<\delta$.
        Moreover,
        \[
            \lim\limits_{x\to\infty}f(x)=L
        \]
        means
        \[
            \forall\varepsilon>0,
            \quad
            \exists M>0,
            \quad
            \forall x\in D,
            \qquad
            x>M\Rightarrow |f(x)-L|<\varepsilon.
        \]
        Finally, $\lim\limits_{x\to a}f(x)=\infty$ means
        \[
            \forall M>0,
            \quad
            \exists\delta>0,
            \quad
            \forall x\in D,
            \qquad
            0<|x-a|<\delta\Rightarrow f(x)>M.
        \]
    \end{definition}

    \begin{theorem}[Squeeze Theorem; Sandwich Theorem]
        Suppose that $g(x)\leq f(x)\leq h(x)$ for all $x$ sufficiently close
        to $a$, except possibly at $a$, and that
        \[
            \lim\limits_{x\to a}g(x)
            =
            \lim\limits_{x\to a}h(x)
            =L.
        \]
        Then
        \[
            \lim\limits_{x\to a}f(x)=L.
        \]
    \end{theorem}

    \begin{definition}[Continuity at a Point]
        Let $f\colon D\to\R$ and $a\in D$. The function $f$ is continuous at $a$ if
        \[
            \lim\limits_{x\to a}f(x)=f(a).
        \]
        Equivalently,
        \[
            \forall\varepsilon>0,
            \quad
            \exists\delta>0,
            \quad
            \forall x\in D,
            \qquad
            |x-a|<\delta
            \Rightarrow
            |f(x)-f(a)|<\varepsilon.
        \]
        The function is continuous on $D$ if it is continuous at every point of
        $D$.
    \end{definition}

    \begin{definition}[Derivative]
        Let $f$ be defined near $a$. The derivative of $f$ at $a$ is
        \[
            f'(a)
            =
            \lim\limits_{h\to 0}\frac{f(a+h)-f(a)}{h},
        \]
        provided this limit exists as a finite real number. Thus $f'(a)=L$ if
        and only if
        \[
            \forall\varepsilon>0,
            \quad
            \exists\delta>0,
            \quad
            \forall h\in\R,
            \qquad
            0<|h|<\delta
            \Rightarrow
            \left|
                \frac{f(a+h)-f(a)}{h}-L
            \right|<\varepsilon.
        \]
    \end{definition}

    \begin{formula}[Standard Trigonometric and Exponential Limits]
        The following limits are fundamental:
        \[
            \lim\limits_{x\to 0}\frac{\sin x}{x}=1,
            \qquad
            \lim\limits_{x\to 0}\frac{1-\cos x}{x^2}=\frac12,
        \]
        \[
            \lim\limits_{x\to 0}\frac{e^x-1}{x}=1,
            \qquad
            \lim\limits_{x\to 0}\frac{\ln(1+x)}{x}=1,
        \]
        and
        \[
            \lim\limits_{x\to 0}(1+x)^{1/x}=e.
        \]
    \end{formula}

    \begin{strategy}
        These limits are often faster than L'H\^opital's Rule. The
        expressions $\sin x$, $1-\cos x$, $e^x-1$, $\ln(1+x)$, and
        $(1+x)^{1/x}$ are strong signals for these standard patterns.
    \end{strategy}

    \begin{formula}[Elementary Function Reference Table]
        The following domains, ranges, symmetries, and periods are standard.

        {\small
        \renewcommand{\arraystretch}{1.18}
        \begin{longtable}{@{}L{0.17\textwidth}L{0.24\textwidth}L{0.20\textwidth}L{0.13\textwidth}L{0.12\textwidth}@{}}
            \toprule
            Function & Domain & Range & Symmetry & Period \\
            \midrule
            \endfirsthead

            \toprule
            Function & Domain & Range & Symmetry & Period \\
            \midrule
            \endhead

            \bottomrule
            \endlastfoot

            $e^x$ & $\R$ & $(0,\infty)$ & neither & none \\
            $\ln x$ & $(0,\infty)$ & $\R$ & neither & none \\
            $\sin x$ & $\R$ & $[-1,1]$ & odd & $2\pi$ \\
            $\cos x$ & $\R$ & $[-1,1]$ & even & $2\pi$ \\
            $\tan x$ & $\R\setminus\{\frac{\pi}{2}+k\pi\mid k\in\Z\}$ & $\R$ & odd & $\pi$ \\
            $\sec x$ & $\R\setminus\{\frac{\pi}{2}+k\pi\mid k\in\Z\}$ & $(-\infty,-1]\cup[1,\infty)$ & even & $2\pi$ \\
            $\csc x$ & $\R\setminus\{k\pi\mid k\in\Z\}$ & $(-\infty,-1]\cup[1,\infty)$ & odd & $2\pi$ \\
            $\cot x$ & $\R\setminus\{k\pi\mid k\in\Z\}$ & $\R$ & odd & $\pi$ \\
            $\arcsin x$ & $[-1,1]$ & $[-\frac{\pi}{2},\frac{\pi}{2}]$ & odd & none \\
            $\arccos x$ & $[-1,1]$ & $[0,\pi]$ & neither & none \\
            $\arctan x$ & $\R$ & $(-\frac{\pi}{2},\frac{\pi}{2})$ & odd & none \\
            $\arcsec x$ & $(-\infty,-1]\cup[1,\infty)$ & $[0,\pi]\setminus\{\frac{\pi}{2}\}$ & neither & none \\
            $\arccsc x$ & $(-\infty,-1]\cup[1,\infty)$ & $[-\frac{\pi}{2},\frac{\pi}{2}]\setminus\{0\}$ & odd & none \\
            $\arccot x$ & $\R$ & $(0,\pi)$ & neither & none \\
            $\sinh x$ & $\R$ & $\R$ & odd & none \\
            $\cosh x$ & $\R$ & $[1,\infty)$ & even & none \\
            $\tanh x$ & $\R$ & $(-1,1)$ & odd & none \\
            $\sech x$ & $\R$ & $(0,1]$ & even & none \\
            $\csch x$ & $\R\setminus\{0\}$ & $\R\setminus\{0\}$ & odd & none \\
            $\coth x$ & $\R\setminus\{0\}$ & $(-\infty,-1)\cup(1,\infty)$ & odd & none \\
            $\arcsinh x$ & $\R$ & $\R$ & odd & none \\
            $\arccosh x$ & $[1,\infty)$ & $[0,\infty)$ & neither & none \\
            $\arctanh x$ & $(-1,1)$ & $\R$ & odd & none \\
        \end{longtable}
        }
    \end{formula}

    \begin{formula}[Exact Trigonometric Values Table]
        For the standard first-quadrant angles,

        {\small
        \renewcommand{\arraystretch}{1.18}
        \begin{longtable}{@{}L{0.16\textwidth}L{0.16\textwidth}L{0.16\textwidth}L{0.16\textwidth}L{0.16\textwidth}@{}}
            \toprule
            $x$ & $0$ & $\frac{\pi}{6}$ & $\frac{\pi}{4}$ & $\frac{\pi}{3}$ \\
            \midrule
            \endfirsthead

            \toprule
            $x$ & $0$ & $\frac{\pi}{6}$ & $\frac{\pi}{4}$ & $\frac{\pi}{3}$ \\
            \midrule
            \endhead

            \bottomrule
            \endlastfoot

            $\sin x$ & $0$ & $\frac12$ & $\frac{\sqrt2}{2}$ & $\frac{\sqrt3}{2}$ \\
            $\cos x$ & $1$ & $\frac{\sqrt3}{2}$ & $\frac{\sqrt2}{2}$ & $\frac12$ \\
            $\tan x$ & $0$ & $\frac{\sqrt3}{3}$ & $1$ & $\sqrt3$ \\
            $\sec x$ & $1$ & $\frac{2\sqrt3}{3}$ & $\sqrt2$ & $2$ \\
            $\csc x$ & undefined & $2$ & $\sqrt2$ & $\frac{2\sqrt3}{3}$ \\
            $\cot x$ & undefined & $\sqrt3$ & $1$ & $\frac{\sqrt3}{3}$ \\
        \end{longtable}
        }
        Moreover,
        \[
            \sin\frac{\pi}{2}=1,
            \qquad
            \cos\frac{\pi}{2}=0,
        \]
        and the remaining exact values follow from parity, periodicity, and
        reference-angle identities.
    \end{formula}

    \begin{formula}[Trigonometric Identity and Synthesis Table]
        The principal trigonometric identities are

        {\small
        \renewcommand{\arraystretch}{1.18}
        \begin{longtable}{@{}L{0.25\textwidth}L{0.67\textwidth}@{}}
            \toprule
            Type & Identities \\
            \midrule
            \endfirsthead

            \toprule
            Type & Identities \\
            \midrule
            \endhead

            \bottomrule
            \endlastfoot

            Reciprocal &
            $\displaystyle \sec x=\frac{1}{\cos x},\quad
            \csc x=\frac{1}{\sin x},\quad
            \cot x=\frac{1}{\tan x}$ \\
            Quotient &
            $\displaystyle \tan x=\frac{\sin x}{\cos x},\quad
            \cot x=\frac{\cos x}{\sin x}$ \\
            Pythagorean &
            $\sin^2x+\cos^2x=1$, $1+\tan^2x=\sec^2x$,
            $1+\cot^2x=\csc^2x$ \\
            Angle addition &
            $\sin(x\pm y)=\sin x\cos y\pm\cos x\sin y$ \\
            Angle addition &
            $\cos(x\pm y)=\cos x\cos y\mp\sin x\sin y$ \\
            Angle addition &
            $\displaystyle\tan(x\pm y)=
            \frac{\tan x\pm\tan y}{1\mp\tan x\tan y}$ \\
            Double angle &
            $\sin 2x=2\sin x\cos x$ \\
            Double angle &
            $\cos 2x=\cos^2x-\sin^2x=2\cos^2x-1=1-2\sin^2x$ \\
            Double angle &
            $\displaystyle\tan 2x=\frac{2\tan x}{1-\tan^2x}$ \\
            Half angle &
            $\displaystyle\sin^2\frac{x}{2}=\frac{1-\cos x}{2},\quad
            \cos^2\frac{x}{2}=\frac{1+\cos x}{2}$ \\
            Product-to-sum &
            $\displaystyle2\sin x\sin y=\cos(x-y)-\cos(x+y)$ \\
            Product-to-sum &
            $\displaystyle2\cos x\cos y=\cos(x-y)+\cos(x+y)$ \\
            Product-to-sum &
            $\displaystyle2\sin x\cos y=\sin(x+y)+\sin(x-y)$ \\
            Sum-to-product &
            $\displaystyle\sin x+\sin y=
            2\sin\frac{x+y}{2}\cos\frac{x-y}{2}$ \\
            Sum-to-product &
            $\displaystyle\cos x+\cos y=
            2\cos\frac{x+y}{2}\cos\frac{x-y}{2}$ \\
            Sum-to-product &
            $\displaystyle\sin x-\sin y=
            2\cos\frac{x+y}{2}\sin\frac{x-y}{2}$ \\
            Sum-to-product &
            $\displaystyle\cos x-\cos y=
            -2\sin\frac{x+y}{2}\sin\frac{x-y}{2}$ \\
            Hyperbolic &
            $\cosh^2x-\sinh^2x=1$, $1-\tanh^2x=\sech^2x$ \\
            Hyperbolic &
            $\displaystyle\tanh x=\frac{\sinh x}{\cosh x},\quad
            \coth x=\frac{\cosh x}{\sinh x}$ \\
            Inverse hyperbolic &
            $\displaystyle\arcsinh x=\ln\!\left(x+\sqrt{x^2+1}\right)$ \\
            Inverse hyperbolic &
            $\displaystyle\arccosh x=\ln\!\left(x+\sqrt{x^2-1}\right)\ (x\geq1)$ \\
            Inverse hyperbolic &
            $\displaystyle\arctanh x=\frac12\ln\!\left(\frac{1+x}{1-x}\right)\ (|x|<1)$ \\
            Synthesis &
            $\displaystyle a\cos x+b\sin x
            =R\cos(x-\alpha),\quad
            R=\sqrt{a^2+b^2},\quad
            \cos\alpha=\frac{a}{R},\quad
            \sin\alpha=\frac{b}{R}\quad(R>0)$ \\
        \end{longtable}
        }
    \end{formula}

    \begin{theorem}[L'H\^opital's Rule]
        Let $f$ and $g$ be differentiable on some punctured interval around $a$, and
        suppose that $g'(x)\neq 0$ there. Assume either
        \[
            f(x)\to 0,
            \qquad
            g(x)\to 0,
        \]
        or
        \[
            |f(x)|\to\infty,
            \qquad
            |g(x)|\to\infty,
        \]
        as $x\to a$. If
        \[
            \lim\limits_{x\to a}\frac{f'(x)}{g'(x)}=L
        \]
        exists in the extended real numbers, then
        \[
            \lim\limits_{x\to a}\frac{f(x)}{g(x)}=L.
        \]
    \end{theorem}

    \begin{strategy}
        Apply L'H\^opital's Rule only after identifying a $0/0$ or
        $\infty/\infty$ form. Rewrite $0\cdot\infty$ and $\infty-\infty$ as a
        quotient first. For $0^0$, $1^\infty$, or $\infty^0$, take logarithms
        before differentiating.
    \end{strategy}

    \begin{formula}[Derivative and Antiderivative Table]
        The following table records the basic one-variable forms. Affine changes
        of variable are handled by the chain rule and substitution. The constant
        of integration is omitted from the right column.

        {\small
        \renewcommand{\arraystretch}{1.18}
        \begin{longtable}{@{}L{0.26\textwidth}L{0.29\textwidth}L{0.35\textwidth}@{}}
            \toprule
            $f(x)$ & $f'(x)$ & $\displaystyle\int f(x)\dd x$ \\
            \midrule
            \endfirsthead

            \toprule
            $f(x)$ & $f'(x)$ & $\displaystyle\int f(x)\dd x$ \\
            \midrule
            \endhead

            \bottomrule
            \endlastfoot

            $c$ & $0$ & $cx$ \\
            $x$ & $1$ & $\displaystyle\frac{x^2}{2}$ \\
            $x^a$ & $ax^{a-1}$ & $\displaystyle\frac{x^{a+1}}{a+1}\quad(a\neq -1)$ \\
            $\displaystyle\frac{1}{x}$ & $\displaystyle -\frac{1}{x^2}$ & $\ln|x|$ \\
            $e^x$ & $e^x$ & $e^x$ \\
            $a^x$ & $a^x\ln a$ & $\displaystyle\frac{a^x}{\ln a}\quad(a>0,\ a\neq 1)$ \\
            $\ln x$ & $\displaystyle\frac{1}{x}$ & $x\ln x-x$ \\
            $\lg x$ & $\displaystyle\frac{1}{x\ln 10}$ & $\displaystyle\frac{x\ln x-x}{\ln 10}$ \\
            $\lb x$ & $\displaystyle\frac{1}{x\ln 2}$ & $\displaystyle\frac{x\ln x-x}{\ln 2}$ \\
            $\log_a x$ & $\displaystyle\frac{1}{x\ln a}$ & $\displaystyle\frac{x\ln x-x}{\ln a}\quad(a>0,\ a\neq 1)$ \\
            $x\ln x$ & $\ln x+1$ & $\displaystyle\frac{x^2}{2}\ln x-\frac{x^2}{4}$ \\
            $\displaystyle\frac{\ln x}{x}$ & $\displaystyle\frac{1-\ln x}{x^2}$ & $\displaystyle\frac{1}{2}(\ln x)^2$ \\
            $\displaystyle\frac{1}{x\ln x}$ & $\displaystyle-\frac{\ln x+1}{x^2(\ln x)^2}$ & $\ln|\ln x|$ \\
            $\sin x$ & $\cos x$ & $-\cos x$ \\
            $\cos x$ & $-\sin x$ & $\sin x$ \\
            $\tan x$ & $\sec^2x$ & $-\ln|\cos x|$ \\
            $\cot x$ & $-\csc^2x$ & $\ln|\sin x|$ \\
            $\sec x$ & $\sec x\tan x$ & $\ln|\sec x+\tan x|$ \\
            $\csc x$ & $-\csc x\cot x$ & $\ln|\csc x-\cot x|$ \\
            $\sec^2x$ & $2\sec^2x\tan x$ & $\tan x$ \\
            $\csc^2x$ & $-2\csc^2x\cot x$ & $-\cot x$ \\
            $\sec x\tan x$ & $\sec x\tan^2x+\sec^3x$ & $\sec x$ \\
            $\csc x\cot x$ & $-\csc x\cot^2x-\csc^3x$ & $-\csc x$ \\
            $\arcsin x$ & $\displaystyle\frac{1}{\sqrt{1-x^2}}$ & $x\arcsin x+\sqrt{1-x^2}$ \\
            $\arccos x$ & $\displaystyle-\frac{1}{\sqrt{1-x^2}}$ & $x\arccos x-\sqrt{1-x^2}$ \\
            $\arctan x$ & $\displaystyle\frac{1}{1+x^2}$ & $x\arctan x-\displaystyle\frac{1}{2}\ln(1+x^2)$ \\
            $\sinh x$ & $\cosh x$ & $\cosh x$ \\
            $\cosh x$ & $\sinh x$ & $\sinh x$ \\
            $\tanh x$ & $\sech^2x$ & $\ln(\cosh x)$ \\
            $\coth x$ & $-\csch^2x$ & $\ln|\sinh x|$ \\
            $\sech x$ & $-\sech x\tanh x$ & $\arctan(\sinh x)$ \\
            $\csch x$ & $-\csch x\coth x$ & $\ln\left|\tanh\frac{x}{2}\right|$ \\
            $\arcsinh x$ & $\displaystyle\frac{1}{\sqrt{x^2+1}}$ & $x\arcsinh x-\sqrt{x^2+1}$ \\
            $\arccosh x$ & $\displaystyle\frac{1}{\sqrt{x-1}\sqrt{x+1}}$ & $x\arccosh x-\sqrt{x^2-1}$ \\
            $\arctanh x$ & $\displaystyle\frac{1}{1-x^2}$ & $x\arctanh x+\displaystyle\frac{1}{2}\ln(1-x^2)$ \\
            $\displaystyle\frac{1}{x^2+a^2}$ & $\displaystyle-\frac{2x}{(x^2+a^2)^2}$ & $\displaystyle\frac{1}{a}\arctan\frac{x}{a}\quad(a>0)$ \\
            $\displaystyle\frac{1}{x^2-a^2}$ & $\displaystyle-\frac{2x}{(x^2-a^2)^2}$ & $\displaystyle\frac{1}{2a}\ln\left|\frac{x-a}{x+a}\right|\quad(a>0)$ \\
            $\displaystyle\frac{1}{a^2-x^2}$ & $\displaystyle\frac{2x}{(a^2-x^2)^2}$ & $\displaystyle\frac{1}{2a}\ln\left|\frac{a+x}{a-x}\right|\quad(a>0)$ \\
            $\displaystyle\frac{1}{\sqrt{a^2-x^2}}$ & $\displaystyle\frac{x}{(a^2-x^2)^{3/2}}$ & $\displaystyle\arcsin\frac{x}{a}\quad(a>0)$ \\
            $\displaystyle\frac{1}{\sqrt{x^2+a^2}}$ & $\displaystyle-\frac{x}{(x^2+a^2)^{3/2}}$ & $\ln\left|x+\sqrt{x^2+a^2}\right|$ \\
            $\displaystyle\frac{1}{\sqrt{x^2-a^2}}$ & $\displaystyle-\frac{x}{(x^2-a^2)^{3/2}}$ & $\ln\left|x+\sqrt{x^2-a^2}\right|\quad(a>0)$ \\
            $\displaystyle\frac{x}{\sqrt{x^2+a^2}}$ & $\displaystyle\frac{a^2}{(x^2+a^2)^{3/2}}$ & $\sqrt{x^2+a^2}$ \\
            $\displaystyle\frac{x}{\sqrt{a^2-x^2}}$ & $\displaystyle\frac{a^2}{(a^2-x^2)^{3/2}}$ & $-\sqrt{a^2-x^2}$ \\
            $\sqrt{a^2-x^2}$ & $\displaystyle-\frac{x}{\sqrt{a^2-x^2}}$ & $\displaystyle\frac{x}{2}\sqrt{a^2-x^2}+\frac{a^2}{2}\arcsin\frac{x}{a}$ \\
            $\sqrt{x^2+a^2}$ & $\displaystyle\frac{x}{\sqrt{x^2+a^2}}$ & $\displaystyle\frac{x}{2}\sqrt{x^2+a^2}+\frac{a^2}{2}\ln\left|x+\sqrt{x^2+a^2}\right|$ \\
            $\sqrt{x^2-a^2}$ & $\displaystyle\frac{x}{\sqrt{x^2-a^2}}$ & $\displaystyle\frac{x}{2}\sqrt{x^2-a^2}-\frac{a^2}{2}\ln\left|x+\sqrt{x^2-a^2}\right|$ \\
        \end{longtable}
        }
    \end{formula}

    \begin{tactic}
        An antiderivative is often exposed by a simple substitution, symmetry,
        completing the square, or integration by parts. Try those before expanding into a long calculation.
    \end{tactic}

    \begin{theorem}[Intermediate Value Theorem]
        Let $f$ be continuous on $[a,b]$. If $N$ is any number between
        $f(a)$ and $f(b)$, then there exists $c\in[a,b]$ s.t.
        \[
            f(c)=N.
        \]
    \end{theorem}

    \begin{strategy}
        The Intermediate Value Theorem is the theorem behind many existence
        arguments. To show that an equation has a solution, define a continuous
        function and show that its values cross the desired level.
    \end{strategy}

    \begin{theorem}[Extreme Value Theorem; Weierstrass Extreme Value Theorem]
        Let $f$ be continuous on $[a,b]$. Then $f$ attains both a maximum and
        a minimum on $[a,b]$.
    \end{theorem}

    \begin{remark}
        The closed interval condition is important. A continuous function on an
        open interval need not attain a maximum or a minimum.
    \end{remark}

    \begin{theorem}[Rolle's Theorem]
        Let $f$ be continuous on $[a,b]$ and differentiable on $(a,b)$. If
        $f(a)=f(b)$, then there exists $c\in(a,b)$ s.t.
        \[
            f'(c)=0.
        \]
    \end{theorem}

    \begin{theorem}[Mean Value Theorem; Lagrange Mean Value Theorem]
        Let $f$ be continuous on $[a,b]$ and differentiable on $(a,b)$. Then
        there exists $c\in(a,b)$ s.t.
        \[
            f'(c)
            =
            \frac{f(b)-f(a)}{b-a}.
        \]
    \end{theorem}

    \begin{idea}
        The Mean Value Theorem says that somewhere on the interval, the
        instantaneous rate of change equals the average rate of change.
    \end{idea}

    \begin{theorem}[Cauchy's Mean Value Theorem; Extended Mean Value Theorem]
        Let $f$ and $g$ be continuous on $[a,b]$ and differentiable on
        $(a,b)$. Then there exists $c\in(a,b)$ s.t.
        \[
            \bigl(f(b)-f(a)\bigr)g'(c)
            =
            \bigl(g(b)-g(a)\bigr)f'(c).
        \]
        If $g'(x)\neq 0$ for every $x\in(a,b)$, then $g(b)\neq g(a)$
        and the result may be written as
        \[
            \frac{f'(c)}{g'(c)}
            =
            \frac{f(b)-f(a)}{g(b)-g(a)}.
        \]
    \end{theorem}

    \begin{theorem}[Derivative of an Inverse Function]
        Let $I\subseteq\R$ be an interval, let $f\colon I\to\R$ be one-to-one, and
        suppose that $f$ is differentiable at an interior point $x_0\in I$ with
        $f'(x_0)\neq0$. If the inverse $g=f^{-1}$ is continuous at
        $y_0=f(x_0)$, then $g$ is differentiable at $y_0$ and
        \[
            g'(y_0)=\frac{1}{f'(x_0)}.
        \]
        In particular, the continuity hypothesis is automatic when $f$ is
        continuous and strictly monotone on an interval.
    \end{theorem}

    \begin{definition}[Riemann Integral]
        Let $f\colon [a,b]\to\R$ be bounded. For a partition
        \[
            P:\quad a=x_0<x_1<\cdots<x_n=b,
        \]
        choose $\xi_i\in[x_{i-1},x_i]$ and form the Riemann sum
        \[
            \sum_{i=1}^{n}f(\xi_i)(x_i-x_{i-1}).
        \]
        If these sums approach the same limit $I$ as the mesh
        $\max_i(x_i-x_{i-1})\to0$, independently of the sample points, then
        $f$ is Riemann integrable and
        \[
            I=\int_a^b f(x)\dd x.
        \]
        Every continuous function on a closed bounded interval is Riemann
        integrable.
    \end{definition}

    \begin{theorem}[Fundamental Theorem of Calculus]
        If $f$ is continuous on $[a,b]$ and
        \[
            F(x)=\int_a^x f(t)\dd t,
        \]
        then $F'(x)=f(x)$. Conversely, if $F'=f$ on $[a,b]$, then
        \[
            \int_a^b f(x)\dd x=F(b)-F(a).
        \]
    \end{theorem}

    \begin{theorem}[Taylor's Theorem]
        Let $f$ be $(n+1)$ times differentiable on an open interval containing
        the closed interval with endpoints $a$ and $x$. Then there exists a
        point $\xi$ strictly between $a$ and $x$ s.t.
        \[
            f(x)
            =
            \sum_{k=0}^{n}\frac{f^{(k)}(a)}{k!}(x-a)^k
            +
            \frac{f^{(n+1)}(\xi)}{(n+1)!}(x-a)^{n+1}.
        \]
        The final term is the Lagrange remainder.
    \end{theorem}

    \begin{corollary}[Taylor Error Bound]
        If
        \[
            |f^{(n+1)}(t)|\leq M
        \]
        for every $t$ between $a$ and $x$, then the Taylor remainder satisfies
        \[
            |R_n(x)|
            \leq
            \frac{M}{(n+1)!}|x-a|^{n+1}.
        \]
    \end{corollary}

    \begin{strategy}
        Taylor expansion is one of the fastest tools for limits, approximations,
        and hidden series evaluations. Expressions involving $e^x$, $\sin x$,
        $\cos x$, $\ln(1+x)$, or $\arctan x$ often invite Taylor expansion.
    \end{strategy}

    \begin{definition}[Power Series and Radius of Convergence]
        A power series centered at $a$ is a series of the form
        \[
            \sum_{n=0}^{\infty}c_n(x-a)^n.
        \]
        There is a radius of convergence $R\in[0,\infty]$ s.t. the series
        converges absolutely for $|x-a|<R$ and diverges for $|x-a|>R$; behavior
        at the boundary $|x-a|=R$ must be checked separately. A Taylor series is
        the power series with coefficients
        \[
            c_n=\frac{f^{(n)}(a)}{n!}.
        \]
        When $a=0$, it is called a Maclaurin series.
    \end{definition}

    \begin{formula}[Maclaurin Series]
        The most important Maclaurin series are
        \[
            e^x
            =
            \sum_{n=0}^{\infty}\frac{x^n}{n!},
        \]
        \[
            \sin x
            =
            \sum_{n=0}^{\infty}(-1)^n\frac{x^{2n+1}}{(2n+1)!},
            \qquad
            \cos x
            =
            \sum_{n=0}^{\infty}(-1)^n\frac{x^{2n}}{(2n)!},
        \]
        \[
            \frac{1}{1-x}
            =
            \sum_{n=0}^{\infty}x^n
            \qquad (|x|<1),
        \]
        \[
            \ln(1+x)
            =
            \sum_{n=1}^{\infty}(-1)^{n+1}\frac{x^n}{n}
            \qquad (-1<x\leq 1),
        \]
        and
        \[
            \arctan x
            =
            \sum_{n=0}^{\infty}(-1)^n\frac{x^{2n+1}}{2n+1}
            \qquad (|x|\leq 1).
        \]
        Also,
        \[
            -\ln(1-x)
            =
            \sum_{n=1}^{\infty}\frac{x^n}{n}
            \qquad (-1\leq x<1),
        \]
        \[
            \sinh x
            =
            \sum_{n=0}^{\infty}\frac{x^{2n+1}}{(2n+1)!},
            \qquad
            \cosh x
            =
            \sum_{n=0}^{\infty}\frac{x^{2n}}{(2n)!},
        \]
        and
        \[
            \arctanh x
            =
            \sum_{n=0}^{\infty}\frac{x^{2n+1}}{2n+1}
            \qquad (|x|<1).
        \]
    \end{formula}

    \begin{formula}[Binomial Series]
        For any real or complex number $\beta$,
        \[
            (1+x)^\beta
            =
            \sum_{n=0}^{\infty}\binom{\beta}{n}x^n
            \qquad (|x|<1),
        \]
        where
        \[
            \binom{\beta}{n}
            =
            \frac{\beta(\beta-1)\cdots(\beta-n+1)}{n!}.
        \]
    \end{formula}

    \begin{formula}[Chebyshev Polynomials]
        The Chebyshev polynomials of the first kind are characterized by
        \[
            T_n(\cos\theta)=\cos(n\theta).
        \]
        They satisfy
        \[
            T_0(x)=1,
            \qquad
            T_1(x)=x,
            \qquad
            T_{n+1}(x)=2xT_n(x)-T_{n-1}(x).
        \]
        Hence multiple-angle trigonometric expressions can be converted into
        polynomial calculations almost immediately; for example,
        \[
            \cos(3\theta)=4\cos^3\theta-3\cos\theta.
        \]
    \end{formula}

    \begin{formula}[Classical Series Values]
        The following standard series values are useful in computation:
        \[
            \sum_{n=1}^{\infty}\frac{1}{n^2}
            =
            \frac{\pi^2}{6},
            \qquad
            \sum_{n=1}^{\infty}\frac{1}{n^4}
            =
            \frac{\pi^4}{90},
        \]
        \[
            \sum_{n=1}^{\infty}\frac{(-1)^{n+1}}{n}
            =
            \ln 2,
            \qquad
            \sum_{n=0}^{\infty}\frac{(-1)^n}{2n+1}
            =
            \frac{\pi}{4}.
        \]
    \end{formula}

    \begin{formula}[Geometric-Series Differentiation]
        Starting from
        \[
            \frac{1}{1-x}
            =
            \sum_{n=0}^{\infty}x^n
            \qquad (|x|<1),
        \]
        term-by-term differentiation gives
        \[
            \frac{1}{(1-x)^2}
            =
            \sum_{n=1}^{\infty}nx^{n-1}
            \qquad (|x|<1),
        \]
        and hence
        \[
            \sum_{n=1}^{\infty}nx^n
            =
            \frac{x}{(1-x)^2}
            \qquad (|x|<1).
        \]
    \end{formula}

    \begin{formula}[Telescoping Series]
        If
        \[
            a_n=b_n-b_{n+1},
        \]
        then the partial sums satisfy
        \[
            \sum_{n=1}^{N}a_n
            =
            b_1-b_{N+1}.
        \]
        Hence, if $b_n\to L$, then
        \[
            \sum_{n=1}^{\infty}a_n=b_1-L.
        \]
    \end{formula}

    \begin{formula}[Abel's Summation Formula; Summation by Parts]
        Let
        \[
            A_k=\sum_{j=m}^{k}a_j.
        \]
        Then
        \[
            \sum_{k=m}^{n}a_kb_k
            =
            A_nb_n
            +
            \sum_{k=m}^{n-1}A_k(b_k-b_{k+1}).
        \]
        This is the discrete analogue of integration by parts and is useful when
        one factor has simple partial sums while the other changes slowly.
    \end{formula}

    \begin{formula}[Euler--Maclaurin Summation Formula]
        Let $p\geq1$ and let $f\in C^{2p}([m,n])$, where $m<n$ are integers.
        Then
        \[
            \sum_{k=m}^{n}f(k)
            =
            \int_m^n f(x)\dd x
            +
            \frac{f(m)+f(n)}{2}
            +
            \sum_{r=1}^{p-1}
            \frac{B_{2r}}{(2r)!}
            \left(f^{(2r-1)}(n)-f^{(2r-1)}(m)\right)
            +R_p,
        \]
        where $B_{2r}$ are the Bernoulli numbers and
        \[
            |R_p|
            \leq
            \frac{2\zeta(2p)}{(2\pi)^{2p}}
            \int_m^n|f^{(2p)}(x)|\dd x.
        \]
        The case $p=1$ gives the trapezoidal approximation together with a
        rigorous error bound.
    \end{formula}

    \begin{formula}[Asymptotic Growth Hierarchy]
        For fixed $p,q>0$ and $c>1$, as $n\to\infty$,
        \[
            (\ln n)^p=o(n^q),
            \qquad
            n^q=o(c^n),
            \qquad
            c^n=o(n!),
            \qquad
            n!=o(n^n).
        \]
    \end{formula}

    \begin{strategy}
        For fast series calculations, first try to reduce the expression to a
        geometric series, a differentiated geometric series, a Maclaurin series,
        or a known special value such as $\pi^2/6$, $\pi^4/90$, $\ln 2$, or
        $\pi/4$.
    \end{strategy}

    \begin{formula}[Integration by Parts]
        If $u$ and $v$ are differentiable functions, then
        \[
            \int u\dd v
            =
            uv-\int v\dd u.
        \]
        Equivalently,
        \[
            \int_a^b u(x)v'(x)\dd x
            =
            \left[u(x)v(x)\right]_a^b
            -
            \int_a^b u'(x)v(x)\dd x.
        \]
    \end{formula}

    \begin{theorem}[Leibniz Integral Rule; Differentiation under the Integral Sign]
        Let $I$ be an open interval, let $a,b\colon I\to\R$ be continuously
        differentiable, and suppose that $f$ and $\partial f/\partial t$ are
        continuous on
        \[
            \{(x,t)\mid t\in I,\ a(t)\leq x\leq b(t)\}.
        \]
        Then
        \[
            \frac{\dd}{\dd t}
            \int_{a(t)}^{b(t)}f(x,t)\dd x
            =
            f(b(t),t)b'(t)-f(a(t),t)a'(t)
            +
            \int_{a(t)}^{b(t)}
            \frac{\partial f}{\partial t}(x,t)\dd x.
        \]
        In particular, if the endpoints are constant, only the integral of the
        partial derivative remains.
    \end{theorem}

    \begin{tactic}
        Introducing a parameter and differentiating under the integral sign is
        often called Feynman's trick. It is especially effective when direct
        integration is difficult but the parameter derivative is elementary.
    \end{tactic}

    \begin{formula}[Frullani's Integral]
        Let $f\colon [0,\infty)\to\R$ be continuous, suppose that the finite limits
        $f(0)$ and $f(\infty)\coloneqq\lim\limits_{x\to\infty}f(x)$ exist, and assume
        that the improper integral below converges. For $a,b>0$,
        \[
            \int_0^\infty\frac{f(ax)-f(bx)}{x}\dd x
            =
            \bigl(f(0)-f(\infty)\bigr)\ln\frac{b}{a}.
        \]
        A standard sufficient hypothesis is that $f$ is continuously
        differentiable and $f'\in L^1(0,\infty)$.
    \end{formula}

    \begin{formula}[Arc Length Formula]
        The length of the graph $y=f(x)$ on $[a,b]$ is
        \[
            L
            =
            \int_a^b \sqrt{1+(f'(x))^2}\dd x.
        \]
        More generally, if a curve is parametrized by
        $\mathbf{r}(t)=(x(t),y(t),z(t))$, then
        \[
            L
            =
            \int_a^b \|\mathbf{r}'(t)\|\dd t.
        \]
    \end{formula}

    \begin{formula}[Vi\`ete's Infinite Product]
        For every real $x$,
        \[
            \frac{\sin x}{x}
            =
            \prod_{k=1}^{\infty}
            \cos\left(\frac{x}{2^k}\right),
        \]
        where the value at $x=0$ is interpreted by continuity. Setting
        $x=\pi/2$ gives Vi\`ete's classical product for $2/\pi$.
    \end{formula}

    \begin{idea}
        Vi\`ete's formula is a product formula for $\pi$. The key identity is
        the double-angle formula
        \[
            \sin(2x)=2\sin x\cos x.
        \]
    \end{idea}

    \begin{formula}[Euler's Formula]
        For every real number $x$,
        \[
            e^{ix}
            =
            \cos x+i\sin x.
        \]
        In particular,
        \[
            e^{i\pi}+1=0.
        \]
    \end{formula}

    \begin{remark}
        The identity $e^{i\pi}+1=0$ is called Euler's identity. It connects
        the constants $e$, $i$, $\pi$, $1$, and $0$ in a single formula.
    \end{remark}

    \begin{formula}[Gaussian Integral]
        The Gaussian integral is
        \[
            \int_{-\infty}^{\infty}e^{-x^2}\dd x
            =
            \sqrt{\pi}.
        \]
        Equivalently,
        \[
            \int_0^\infty e^{-x^2}\dd x
            =
            \frac{\sqrt{\pi}}{2}.
        \]
    \end{formula}

    \begin{remark}
        The Gaussian integral is one of the most famous integrals in calculus.
        It is closely connected to the normal distribution in probability.
    \end{remark}

    \begin{formula}[Fresnel Integrals]
        The Fresnel integrals satisfy
        \[
            \int_0^{\infty}\cos(x^2)\dd x
            =
            \int_0^{\infty}\sin(x^2)\dd x
            =
            \sqrt{\frac{\pi}{8}}.
        \]
    \end{formula}

    \begin{remark}
        Fresnel integrals appear in wave optics, diffraction, and the Cornu
        spiral. They are memorable examples of non-elementary definite
        integrals with clean values.
    \end{remark}

    \begin{formula}[Sine Integral; Dirichlet Integral]
        The sine integral is
        \[
            \operatorname{Si}(x)
            =
            \int_0^x \frac{\sin t}{t}\dd t.
        \]
        The classical Dirichlet integral is
        \[
            \int_0^{\infty}\frac{\sin x}{x}\dd x
            =
            \frac{\pi}{2}.
        \]
    \end{formula}

    \begin{formula}[Arctangent Integral]
        For $a>0$,
        \[
            \int_{-\infty}^{\infty}\frac{1}{x^2+a^2}\dd x
            =
            \frac{\pi}{a},
            \qquad
            \int_0^{\infty}\frac{1}{x^2+a^2}\dd x
            =
            \frac{\pi}{2a}.
        \]
        More generally,
        \[
            \int_{-\infty}^{\infty}\frac{c}{(x-b)^2+a^2}\dd x
            =
            \frac{c\pi}{a}.
        \]
    \end{formula}

    \begin{definition}[Gamma Function]
        For $\Real(z)>0$, the Gamma function is defined by
        \[
            \Gamma(z)
            =
            \int_0^\infty t^{z-1}e^{-t}\dd t.
        \]
        It extends meromorphically to $\C$, with simple poles at the
        nonpositive integers.
    \end{definition}

    \begin{proposition}[Gamma Function and Factorials]
        The Gamma function satisfies
        \[
            \Gamma(z+1)=z\Gamma(z)
        \]
        wherever both sides are defined.
        In particular, for every positive integer $n$,
        \[
            \Gamma(n)=(n-1)!.
        \]
        Also,
        \[
            \Gamma\left(\frac12\right)=\sqrt{\pi}.
        \]
    \end{proposition}

    \begin{definition}[Beta Function; Euler Beta Integral]
        For $x,y>0$, the Beta function is defined by
        \[
            B(x,y)
            =
            \int_0^1 t^{x-1}(1-t)^{y-1}\dd t.
        \]
    \end{definition}

    \begin{formula}[Beta--Gamma Identity]
        The Beta and Gamma functions are related by
        \[
            B(x,y)
            =
            \frac{\Gamma(x)\Gamma(y)}{\Gamma(x+y)}.
        \]
        In particular, for nonnegative integers $m,n$,
        \[
            \int_0^1 x^m(1-x)^n\dd x
            =
            B(m+1,n+1)
            =
            \frac{m!n!}{(m+n+1)!}
            =
            \frac{1}{(m+n+1)\binom{m+n}{m}}.
        \]
    \end{formula}

    \begin{formula}[Gamma-Type Exponential Integrals]
        For $a>0$ and every nonnegative integer $n$,
        \[
            \int_0^{\infty}x^n e^{-ax}\dd x
            =
            \frac{n!}{a^{n+1}}.
        \]
        More generally, for $\Real(s)>0$,
        \[
            \int_0^{\infty}x^{s-1}e^{-ax}\dd x
            =
            \frac{\Gamma(s)}{a^s}.
        \]
    \end{formula}

    \begin{formula}[Damped Sine and Cosine Integrals]
        For $a>0$,
        \[
            \int_0^{\infty}e^{-ax}\cos bx\dd x
            =
            \frac{a}{a^2+b^2},
            \qquad
            \int_0^{\infty}e^{-ax}\sin bx\dd x
            =
            \frac{b}{a^2+b^2}.
        \]
    \end{formula}

    \begin{strategy}
        Many fast integral problems reduce to this formula after completing a
        square or translating the variable. A shift such as $u=x-b$ does not
        change the infinite limits.
    \end{strategy}

    \begin{formula}[Symmetry Trick for Definite Integrals]
        If $f$ is integrable on $[a,b]$, then
        \[
            \int_a^b f(x)\dd x
            =
            \int_a^b f(a+b-x)\dd x.
        \]
        Hence, if
        \[
            f(x)+f(a+b-x)=C
        \]
        is constant, then
        \[
            \int_a^b f(x)\dd x
            =
            \frac{C(b-a)}{2}.
        \]
    \end{formula}

    \begin{strategy}
        This is one of the most useful Integration-Bee-style tricks. Typical
        signals are intervals such as $[0,\pi/2]$, expressions involving
        $\tan x$ and $\cot x$, or a denominator that becomes complementary after
        the substitution $x\mapsto a+b-x$.
    \end{strategy}

    \begin{formula}[Parity and Periodicity Shortcuts for Integrals]
        If $f$ is integrable on $[-a,a]$, then
        \[
            f(-x)=-f(x)
            \quad\Longrightarrow\quad
            \int_{-a}^{a}f(x)\dd x=0,
        \]
        and
        \[
            f(-x)=f(x)
            \quad\Longrightarrow\quad
            \int_{-a}^{a}f(x)\dd x
            =
            2\int_{0}^{a}f(x)\dd x.
        \]
        If $f$ has period $T$, then for every $c$,
        \[
            \int_c^{c+T}f(x)\dd x
            =
            \int_0^T f(x)\dd x.
        \]
    \end{formula}

    \begin{formula}[Reciprocal-Substitution Identity]
        If the improper integrals converge, the substitution $x=1/t$ gives
        \[
            \int_0^\infty f(x)\dd x
            =
            \int_0^\infty \frac{f(1/x)}{x^2}\dd x.
        \]
        Therefore,
        \[
            \int_0^\infty f(x)\dd x
            =
            \frac12
            \int_0^\infty
            \left(f(x)+\frac{f(1/x)}{x^2}\right)\dd x.
        \]
        This often reveals a cancellation or symmetry hidden in an integral on
        $(0,\infty)$.
    \end{formula}

    \begin{formula}[Weierstrass Substitution; Tangent Half-Angle Substitution]
        The substitution
        \[
            t=\tan\frac{x}{2}
        \]
        converts rational expressions in $\sin x$ and $\cos x$ into rational
        functions of $t$:
        \[
            \sin x=\frac{2t}{1+t^2},
            \qquad
            \cos x=\frac{1-t^2}{1+t^2},
            \qquad
            \dd x=\frac{2\dd t}{1+t^2}.
        \]
    \end{formula}

    \begin{algorithm}[Fast Integral Recognition Checklist]
        Before expanding or integrating directly, check in this order:
        \begin{bookitems}
            \item parity, periodicity, and interval symmetry;
            \item the substitutions $x\mapsto a-x$, $x\mapsto 1/x$, or
                  $x\mapsto a/x$;
            \item a hidden derivative for substitution or integration by parts;
            \item a parameter that permits differentiation under the integral sign;
            \item Beta--Gamma, Gaussian, Wallis, or a standard trigonometric
                  integral;
            \item the tangent half-angle substitution for rational trigonometric
                  expressions.
        \end{bookitems}
    \end{algorithm}

    \begin{formula}[Wallis Product]
        The Wallis product is
        \[
            \frac{\pi}{2}
            =
            \prod_{n=1}^{\infty}
            \frac{(2n)(2n)}{(2n-1)(2n+1)}
            =
            \frac{2}{1}\cdot\frac{2}{3}\cdot
            \frac{4}{3}\cdot\frac{4}{5}\cdot
            \frac{6}{5}\cdot\frac{6}{7}\cdots.
        \]
    \end{formula}

    \begin{formula}[Euler's Reflection Formula]
        For $z\in\C\setminus\Z$,
        \[
            \Gamma(z)\Gamma(1-z)
            =
            \frac{\pi}{\sin(\pi z)}.
        \]
    \end{formula}

    \begin{formula}[Euler Beta Integral on the Positive Half-Line]
        If $p>1$, then
        \[
            \int_0^\infty\frac{\dd x}{1+x^p}
            =
            \frac{1}{p}
            B\left(\frac{1}{p},1-\frac{1}{p}\right)
            =
            \frac{\pi}{p}\csc\left(\frac{\pi}{p}\right).
        \]
        The substitution
        \[
            t=\frac{x^p}{1+x^p}
        \]
        converts the improper integral into an Euler beta integral.
    \end{formula}

    \begin{formula}[Stirling's Formula]
        As $n\to\infty$,
        \[
            n!
            \sim
            \sqrt{2\pi n}\left(\frac{n}{e}\right)^n.
        \]
    \end{formula}

    \begin{remark}
        Stirling's formula is the standard approximation for large factorials.
        It often appears in counting, probability, and asymptotic estimation.
    \end{remark}

    \begin{formula}[Trapezoidal Rule]
        The trapezoidal approximation is
        \[
            \int_a^b f(x)\dd x
            \approx
            \frac{b-a}{2}\bigl(f(a)+f(b)\bigr).
        \]
        It has degree of precision $1$.
    \end{formula}

    \begin{formula}[Simpson's Rule]
        Simpson's rule is
        \[
            \int_a^b f(x)\dd x
            \approx
            \frac{b-a}{6}
            \left(
                f(a)+4f\left(\frac{a+b}{2}\right)+f(b)
            \right).
        \]
        It has degree of precision $3$.
    \end{formula}

    \begin{remark}
        Simpson's rule uses quadratic interpolation, but the resulting formula
        is exact for every polynomial of degree at most $3$.
    \end{remark}

    \begin{theorem}[$p$-Series Test]
        The series
        \[
            \sum_{n=1}^{\infty}\frac{1}{n^p}
        \]
        converges if and only if $p>1$.
    \end{theorem}

    \begin{theorem}[Direct Comparison Test]
        Let $0\leq a_n\leq b_n$ for all sufficiently large $n$. If
        $\sum b_n$ converges, then $\sum a_n$ converges. If $\sum a_n$
        diverges, then $\sum b_n$ diverges.
    \end{theorem}

    \begin{proposition}[Absolute Convergence Implies Convergence]
        If
        \[
            \sum_{n=1}^{\infty}|a_n|
        \]
        converges, then $\sum_{n=1}^{\infty}a_n$ converges.
    \end{proposition}

    \begin{theorem}[Integral Test]
        Suppose $f$ is positive, continuous, and decreasing on $[1,\infty)$,
        and let $a_n=f(n)$. Then
        \[
            \sum_{n=1}^{\infty}a_n
        \]
        and
        \[
            \int_1^{\infty}f(x)\dd x
        \]
        either both converge or both diverge.
    \end{theorem}

    \begin{theorem}[Limit Comparison Test]
        Let $a_n,b_n>0$. If
        \[
            \lim\limits_{n\to\infty}\frac{a_n}{b_n}=c
        \]
        for some finite number $c>0$, then
        \[
            \sum_{n=1}^{\infty}a_n
            \qquad\text{and}\qquad
            \sum_{n=1}^{\infty}b_n
        \]
        converge or diverge together.
    \end{theorem}

    \begin{theorem}[Ratio Test; d'Alembert's Ratio Test]
        For a series $\sum a_n$ with $a_n\neq0$ eventually, set
        \[
            L=\limsup_{n\to\infty}
            \left|\frac{a_{n+1}}{a_n}\right|.
        \]
        If $L<1$, then $\sum a_n$ converges absolutely. If
        \[
            \liminf_{n\to\infty}
            \left|\frac{a_{n+1}}{a_n}\right|>1,
        \]
        then $a_n\not\to0$ and the series diverges. In the common case in which
        the ordinary limit
        \[
            \lim\limits_{n\to\infty}
            \left|\frac{a_{n+1}}{a_n}\right|=L
        \]
        exists, this reduces to the familiar rule: $L<1$ gives absolute
        convergence, $L>1$ gives divergence, and $L=1$ is inconclusive.
    \end{theorem}

    \begin{theorem}[Root Test; Cauchy's Root Test]
        For a series $\sum a_n$, set
        \[
            L=\limsup_{n\to\infty}\sqrt[n]{|a_n|}.
        \]
        If $L<1$, then $\sum a_n$ converges absolutely. If $L>1$, then the
        series diverges. If $L=1$, the test is inconclusive. When the ordinary
        limit $\lim\limits_{n\to\infty}\sqrt[n]{|a_n|}$ exists, it equals the
        displayed $\limsup$, giving the usual limit form of the test.
    \end{theorem}

    \begin{theorem}[Cauchy--Hadamard Theorem]
        For a power series
        \[
            \sum_{n=0}^{\infty}a_n(x-a)^n,
        \]
        the radius of convergence $R$ is given by
        \[
            \frac{1}{R}
            =
            \limsup_{n\to\infty}\sqrt[n]{|a_n|}.
        \]
        Equivalently, the series converges absolutely for $|x-a|<R$ and
        diverges for $|x-a|>R$.
    \end{theorem}

    \begin{remark}
        This is the standard theorem behind the phrase ``radius of convergence''
        for a power series. It is also the source of the name
        Cauchy--Hadamard.
    \end{remark}

    \begin{theorem}[Termwise Differentiation and Integration of Power Series]
        Suppose
        \[
            f(x)=\sum_{n=0}^{\infty}a_n(x-a)^n
        \]
        has radius of convergence $R>0$. For $|x-a|<R$,
        \[
            f'(x)
            =
            \sum_{n=1}^{\infty}n a_n(x-a)^{n-1},
        \]
        and
        \[
            \int_a^x f(t)\dd t
            =
            \sum_{n=0}^{\infty}\frac{a_n}{n+1}(x-a)^{n+1}.
        \]
        The differentiated and integrated series have the same radius of
        convergence $R$. Convergence at the endpoints must be checked
        separately.
    \end{theorem}

    \begin{theorem}[Alternating Series Test; Leibniz Test]
        If
        \[
            a_1\geq a_2\geq a_3\geq\cdots\geq 0
        \]
        and $a_n\to 0$, then the alternating series
        \[
            \sum_{n=1}^{\infty}(-1)^{n+1}a_n
        \]
        converges.
    \end{theorem}

    \begin{formula}[Alternating Series Error Estimate]
        Under the hypotheses of the Alternating Series Test, if
        \[
            S=\sum_{n=1}^{\infty}(-1)^{n+1}a_n,
            \qquad
            S_N=\sum_{n=1}^{N}(-1)^{n+1}a_n,
        \]
        then
        \[
            |S-S_N|\leq a_{N+1}.
        \]
    \end{formula}

    \begin{strategy}
        Factorials and powers often suggest the Ratio Test, $n$th powers suggest the Root Test, and
        alternating signs with decreasing terms suggest the Alternating Series
        Test.
    \end{strategy}

    \begin{theorem}[Divergence Test; Term Test]
        If the series
        \[
            \sum_{n=1}^{\infty}a_n
        \]
        converges, then
        \[
            \lim\limits_{n\to\infty}a_n=0.
        \]
        Equivalently, if $\lim\limits_{n\to\infty}a_n$ does not exist or is not $0$,
        then $\sum_{n=1}^{\infty}a_n$ diverges.
    \end{theorem}

    \begin{strategy}
        Before trying a sophisticated convergence test, first check whether the
        terms go to $0$. If the terms do not tend to $0$, the series diverges
        immediately.
    \end{strategy}

    \begin{algorithm}[Fast Series Checklist]
        For a convergence or summation problem:
        \begin{bookitems}
            \item Check the term limit first.
            \item Look for telescoping, a geometric series, or a known power
                  series.
            \item For positive terms, compare only the dominant asymptotic
                  factors.
            \item Use the Ratio Test for factorials and exponentials, and the
                  Root Test for expressions raised to the $n$th power.
            \item For alternating terms, check monotonicity and use the first
                  omitted term to control the error.
        \end{bookitems}
    \end{algorithm}

    \section*{Multivariable Calculus}

    \begin{definition}[Open Sets in Euclidean Space]
        For $\mathbf{a}\in\R^n$ and $r>0$, the open ball of radius $r$ centered at
        $\mathbf{a}$ is
        \[
            B_r(\mathbf{a})
            =
            \{\mathbf{x}\in\R^n\mid \|\mathbf{x}-\mathbf{a}\|<r\},
        \]
        where the Euclidean norm is
        \[
            \|\mathbf{x}\|
            =
            \sqrt{x_1^2+\cdots+x_n^2}.
        \]
        A set $U\subseteq\R^n$ is open if every $\mathbf{a}\in U$ has some
        $r>0$ s.t.
        \[
            B_r(\mathbf{a})\subseteq U.
        \]
        The Real Analysis section later gives the corresponding definition in an
        arbitrary metric space.
    \end{definition}

    \begin{definition}[Scalar Field, Vector Field, and Partial Derivative]
        A scalar field on $U\subseteq\R^n$ is a function $f\colon U\to\R$, while a
        vector field is a function $\mathbf{F}\colon U\to\R^m$. For a scalar-valued
        function $f(x_1,\dots,x_n)$, the partial derivative with respect to
        $x_i$ is the one-variable derivative obtained by varying $x_i$ while
        holding the other variables fixed:
        \[
            \frac{\partial f}{\partial x_i}(\mathbf{a})
            =
            \lim\limits_{h\to0}
            \frac{f(\mathbf{a}+h\mathbf{e}_i)-f(\mathbf{a})}{h},
        \]
        whenever the limit exists.
    \end{definition}

    \begin{definition}[Gradient, Divergence, Curl, and Laplacian]
        For a scalar-valued function $f(x,y,z)$ and a vector field
        $\mathbf{F}=(P,Q,R)^{\mathsf T}$, we write
        \[
            \nabla f
            =
            \begin{pmatrix}
                \partial f/\partial x\\
                \partial f/\partial y\\
                \partial f/\partial z
            \end{pmatrix},
        \]
        \[
            \nabla\cdot\mathbf{F}
            =
            \frac{\partial P}{\partial x}
            +
            \frac{\partial Q}{\partial y}
            +
            \frac{\partial R}{\partial z},
        \]
        \[
            \nabla\times\mathbf{F}
            =
            \begin{vmatrix}
                \mathbf{i} & \mathbf{j} & \mathbf{k} \\
                \partial/\partial x & \partial/\partial y & \partial/\partial z \\
                P & Q & R
            \end{vmatrix}
            =
            \begin{pmatrix}
                \partial R/\partial y-\partial Q/\partial z\\
                \partial P/\partial z-\partial R/\partial x\\
                \partial Q/\partial x-\partial P/\partial y
            \end{pmatrix},
        \]
        and
        \[
            \Delta f
            =
            \nabla^2f
            =
            \nabla\cdot(\nabla f)
            =
            f_{xx}+f_{yy}+f_{zz}.
        \]
    \end{definition}

    \begin{idea}
        The gradient points in the direction of steepest increase. Divergence
        measures source or sink behavior. Curl measures local rotation. The
        Laplacian measures how a value differs from its surrounding average.
    \end{idea}

    \begin{theorem}[Clairaut's Theorem; Mixed Partial Derivative Theorem]
        If $f$ has continuous second partial derivatives near a point, then the
        mixed partial derivatives agree at that point:
        \[
            \frac{\partial^2 f}{\partial x\partial y}
            =
            \frac{\partial^2 f}{\partial y\partial x}.
        \]
    \end{theorem}

    \begin{definition}[Jacobian Matrix and Jacobian Determinant]
        Let $\mathbf{F}\colon U\subseteq\R^n\to\R^m$ have component functions
        $F_1,\dots,F_m$. Its Jacobian matrix at $\mathbf{a}$ is
        \[
            \mathbf{J}_{\mathbf{F}}(\mathbf{a})
            =
            \left(\frac{\partial F_i}{\partial x_j}(\mathbf{a})\right)
            _{1\leq i\leq m,\,1\leq j\leq n}.
        \]
        When $m=n$, its determinant
        \[
            \Jac(\mathbf{F})(\mathbf{a})
            =
            \det\!\left(\mathbf{J}_{\mathbf{F}}(\mathbf{a})\right)
        \]
        is the Jacobian determinant.
    \end{definition}

    \begin{definition}[Differentiability; Total Derivative; Fr\'echet Derivative]
        Let $\mathbf{F}\colon U\subseteq\R^n\to\R^m$, where $U$ is open, and
        let $\mathbf{a}\in U$. The map $\mathbf{F}$ is differentiable at
        $\mathbf{a}$ if there exists a linear map
        $D\mathbf{F}(\mathbf{a}):\R^n\to\R^m$ s.t.
        \[
            \mathbf{F}(\mathbf{a}+\mathbf{h})
            =
            \mathbf{F}(\mathbf{a})
            +D\mathbf{F}(\mathbf{a})\mathbf{h}
            +\mathbf{r}(\mathbf{h}),
        \]
        where
        \[
            \lim\limits_{\mathbf{h}\to\mathbf{0}}
            \frac{\|\mathbf{r}(\mathbf{h})\|}{\|\mathbf{h}\|}
            =0.
        \]
        Equivalently, in compact symbolic form,
        {\small
        \[
            \forall\varepsilon>0\;\exists\delta>0\;\forall\mathbf{h}\in\R^n:\quad
            \bigl(\mathbf{a}+\mathbf{h}\in U\land0<\|\mathbf{h}\|<\delta\bigr)
            \Rightarrow
            \|\mathbf{F}(\mathbf{a}+\mathbf{h})-\mathbf{F}(\mathbf{a})
            -D\mathbf{F}(\mathbf{a})\mathbf{h}\|<\varepsilon\|\mathbf{h}\|.
        \]
        }
        The linear map $D\mathbf{F}(\mathbf{a})$ is the total derivative, or
        Fr\'echet derivative. Here linear means that it preserves linear
        combinations. In standard coordinates, its matrix is precisely the
        Jacobian matrix $\mathbf{J}_{\mathbf{F}}(\mathbf{a})$ defined above.
    \end{definition}


    \begin{definition}[Directional Derivative]
        Let $f\colon U\subseteq\R^n\to\R$ be differentiable at
        $\mathbf{a}\in U$, and let $\mathbf{u}$ be a unit vector. The
        directional derivative of $f$ at $\mathbf{a}$ in the direction
        $\mathbf{u}$ is
        \[
            D_{\mathbf{u}}f(\mathbf{a})
            =
            \lim\limits_{t\to0}
            \frac{f(\mathbf{a}+t\mathbf{u})-f(\mathbf{a})}{t}
            =
            \nabla f(\mathbf{a})\cdot\mathbf{u}.
        \]
        Consequently,
        \[
            |D_{\mathbf{u}}f(\mathbf{a})|
            \leq
            \|\nabla f(\mathbf{a})\|,
        \]
        with maximum increase in the direction of
        $\nabla f(\mathbf{a})$ when the gradient is nonzero.
    \end{definition}

    \begin{theorem}[Multivariable Chain Rule]
        Let $\mathbf{F}\colon \R^n\to\R^m$ and $\mathbf{G}\colon \R^m\to\R^k$ be
        differentiable. Then
        \[
            D(\mathbf{G}\circ\mathbf{F})(\mathbf{x})
            =
            D\mathbf{G}(\mathbf{F}(\mathbf{x}))D\mathbf{F}(\mathbf{x}).
        \]
    \end{theorem}

    \begin{theorem}[Fubini's Theorem]
        Let $R\subseteq\R^2$ be a bounded region of the form
        \[
            R=\{(x,y)\mid a\leq x\leq b,\ g_1(x)\leq y\leq g_2(x)\},
        \]
        where $g_1$ and $g_2$ are continuous and $g_1\leq g_2$. If $f$ is
        continuous on $R$, then
        \[
            \iint_R f(x,y)\dd A
            =
            \int_a^b\int_{g_1(x)}^{g_2(x)}f(x,y)\dd y\dd x.
        \]
        If the same region is also described as
        \[
            R=\{(x,y)\mid c\leq y\leq d,\ h_1(y)\leq x\leq h_2(y)\},
        \]
        then the integral may equally be evaluated in the reverse order.
        Measure-theoretic extensions allow weaker hypotheses; Tonelli's Theorem
        is introduced later in this chapter after the relevant measure concepts.
    \end{theorem}

    \begin{theorem}[Change of Variables Formula]
        Let $U,V\subseteq\R^n$ be open, and let
        $\mathbf{F}\colon U\to V$ be a $C^1$ diffeomorphism, meaning a bijection for
        which both $\mathbf{F}$ and $\mathbf{F}^{-1}$ are continuously
        differentiable. If $D\subseteq U$ is a region for which the integrals
        below exist and $f$ is integrable on $\mathbf{F}(D)$, then
        \[
            \int_{\mathbf{F}(D)}f(\mathbf{x})\dd\mathbf{x}
            =
            \int_D f(\mathbf{F}(\mathbf{u}))
            \left|\det D\mathbf{F}(\mathbf{u})\right|\dd\mathbf{u}.
        \]
        The absolute value is essential because the integral measures unsigned
        volume, independently of orientation.
    \end{theorem}

    \begin{remark}
        The Jacobian determinant is the local area or volume scaling factor of
        a change of variables.
    \end{remark}

    \begin{formula}[Polar Coordinates]
        For polar coordinates $(r,\theta)$,
        \[
            x=r\cos\theta,
            \qquad
            y=r\sin\theta,
            \qquad
            r\geq0,
            \qquad
            0\leq\theta<2\pi.
        \]
        The area element is
        \[
            \dd A=r\dd r\dd\theta.
        \]
        The area enclosed by a polar curve $r=r(\theta)$ is
        \[
            A
            =
            \frac12\int_{\alpha}^{\beta}r(\theta)^2\dd\theta.
        \]
        Its arc length is
        \[
            L
            =
            \int_{\alpha}^{\beta}
            \sqrt{r(\theta)^2+\left(r'(\theta)\right)^2}\dd\theta.
        \]
    \end{formula}

    \begin{formula}[Cylindrical and Spherical Coordinates]
        Cylindrical coordinates $(r,\theta,z)$ satisfy
        \[
            x=r\cos\theta,
            \qquad
            y=r\sin\theta,
            \qquad
            z=z,
        \]
        with volume element
        \[
            \dd V=r\dd r\dd\theta\dd z.
        \]
        Spherical coordinates $(\rho,\theta,\phi)$ satisfy
        \[
            x=\rho\sin\phi\cos\theta,
            \qquad
            y=\rho\sin\phi\sin\theta,
            \qquad
            z=\rho\cos\phi,
        \]
        where
        \[
            \rho\geq0,
            \qquad
            0\leq\theta<2\pi,
            \qquad
            0\leq\phi\leq\pi.
        \]
        The volume element is
        \[
            \dd V
            =
            \rho^2\sin\phi\dd\rho\dd\theta\dd\phi.
        \]
        Here $\theta$ is the azimuthal angle measured from the positive
        $x$-axis in the $xy$-plane, and $\phi$ is the polar angle measured from
        the positive $z$-axis.
    \end{formula}

    \begin{formula}[Differential Operators in Standard Coordinates]
        Let
        \[
            \mathbf{F}
            =F_r\mathbf{e}_r+F_\theta\mathbf{e}_\theta+F_z\mathbf{e}_z
        \]
        in cylindrical coordinates. Then
        \[
            \nabla f
            =
            \frac{\partial f}{\partial r}\mathbf{e}_r
            +\frac{1}{r}\frac{\partial f}{\partial\theta}\mathbf{e}_\theta
            +\frac{\partial f}{\partial z}\mathbf{e}_z,
        \]
        \[
            \nabla\cdot\mathbf{F}
            =
            \frac{1}{r}\frac{\partial}{\partial r}(rF_r)
            +\frac{1}{r}\frac{\partial F_\theta}{\partial\theta}
            +\frac{\partial F_z}{\partial z},
        \]
        and
        \[
            \Delta f
            =
            \frac{1}{r}\frac{\partial}{\partial r}
            \left(r\frac{\partial f}{\partial r}\right)
            +\frac{1}{r^2}\frac{\partial^2f}{\partial\theta^2}
            +\frac{\partial^2f}{\partial z^2}.
        \]

        In spherical coordinates $(\rho,\theta,\phi)$, where $\theta$ is the
        azimuthal angle and $\phi$ is the polar angle, write
        \[
            \mathbf{F}
            =F_\rho\mathbf{e}_\rho
            +F_\theta\mathbf{e}_\theta
            +F_\phi\mathbf{e}_\phi.
        \]
        Then
        \[
            \nabla f
            =
            \frac{\partial f}{\partial\rho}\mathbf{e}_\rho
            +\frac{1}{\rho\sin\phi}
             \frac{\partial f}{\partial\theta}\mathbf{e}_\theta
            +\frac{1}{\rho}
             \frac{\partial f}{\partial\phi}\mathbf{e}_\phi,
        \]
        \[
            \nabla\cdot\mathbf{F}
            =
            \frac{1}{\rho^2}\frac{\partial}{\partial\rho}(\rho^2F_\rho)
            +\frac{1}{\rho\sin\phi}\frac{\partial F_\theta}{\partial\theta}
            +\frac{1}{\rho\sin\phi}
             \frac{\partial}{\partial\phi}(\sin\phi\,F_\phi),
        \]
        and
        \[
            \Delta f
            =
            \frac{1}{\rho^2}\frac{\partial}{\partial\rho}
            \left(\rho^2\frac{\partial f}{\partial\rho}\right)
            +\frac{1}{\rho^2\sin^2\phi}
             \frac{\partial^2f}{\partial\theta^2}
            +\frac{1}{\rho^2\sin\phi}\frac{\partial}{\partial\phi}
            \left(\sin\phi\frac{\partial f}{\partial\phi}\right).
        \]
    \end{formula}

    \begin{theorem}[Inverse Function Theorem]
        Let $\mathbf{F}\colon \R^n\to\R^n$ be continuously differentiable near
        $\mathbf{a}$. If
        \[
            \det D\mathbf{F}(\mathbf{a})\neq 0,
        \]
        then $\mathbf{F}$ has a continuously differentiable local inverse near
        $\mathbf{a}$.
    \end{theorem}

    \begin{theorem}[Implicit Function Theorem]
        Let $F(x,y)$ be continuously differentiable near $(a,b)$, with
        \[
            F(a,b)=0,
            \qquad
            \frac{\partial F}{\partial y}(a,b)\neq 0.
        \]
        Then, near $(a,b)$, the equation $F(x,y)=0$ determines $y$ as a
        differentiable function of $x$, and
        \[
            \diff{y}{x}
            =
            -\frac{F_x}{F_y}.
        \]
    \end{theorem}

    \begin{formula}[Second Derivative Test and Hessian]
        A critical point of a differentiable function $f$ is a point where
        $\nabla f=\mathbf{0}$. For a twice differentiable function $f(x,y)$, the
        Hessian matrix is
        \[
            \mathbf{H}_f
            =
            \begin{pmatrix}
                f_{xx} & f_{xy} \\
                f_{yx} & f_{yy}
            \end{pmatrix}.
        \]
        At a critical point $(a,b)$, let
        \[
            \Delta_H
            =
            \det\mathbf{H}_f(a,b)
            =
            f_{xx}(a,b)f_{yy}(a,b)-f_{xy}(a,b)^2,
        \]
        where the last equality uses equality of the mixed partial derivatives
        under the usual continuity hypothesis.
        If $\Delta_H>0$ and $f_{xx}(a,b)>0$, the point is a local minimum. If
        $\Delta_H>0$ and $f_{xx}(a,b)<0$, it is a local maximum. If
        $\Delta_H<0$, it is a saddle point. If $\Delta_H=0$, the test is
        inconclusive.
    \end{formula}

    \begin{theorem}[Lagrange Multipliers]
        Suppose that $f$ has a local maximum or minimum subject to the constraint
        \[
            g(x_1,\dots,x_n)=0
        \]
        at a point where $\nabla g\neq\mathbf{0}$. Then there exists a scalar
        $\lambda$ s.t.
        \[
            \nabla f=\lambda\nabla g.
        \]
    \end{theorem}

    \begin{idea}
        At a constrained optimum, the level surface of $f$ is tangent to the
        constraint surface. Thus their normal vectors are parallel.
    \end{idea}

    \begin{theorem}[Euler's Homogeneous Function Theorem]
        Suppose $f\colon \R^n\to\R$ is differentiable and homogeneous of degree
        $k$, meaning
        \[
            f(t\mathbf{x})=t^k f(\mathbf{x})
        \]
        for $t>0$. Then
        \[
            \mathbf{x}\cdot\nabla f(\mathbf{x})=k f(\mathbf{x}).
        \]
        In coordinates,
        \[
            \sum_{i=1}^{n}x_i\frac{\partial f}{\partial x_i}=kf.
        \]
    \end{theorem}

    \begin{formula}[Parametrized Curves]
        Let $\mathbf{r}(t)$ be a twice differentiable curve. Its velocity,
        speed, and acceleration are
        \[
            \mathbf{v}(t)=\mathbf{r}'(t),
            \qquad
            \|\mathbf{v}(t)\|=\|\mathbf{r}'(t)\|,
            \qquad
            \mathbf{a}(t)=\mathbf{r}''(t).
        \]
        If $\mathbf{r}'(t_0)\neq\mathbf{0}$, the tangent line at $t_0$ is
        \[
            \mathbf{r}(t_0)+s\mathbf{r}'(t_0),
            \qquad
            s\in\R.
        \]
        The speed is constant if and only if
        \[
            \mathbf{r}'(t)\cdot\mathbf{r}''(t)=0.
        \]
    \end{formula}

    \begin{definition}[Frenet--Serret Apparatus]
        Let $\mathbf{r}(s)$ be a $C^3$ regular space curve parametrized by arc
        length, and suppose its curvature is nonzero. Define the unit tangent,
        principal normal, and binormal vectors by
        \[
            \mathbf{T}=\mathbf{r}'(s),
            \qquad
            \kappa=\|\mathbf{T}'(s)\|,
            \qquad
            \mathbf{N}=\frac{\mathbf{T}'(s)}{\kappa},
            \qquad
            \mathbf{B}=\mathbf{T}\times\mathbf{N}.
        \]
        The torsion is
        \[
            \tau=-\mathbf{B}'(s)\cdot\mathbf{N}(s).
        \]
        The collection
        \[
            \{\mathbf{T},\mathbf{N},\mathbf{B},\kappa,\tau\}
        \]
        is called the Frenet--Serret apparatus. It satisfies
        \[
            \mathbf{T}'=\kappa\mathbf{N},
            \qquad
            \mathbf{N}'=-\kappa\mathbf{T}+\tau\mathbf{B},
            \qquad
            \mathbf{B}'=-\tau\mathbf{N}.
        \]
    \end{definition}

    \begin{definition}[Scalar Line Integral]
        If a curve $C$ is parametrized by $\mathbf{r}(t)$ for $a\leq t\leq b$,
        then
        \[
            \int_C f\dd s
            =
            \int_a^b f(\mathbf{r}(t))\|\mathbf{r}'(t)\|\dd t.
        \]
        The value is unchanged by any regular reparametrization.
    \end{definition}

    \begin{definition}[Line Integral]
        If a curve $C$ is parametrized by $\mathbf{r}(t)$ for $a\leq t\leq b$,
        then the line integral of a vector field $\mathbf{F}$ along $C$ is
        \[
            \int_C \mathbf{F}\cdot\dd\mathbf{r}
            =
            \int_a^b \mathbf{F}(\mathbf{r}(t))\cdot\mathbf{r}'(t)\dd t.
        \]
    \end{definition}

    \begin{theorem}[Fundamental Theorem for Line Integrals; Gradient Theorem]
        If $\mathbf{F}=\nabla f$ and $C$ is a smooth curve from $\mathbf{a}$ to
        $\mathbf{b}$, then
        \[
            \int_C \mathbf{F}\cdot\dd\mathbf{r}
            =
            f(\mathbf{b})-f(\mathbf{a}).
        \]
        Hence the integral is path independent.
    \end{theorem}

    \begin{proposition}[Conservative-Field Test]
        Let $U\subseteq\R^n$ be open and simply connected, and let
        $\mathbf{F}\colon U\to\R^n$ be continuously differentiable. Then
        $\mathbf{F}$ is conservative if and only if its Jacobian matrix is
        symmetric:
        \[
            \frac{\partial F_i}{\partial x_j}
            =
            \frac{\partial F_j}{\partial x_i}
            \qquad(1\leq i,j\leq n).
        \]
        In $\R^3$, this is equivalent to $\nabla\times\mathbf{F}=\mathbf{0}$.
        Simple connectedness cannot in general be omitted.
    \end{proposition}

    \begin{theorem}[Green's Theorem]
        Let $D\subseteq\R^2$ be a bounded region whose positively oriented
        boundary $C=\partial D$ is piecewise smooth, and let $P,Q$ have
        continuous first partial derivatives on an open set containing $D$.
        Then
        \[
            \oint_C P\dd x+Q\dd y
            =
            \iint_D
            \left(
                \frac{\partial Q}{\partial x}
                -
                \frac{\partial P}{\partial y}
            \right)\dd A.
        \]
    \end{theorem}

    \begin{idea}
        Green's Theorem converts circulation along a boundary curve into curl
        over the region inside. It is the two-dimensional ancestor of Stokes'
        Theorem.
    \end{idea}


    \begin{definition}[Surface Area, Scalar Surface Integral, and Flux]
        Let a smooth oriented surface be parametrized by
        \[
            \mathbf{r}(u,v),
            \qquad
            (u,v)\in D.
        \]
        Its surface element is
        \[
            \dd S
            =
            \|\mathbf{r}_u\times\mathbf{r}_v\|\dd u\dd v.
        \]
        Hence
        \[
            \iint_S f\dd S
            =
            \iint_D
            f(\mathbf{r}(u,v))
            \|\mathbf{r}_u\times\mathbf{r}_v\|\dd u\dd v,
        \]
        and the flux of a vector field $\mathbf{F}$ through the chosen
        orientation is
        \[
            \iint_S\mathbf{F}\cdot\mathbf{n}\dd S
            =
            \iint_D
            \mathbf{F}(\mathbf{r}(u,v))
            \cdot
            (\mathbf{r}_u\times\mathbf{r}_v)\dd u\dd v.
        \]

        If $S$ is the graph $z=g(x,y)$ with upward orientation, then
        \[
            \dd S
            =
            \sqrt{1+g_x^2+g_y^2}\dd x\dd y,
        \]
        and
        \[
            \mathbf{n}\dd S
            =
            (-g_x,-g_y,1)\dd x\dd y.
        \]
        If instead $S$ is a regular level surface
        $G(x,y,z)=0$, then a unit normal is
        \[
            \mathbf{n}
            =
            \pm\frac{\nabla G}{\|\nabla G\|}.
        \]
    \end{definition}

    \begin{theorem}[Stokes' Theorem]
        Let $S$ be an oriented piecewise smooth surface with positively oriented
        piecewise smooth boundary $\partial S$, and let $\mathbf{F}$ be
        continuously differentiable on an open set containing $S$. Then
        \[
            \oint_{\partial S}\mathbf{F}\cdot\dd\mathbf{r}
            =
            \iint_S(\nabla\times\mathbf{F})\cdot\mathbf{n}\dd S.
        \]
        The orientation of $\partial S$ is determined from that of $S$ by the
        right-hand rule.
    \end{theorem}

    \begin{idea}
        Stokes' Theorem says that circulation around the boundary equals total
        curl through the surface. It turns a boundary integral into an interior
        integral.
    \end{idea}

    \begin{theorem}[Divergence Theorem; Gauss's Theorem]
        Let $V\subseteq\R^3$ be a bounded solid whose boundary
        $S=\partial V$ is piecewise smooth and oriented outward. If
        $\mathbf{F}$ is continuously differentiable on an open set containing
        $V$, then
        \[
            \iint_S\mathbf{F}\cdot\mathbf{n}\dd S
            =
            \iiint_V\nabla\cdot\mathbf{F}\dd V.
        \]
    \end{theorem}

    \begin{idea}
        The Divergence Theorem says that outward flux through the boundary
        equals total source strength inside the region.
    \end{idea}

    \begin{remark}
        Green's Theorem, Stokes' Theorem, and the Divergence Theorem share the
        same philosophy: a boundary integral is equal to an interior integral.
        This is the geometric heart of vector calculus.
    \end{remark}

    \section*{Real Analysis}

    \begin{theorem}[Least Upper Bound Property; Completeness Axiom of $\R$]
        Every nonempty subset of $\R$ that is bounded above has a least upper
        bound in $\R$. Equivalently, every nonempty subset bounded below has a
        greatest lower bound. This completeness property distinguishes $\R$
        from $\Q$ and underlies the standard convergence theorems of analysis.
    \end{theorem}

    \begin{definition}[Metric Space and Open Ball]
        A metric space is a pair $(X,d)$ in which $d\colon X\times X\to\R$ satisfies,
        for all $x,y,z\in X$,
        \[
            d(x,y)\geq 0,
            \qquad
            d(x,y)=0\Leftrightarrow x=y,
        \]
        \[
            d(x,y)=d(y,x),
            \qquad
            d(x,z)\leq d(x,y)+d(y,z).
        \]
        The open ball of radius $r>0$ centered at $x$ is
        \[
            B_r(x)=\{y\in X\mid d(x,y)<r\}.
        \]
    \end{definition}

    \begin{fact}[Common Metrics]
        For $\mathbf{x},\mathbf{y}\in\R^n$, each of
        \[
            d_1(\mathbf{x},\mathbf{y})
            =
            \sum_{k=1}^{n}|x_k-y_k|,
        \]
        \[
            d_2(\mathbf{x},\mathbf{y})
            =
            \left(
                \sum_{k=1}^{n}|x_k-y_k|^2
            \right)^{1/2},
            \qquad
            d_\infty(\mathbf{x},\mathbf{y})
            =
            \max_{1\leq k\leq n}|x_k-y_k|
        \]
        is a metric. On any set $X$, the discrete metric is
        \[
            d_{\mathrm{disc}}(x,y)
            =
            \begin{cases}
                0, & x=y,\\
                1, & x\neq y.
            \end{cases}
        \]
    \end{fact}

    \begin{definition}[Open and Closed Sets in a Metric Space]
        A subset $U\subseteq X$ is open if
        \[
            \forall x\in U,
            \quad
            \exists r>0,
            \qquad
            B_r(x)\subseteq U.
        \]
        A subset $F\subseteq X$ is closed if $X\setminus F$ is open.
    \end{definition}

    \begin{definition}[Point and Set Taxonomy in a Metric Space]
        Let $A\subseteq(X,d)$. An open cover of $A$ is a collection of open
        sets whose union contains $A$. The standard point types and associated
        sets are summarized below.

        {\small
        \renewcommand{\arraystretch}{1.16}
        \begin{longtable}{@{}L{0.22\textwidth}L{0.70\textwidth}@{}}
            \toprule
            Term & Criterion \\
            \midrule
            \endfirsthead

            \toprule
            Term & Criterion \\
            \midrule
            \endhead

            \bottomrule
            \endlastfoot

            Interior point & $x\in A$ and $\exists r>0$ s.t. $B_r(x)\subseteq A$ \\
            Exterior point & $\exists r>0$ s.t. $B_r(x)\subseteq X\setminus A$ \\
            Boundary point & $\forall r>0$, $B_r(x)\cap A\neq\varnothing$ and
                $B_r(x)\cap(X\setminus A)\neq\varnothing$ \\
            Accumulation point; limit point & $\forall r>0$,
                $(B_r(x)\setminus\{x\})\cap A\neq\varnothing$ \\
            Isolated point & $x\in A$ and $\exists r>0$ s.t.
                $B_r(x)\cap A=\{x\}$ \\
            Interior $A^\circ$ & the set of all interior points of $A$ \\
            Closure $\overline{A}$ & $A$ together with all of its accumulation points;
                equivalently, the smallest closed set containing $A$ \\
            Boundary $\partial A$ & $\overline{A}\setminus A^\circ$ \\
            Dense set & $\overline{A}=X$ \\
            Perfect set & $A$ is closed and has no isolated points; equivalently,
                every point of $A$ is an accumulation point of $A$ \\
            Bounded set & $A\subseteq B_R(x_0)$ for some $x_0\in X$ and $R>0$ \\
            Compact set & every open cover of $A$ has a finite subcover \\
        \end{longtable}
        }

        In a metric space, $A$ is open iff every point of $A$ is interior, and
        $A$ is closed iff it contains all of its accumulation points. In $\R^n$,
        compactness is equivalent to being closed and bounded by the
        Heine--Borel Theorem; this equivalence does not hold in arbitrary metric
        spaces.
    \end{definition}

    \begin{fact}[Standard Subsets of the Real Line]
        In the usual topology on $\R$:
        \begin{itemize}
            \item $\Z$ is closed but not open;
            \item $\Q$ is countable and is neither open nor closed;
            \item the Cantor set is nonempty, compact, perfect, uncountable, and
                  contains no nondegenerate interval.
        \end{itemize}
        These examples are standard tests of the distinctions among countability,
        openness, closedness, and perfectness.
    \end{fact}

    \begin{definition}[Limit of a Map Between Metric Spaces]
        Let $(X,d_X)$ and $(Y,d_Y)$ be metric spaces, let $D\subseteq X$, let
        $a$ be an accumulation point of $D$, and let $f\colon D\to Y$. The statement
        \[
            \lim\limits_{x\to a}f(x)=L
        \]
        means
        \[
            \forall\varepsilon>0,
            \quad
            \exists\delta>0,
            \quad
            \forall x\in D,
            \qquad
            0<d_X(x,a)<\delta
            \Rightarrow
            d_Y\bigl(f(x),L\bigr)<\varepsilon.
        \]
    \end{definition}

    \begin{definition}[Convergence of a Sequence]
        A sequence $(x_n)$ in a metric space $(X,d)$ converges to $x\in X$,
        written $x_n\to x$, if
        \[
            \forall\varepsilon>0,
            \quad
            \exists N\in\N,
            \quad
            \forall n\geq N,
            \qquad
            d(x_n,x)<\varepsilon.
        \]
    \end{definition}

    \begin{definition}[Cauchy Sequence and Complete Metric Space]
        A sequence $(x_n)$ in $(X,d)$ is Cauchy if
        \[
            \forall\varepsilon>0,
            \quad
            \exists N\in\N,
            \quad
            \forall m,n\geq N,
            \qquad
            d(x_m,x_n)<\varepsilon.
        \]
        The metric space is complete if every Cauchy sequence in $X$ converges
        to a point of $X$.
    \end{definition}

    \begin{definition}[Continuity in Metric Spaces]
        Let $(X,d_X)$ and $(Y,d_Y)$ be metric spaces. A map $f\colon X\to Y$ is
        continuous at $x_0\in X$ if
        \[
            \forall\varepsilon>0,
            \quad
            \exists\delta>0,
            \quad
            \forall x\in X,
            \qquad
            d_X(x,x_0)<\delta
            \Rightarrow
            d_Y\bigl(f(x),f(x_0)\bigr)<\varepsilon.
        \]
        It is continuous on $X$ if it is continuous at every $x_0\in X$.
    \end{definition}

    \begin{theorem}[Sequential Criterion for Continuity]
        A map $f\colon X\to Y$ between metric spaces is continuous at $x\in X$ if
        and only if
        \[
            x_n\to x
            \Rightarrow
            f(x_n)\to f(x)
        \]
        for every sequence $(x_n)$ in $X$.
    \end{theorem}

    \begin{definition}[Uniform Continuity]
        A map $f\colon (X,d_X)\to(Y,d_Y)$ is uniformly continuous if
        \[
            \forall\varepsilon>0,
            \quad
            \exists\delta>0,
            \quad
            \forall x,y\in X,
            \qquad
            d_X(x,y)<\delta
            \Rightarrow
            d_Y\bigl(f(x),f(y)\bigr)<\varepsilon.
        \]
        Unlike ordinary continuity, the same $\delta$ must work for every pair
        $x,y\in X$.
    \end{definition}

    \begin{definition}[Pointwise and Uniform Convergence]
        Let $f_n\colon E\to Y$ and $f\colon E\to Y$, where $(Y,d)$ is a metric space. The
        sequence $(f_n)$ converges pointwise to $f$, written
        $f_n\to f$, if
        \[
            \forall x\in E,
            \quad
            \forall\varepsilon>0,
            \quad
            \exists N=N(x,\varepsilon)\in\N,
            \quad
            \forall n\geq N,
            \qquad
            d\bigl(f_n(x),f(x)\bigr)<\varepsilon.
        \]
        It converges uniformly to $f$, written $f_n\rightrightarrows f$, if
        \[
            \forall\varepsilon>0,
            \quad
            \exists N=N(\varepsilon)\in\N,
            \quad
            \forall n\geq N,
            \quad
            \forall x\in E,
            \qquad
            d\bigl(f_n(x),f(x)\bigr)<\varepsilon.
        \]
        For real- or complex-valued functions, this is equivalent to
        \[
            \|f_n-f\|_\infty
            =
            \sup_{x\in E}|f_n(x)-f(x)|
            \longrightarrow 0.
        \]
        The convergence is locally uniform on $E$ if, for every compact set
        $K\subseteq E$, the restrictions $f_n|_K$ converge uniformly to $f|_K$.
        Equivalently, for real- or complex-valued functions,
        \[
            \sup_{x\in K}|f_n(x)-f(x)|
            \longrightarrow 0
        \]
        for every compact $K\subseteq E$.
    \end{definition}

    \begin{definition}[Equicontinuity and Uniform Boundedness]
        A family $\mathcal{F}$ of maps from $(X,d_X)$ to $(Y,d_Y)$ is
        equicontinuous at $x_0\in X$ if
        \[
            \forall\varepsilon>0,
            \quad
            \exists\delta>0,
            \quad
            \forall f\in\mathcal{F},
            \quad
            \forall x\in X,
            \qquad
            d_X(x,x_0)<\delta
            \Rightarrow
            d_Y\bigl(f(x),f(x_0)\bigr)<\varepsilon.
        \]
        A sequence of real-valued functions $(f_n)$ is uniformly bounded if
        \[
            \exists M>0,
            \quad
            \forall n\in\N,
            \quad
            \forall x\in X,
            \qquad
            |f_n(x)|\leq M.
        \]
    \end{definition}

    \begin{definition}[Limit Superior and Limit Inferior]
        For a real sequence $(a_n)$, define
        \[
            \limsup_{n\to\infty} a_n
            =
            \lim\limits_{n\to\infty}\sup_{k\geq n} a_k,
            \qquad
            \liminf_{n\to\infty} a_n
            =
            \lim\limits_{n\to\infty}\inf_{k\geq n} a_k,
        \]
        whenever these extended real limits are allowed.
    \end{definition}

    \begin{theorem}[Monotone Convergence Theorem for Sequences]
        Every bounded monotone real sequence converges. More precisely, every
        increasing sequence bounded above converges to its supremum, and every
        decreasing sequence bounded below converges to its infimum.
    \end{theorem}

    \begin{theorem}[Bolzano--Weierstrass Theorem]
        Every bounded sequence in $\R^n$ has a convergent subsequence.
    \end{theorem}

    \begin{theorem}[Heine--Borel Theorem]
        A subset of $\R^n$ is compact if and only if it is closed and bounded.
    \end{theorem}

    \begin{theorem}[Heine--Cantor Theorem; Uniform Continuity on Compact Sets]
        If $K\subseteq\R^n$ is compact and $f\colon K\to\R^m$ is continuous, then $f$
        is uniformly continuous on $K$.
    \end{theorem}

    \begin{strategy}
        In analysis questions, compactness usually means that one may extract a
        convergent subsequence or that a continuous function attains extrema.
        In $\R^n$, the fast recognition phrase is ``closed and bounded.''
    \end{strategy}

    \begin{theorem}[Cauchy Criterion for Convergence]
        In a complete metric space, a sequence converges if and only if it is
        Cauchy. In particular, for $(a_n)$ in $\R$ or $\C$,
        \[
            (a_n)\text{ converges}
            \Leftrightarrow
            \forall\varepsilon>0,
            \ \exists N\in\N,
            \ \forall m,n\geq N,
            \ |a_m-a_n|<\varepsilon.
        \]
    \end{theorem}

    \begin{theorem}[Uniform Cauchy Criterion]
        Let $(Y,d)$ be complete and let $f_n\colon E\to Y$. Then $(f_n)$ converges
        uniformly on $E$ if and only if
        \[
            \forall\varepsilon>0,
            \quad
            \exists N\in\N,
            \quad
            \forall m,n\geq N,
            \quad
            \forall x\in E,
            \qquad
            d\bigl(f_m(x),f_n(x)\bigr)<\varepsilon.
        \]
    \end{theorem}

    \begin{theorem}[Weierstrass M-Test]
        Let $(f_n)$ be a sequence of functions on a set $E$. If
        \[
            |f_n(x)|\leq M_n
            \quad\text{for all }x\in E
        \]
        and
        \[
            \sum_{n=1}^{\infty}M_n
        \]
        converges, then the series
        \[
            \sum_{n=1}^{\infty} f_n(x)
        \]
        converges uniformly on $E$.
    \end{theorem}

    \begin{theorem}[Uniform Limit Theorem]
        If $(f_n)$ is a sequence of real- or complex-valued continuous functions
        on a metric space $E$ and $f_n\rightrightarrows f$ on $E$, then $f$ is continuous
        on $E$.
    \end{theorem}

    \begin{theorem}[Term-by-Term Differentiation]
        Suppose $(f_n)$ is a sequence of continuously differentiable functions
        on a closed bounded interval $[a,b]$, and suppose $(f_n(x_0))$
        converges for some $x_0\in[a,b]$. If $f_n'\rightrightarrows g$ on
        $[a,b]$, then $f_n\rightrightarrows f$ for a differentiable function
        $f$, and
        \[
            f'=g=\lim\limits_{n\to\infty} f_n'.
        \]
    \end{theorem}

    \begin{theorem}[Arzelà--Ascoli Theorem]
        Every uniformly bounded and equicontinuous sequence of real-valued
        continuous functions on a compact metric space has a uniformly
        convergent subsequence.
    \end{theorem}

    \begin{theorem}[Stone--Weierstrass Theorem]
        A subalgebra of $C(K,\R)$ that contains the constants and separates
        points is dense in $C(K,\R)$ under the supremum norm, provided $K$ is
        compact.
    \end{theorem}

    \begin{theorem}[Banach Fixed-Point Theorem; Contraction Mapping Theorem]
        Let $(X,d)$ be a nonempty complete metric space, and suppose that
        $T\colon X\to X$ is a contraction: there exists $q\in[0,1)$ s.t.
        \[
            \forall x,y\in X,
            \qquad
            d\bigl(T(x),T(y)\bigr)\leq qd(x,y).
        \]
        Then $T$ has a unique fixed point $x^\ast\in X$, and the iteration
        $x_{n+1}=T(x_n)$ converges to $x^\ast$ for every initial point $x_0\in X$.
    \end{theorem}

    \begin{theorem}[Baire Category Theorem]
        A nonempty complete metric space cannot be written as a countable union
        of nowhere dense closed sets. Equivalently, a countable intersection of
        open dense subsets of a complete metric space is dense.
    \end{theorem}

    \begin{definition}[Sigma-Algebra, Measure, Measurability, and Almost-Everywhere Convergence]
        A $\sigma$-algebra $\mathcal{F}$ on a set $X$ is a collection of subsets
        containing $X$ and closed under complements and countable unions. A
        measure is a function $\mu\colon \mathcal{F}\to[0,\infty]$ with
        $\mu(\varnothing)=0$ that is countably additive on pairwise disjoint
        sets. The triple $(X,\mathcal{F},\mu)$ is a measure space.

        A function $f\colon X\to\R$ is measurable if the inverse image of every Borel
        set is in $\mathcal{F}$. A property holds almost everywhere, abbreviated
        a.e., if it fails only on a set of measure zero. Thus $f_n\aeto f$ means
        that $f_n(x)\to f(x)$ for almost every $x$.
    \end{definition}

    \begin{theorem}[Tonelli's Theorem]
        Let $(X,\mu)$ and $(Y,\nu)$ be $\sigma$-finite measure spaces, and let
        $f\colon X\times Y\to[0,\infty]$ be measurable. Then the iterated integrals
        and the product-space integral agree, possibly with value $+\infty$:
        \[
            \int_{X\times Y}f\dd(\mu\times\nu)
            =
            \int_X\left(\int_Y f(x,y)\dd\nu(y)\right)\dd\mu(x)
            =
            \int_Y\left(\int_X f(x,y)\dd\mu(x)\right)\dd\nu(y).
        \]
        In particular, a nonnegative series may be integrated term by term.
    \end{theorem}

    \begin{theorem}[Fatou's Lemma]
        If $f_n\geq 0$ are measurable, then
        \[
            \int \liminf_{n\to\infty}f_n\dd\mu
            \leq
            \liminf_{n\to\infty}\int f_n\dd\mu.
        \]
    \end{theorem}

    \begin{theorem}[Lebesgue Monotone Convergence Theorem; Beppo Levi Theorem]
        If $0\leq f_1\leq f_2\leq\cdots$ and $f_n\aeto f$,
        then
        \[
            \lim\limits_{n\to\infty}\int f_n\dd\mu
            =
            \int f\dd\mu.
        \]
    \end{theorem}

    \begin{theorem}[Lebesgue Dominated Convergence Theorem]
        If $f_n\aeto f$ and there exists an integrable function
        $g$ s.t. $|f_n|\leq g$ for all $n$, then
        \[
            \lim\limits_{n\to\infty}\int f_n\dd\mu
            =
            \int f\dd\mu.
        \]
    \end{theorem}

    \begin{theorem}[Density of Irrational Rotations]
        If $\alpha\in\R\setminus\Q$, then the sequence of fractional parts
        \[
            \{n\alpha\},\qquad n=0,1,2,\dots,
        \]
        is dense in $[0,1]$. Equivalently, the orbit of any point under rotation
        of the circle through an irrational multiple of $2\pi$ is dense in the
        circle.
    \end{theorem}

    \begin{idea}
        Write $\lVert x\rVert_{\R/\Z}$ for the distance from $x$ to the nearest
        integer. By the pigeonhole principle, for every $N\geq2$ there is a
        positive integer $q$ s.t.
        \[
            0<\lVert q\alpha\rVert_{\R/\Z}<\frac1N.
        \]
        Put $\delta=\lVert q\alpha\rVert_{\R/\Z}$ and
        $M=\lfloor1/\delta\rfloor$. Up to reversing orientation on the circle,
        the integer orbit contains
        \[
            0,\delta,2\delta,\dots,M\delta,
        \]
        whose successive circular gaps are at most $\delta<1/N$. Thus the
        two-sided orbit $\{n\alpha\bmod1\mid n\in\Z\}$ is dense.

        Moreover, irrationality gives an unbounded sequence of positive
        integers $q_j$ for which
        \[
            \lVert q_j\alpha\rVert_{\R/\Z}\longrightarrow0.
        \]
        For each fixed $m\geq1$, the nonnegative multiples
        $(q_j-m)\alpha$ converge modulo $1$ to $-m\alpha$, and $q_j-m\geq0$
        for all sufficiently large $j$. Hence the closure of the nonnegative
        orbit contains every negative multiple as well as every nonnegative
        multiple. It therefore contains the dense two-sided orbit, proving the
        claim.
    \end{idea}

    \section*{Complex Analysis}

    \begin{definition}[Complex Differentiability; Holomorphic and Entire Functions]
        Let $U\subseteq\C$ be open and let $f\colon U\to\C$. The complex derivative of
        $f$ at $z_0\in U$ is
        \[
            f'(z_0)
            =
            \lim\limits_{z\to z_0}
            \frac{f(z)-f(z_0)}{z-z_0},
        \]
        provided the limit exists independently of the direction of approach.
        The function is holomorphic on $U$ if it is complex differentiable at
        every point of $U$, and entire if it is holomorphic on all of $\C$.
    \end{definition}

    \begin{theorem}[Cauchy--Riemann Equations]
        Write
        \[
            f(z)=u(x,y)+iv(x,y),
            \qquad
            z=x+iy.
        \]
        If $f$ is complex differentiable at $z_0=(x_0,y_0)$ and the first partial
        derivatives exist there, then
        \[
            \frac{\partial u}{\partial x}(x_0,y_0)
            =
            \frac{\partial v}{\partial y}(x_0,y_0),
            \qquad
            \frac{\partial u}{\partial y}(x_0,y_0)
            =
            -\frac{\partial v}{\partial x}(x_0,y_0).
        \]
        Conversely, if these partial derivatives are continuous near $z_0$ and
        satisfy the equations there, then $f$ is complex differentiable at
        $z_0$.
    \end{theorem}

    \begin{formula}[Cauchy--Riemann Equations in Polar Coordinates]
        Write
        \[
            f(re^{i\theta})=u(r,\theta)+iv(r,\theta),
            \qquad
            r>0.
        \]
        If the relevant first partial derivatives are continuous, then the
        Cauchy--Riemann equations become
        \[
            \frac{\partial u}{\partial r}
            =
            \frac{1}{r}\frac{\partial v}{\partial\theta},
            \qquad
            \frac{\partial v}{\partial r}
            =
            -\frac{1}{r}\frac{\partial u}{\partial\theta}.
        \]
    \end{formula}

    \begin{strategy}
        To test whether a given $f(z)=u(x,y)+iv(x,y)$ can be holomorphic, compute
        the four first partial derivatives and check the Cauchy--Riemann equations.
        Failure at one point immediately rules out holomorphicity there.
    \end{strategy}

    \begin{fact}[Standard Holomorphic and Nonholomorphic Examples]
        Every polynomial, $e^z$, $\sin z$, and $\cos z$ is entire. The function
        $\tan z$ is meromorphic but not entire because it has poles at
        \[
            z=\frac{\pi}{2}+k\pi,
            \qquad
            k\in\Z.
        \]
        The function $z\mapsto\Real(z)$ is not complex differentiable at any
        point, since its difference quotient depends on the direction of
        approach.
    \end{fact}

    \begin{proposition}[Holomorphic Functions Have Harmonic Components]
        If $f=u+iv$ is holomorphic and $u,v$ have continuous second partial
        derivatives, then
        \[
            \Delta u=0,
            \qquad
            \Delta v=0.
        \]
        Thus the real and imaginary parts of a holomorphic function are harmonic
        conjugates locally.
    \end{proposition}

    \begin{formula}[Complex Exponential, Trigonometric Functions, and Logarithm]
        For $z=x+iy$,
        \[
            e^z=e^x(\cos y+i\sin y).
        \]
        The complex sine and cosine satisfy
        \[
            \sin z=\frac{e^{iz}-e^{-iz}}{2i},
            \qquad
            \cos z=\frac{e^{iz}+e^{-iz}}{2}.
        \]
        For $z\neq0$, the complex logarithm is multivalued:
        \[
            \log z
            =
            \ln|z|+i\bigl(\arg z+2\pi k\bigr),
            \qquad
            k\in\Z.
        \]
        A single-valued holomorphic logarithm requires a choice of branch on a
        domain that does not wind around the origin.
    \end{formula}

    \begin{definition}[Complex Contour Integral]
        Let $\gamma\colon [a,b]\to\C$ be a piecewise $C^1$ parametrized contour and
        let $f$ be continuous on the image of $\gamma$. The contour integral is
        defined by
        \[
            \int_\gamma f(z)\dd z
            =
            \int_a^b f(\gamma(t))\gamma'(t)\dd t.
        \]
        If $\gamma$ is closed, the notation $\oint_\gamma f(z)\dd z$ is
        commonly used. Reversing the orientation changes the sign of the
        integral.
    \end{definition}

    \begin{theorem}[Cauchy--Goursat Integral Theorem]
        If $f$ is holomorphic on a simply connected domain $U\subseteq\C$, then
        for every closed piecewise smooth curve $\gamma$ in $U$,
        \[
            \oint_{\gamma}f(z)\dd z=0.
        \]
        Equivalently, contour integrals of $f$ in $U$ are path-independent.
    \end{theorem}

    \begin{theorem}[Cauchy Integral Formula]
        Let $f$ be holomorphic on an open set containing a positively oriented
        simple closed contour $\gamma$ and its interior. If $a$ lies inside
        $\gamma$, then
        \[
            f(a)
            =
            \frac{1}{2\pi i}
            \oint_{\gamma}\frac{f(z)}{z-a}\dd z.
        \]
        More generally, for every integer $n\geq 0$,
        \[
            f^{(n)}(a)
            =
            \frac{n!}{2\pi i}
            \oint_{\gamma}
            \frac{f(z)}{(z-a)^{n+1}}\dd z.
        \]
    \end{theorem}

    \begin{strategy}
        When a contour integral contains $(z-a)^{-m}$ and the remaining factor
        is holomorphic, compare it directly with Cauchy's integral formula before
        attempting a parametrization of the contour.
    \end{strategy}

    \begin{theorem}[Maximum Modulus Principle]
        A nonconstant holomorphic function on a connected open set cannot attain
        a maximum of $|f|$ at an interior point. Consequently, if $f$ is
        holomorphic on a bounded domain and continuous on its closure, then the
        maximum of $|f|$ occurs on the boundary.
    \end{theorem}

    \begin{theorem}[Liouville's Theorem]
        Every bounded entire function is constant.
    \end{theorem}

    \begin{remark}
        Liouville's Theorem gives a short proof of the Fundamental Theorem of
        Algebra: if a nonconstant complex polynomial had no zero, then its
        reciprocal would be a bounded entire function.
    \end{remark}

    \begin{theorem}[Taylor Series for Holomorphic Functions]
        If $f$ is holomorphic on the open disk $|z-a|<R$, then
        \[
            f(z)
            =
            \sum_{n=0}^{\infty}
            \frac{f^{(n)}(a)}{n!}(z-a)^n,
            \qquad
            |z-a|<R.
        \]
        Thus a holomorphic function is analytic: complex differentiability implies
        local representability by a convergent power series.
    \end{theorem}

    \begin{theorem}[Laurent Series]
        If $f$ is holomorphic on an annulus
        \[
            r<|z-a|<R,
        \]
        then it has a unique Laurent expansion
        \[
            f(z)
            =
            \sum_{n=-\infty}^{\infty}c_n(z-a)^n
        \]
        converging throughout the annulus. The coefficient $c_{-1}$ is called the
        residue of $f$ at $a$.
    \end{theorem}

    \begin{definition}[Isolated Singularities and Residues]
        A point $a$ is an isolated singularity of $f$ if $f$ is holomorphic on
        some punctured disk $0<|z-a|<R$ but is not holomorphic at $a$. In the
        Laurent expansion about $a$, the residue is
        \[
            \Res(f;a)=c_{-1}.
        \]
        The singularity is removable if every negative-power coefficient is zero,
        a pole if only finitely many negative-power coefficients are nonzero, and
        essential if infinitely many are nonzero.
    \end{definition}

    \begin{formula}[Residues at Poles]
        If $a$ is a simple pole of $f$, then
        \[
            \Res(f;a)
            =
            \lim\limits_{z\to a}(z-a)f(z).
        \]
        More generally, if $a$ is a pole of order $m$, then
        \[
            \Res(f;a)
            =
            \frac{1}{(m-1)!}
            \lim\limits_{z\to a}
            \frac{\dd^{m-1}}{\dd z^{m-1}}
            \left((z-a)^m f(z)\right).
        \]
        If $f(z)=g(z)/h(z)$, where $g$ and $h$ are holomorphic,
        $h(a)=0$, and $h'(a)\neq 0$, then
        \[
            \Res(f;a)=\frac{g(a)}{h'(a)}.
        \]
    \end{formula}

    \begin{theorem}[Residue Theorem]
        Let $f$ be holomorphic inside and on a positively oriented simple closed
        contour $\gamma$, except for isolated singularities
        $a_1,\dots,a_m$ inside $\gamma$. Then
        \[
            \oint_{\gamma}f(z)\dd z
            =
            2\pi i
            \sum_{k=1}^{m}\Res(f;a_k).
        \]
    \end{theorem}

    \begin{strategy}[Real Integrals by Residues]
        For a rational integral over $\R$, extend the integrand to a complex
        rational function, choose a semicircular contour, sum the residues in the
        appropriate half-plane, and verify that the contribution from the large
        arc vanishes. Trigonometric integrals over $[0,2\pi]$ often become contour
        integrals under the substitution $z=e^{i\theta}$ and
        \[
            \dd\theta=\frac{\dd z}{iz}.
        \]
    \end{strategy}

    \begin{theorem}[Rouch\'e's Theorem]
        Let $f$ and $g$ be holomorphic inside and on a simple closed contour
        $\gamma$. If
        \[
            |g(z)|<|f(z)|
            \qquad
            \text{for every }z\in\gamma,
        \]
        then $f$ and $f+g$ have the same number of zeros inside $\gamma$, counted
        with multiplicity.
    \end{theorem}

    \begin{definition}[Conformal Map]
        A holomorphic function $f$ is conformal at $z_0$ if $f'(z_0)\neq 0$.
        At such a point, $f$ preserves oriented angles and infinitesimal shapes,
        although it may change scale.
    \end{definition}

    \begin{definition}[M\"obius Transformation; Linear Fractional Transformation]
        A M\"obius transformation is a map of the extended complex plane of the
        form
        \[
            T(z)=\frac{az+b}{cz+d},
            \qquad
            ad-bc\neq 0.
        \]
        It is conformal wherever its derivative is finite and nonzero, and it maps
        generalized circles---circles and lines---to generalized circles.
    \end{definition}

    \begin{formula}[Standard Conformal Maps]
        The following transformations are useful to recognize quickly:
        \[
            z\mapsto az+b
            \quad\text{(rotation, dilation, and translation)},
        \]
        \[
            z\mapsto \frac{1}{z}
            \quad\text{(inversion)},
            \qquad
            z\mapsto z^n
            \quad\text{(angle multiplication)},
        \]
        and
        \[
            z\mapsto \frac{z-i}{z+i},
        \]
        which maps the upper half-plane conformally onto the unit disk.
    \end{formula}

    \chapter{Linear Algebra}

    \relatedcourses{(MATH203) Applied Linear Algebra}

    \relatedbooks{Linear Algebra and Its Applications (Gilbert Strang)}

    \begin{definition}[Vector Space]
        A vector space $V$ over a field $F$ is a set equipped with vector
        addition and scalar multiplication satisfying the standard vector-space
        axioms. The scalars belong to $F$; in this book the principal cases are
        $F=\R$ and $F=\C$.
    \end{definition}

    \begin{definition}[Linear Combination, Span, and Linear Independence]
        A vector of the form
        \[
            c_1\mathbf{v}_1+\cdots+c_k\mathbf{v}_k
        \]
        is a linear combination of $\mathbf{v}_1,\dots,\mathbf{v}_k$. Their
        span is
        \[
            \Span(\mathbf{v}_1,\dots,\mathbf{v}_k)
            =
            \left\{
                c_1\mathbf{v}_1+\cdots+c_k\mathbf{v}_k
                :c_1,\dots,c_k\in F
            \right\}.
        \]
        The vectors are linearly independent if
        \[
            c_1\mathbf{v}_1+\cdots+c_k\mathbf{v}_k=\mathbf{0}
        \]
        implies $c_1=\cdots=c_k=0$.
    \end{definition}

    \begin{definition}[Basis and Dimension]
        A list
        \[
            \mathbf{v}_1,\dots,\mathbf{v}_n
        \]
        is a basis of $V$ if it is linearly independent and spans $V$. If $V$
        has a finite basis, every basis has the same number of vectors; this
        number is the dimension of $V$, denoted by $\dim V$.
    \end{definition}

    \begin{theorem}[Basis Extension and Reduction]
        Every linearly independent list in a finite-dimensional vector space
        can be extended to a basis, and every spanning list contains a basis.
        More generally, the assertion that every vector space has a basis is
        equivalent to a standard form of the axiom of choice.
    \end{theorem}

    \begin{definition}[Linear Map; Linear Transformation]
        Let $V$ and $W$ be vector spaces over the same field $F$. A map
        $T\colon V\to W$ is linear if
        \[
            \forall\alpha,\beta\in F,
            \quad
            \forall\mathbf{u},\mathbf{v}\in V,
            \qquad
            T(\alpha\mathbf{u}+\beta\mathbf{v})
            =
            \alpha T(\mathbf{u})+\beta T(\mathbf{v}).
        \]
        Equivalently, $T$ preserves vector addition and scalar
        multiplication. Its kernel and image are
        \[
            \ker T=\{\mathbf{v}\in V\mid T(\mathbf{v})=\mathbf{0}\},
            \qquad
            \image T=\{T(\mathbf{v})\mid\mathbf{v}\in V\}.
        \]
    \end{definition}


    \begin{formula}[Coordinate Vectors, Matrix Representation, and Change of Basis]
        Let
        \[
            \mathcal{B}=(\mathbf{v}_1,\dots,\mathbf{v}_n)
        \]
        be an ordered basis of $V$. The coordinate vector of
        $\mathbf{v}\in V$ relative to $\mathcal{B}$ is the unique column vector
        $[\mathbf{v}]_{\mathcal{B}}$ satisfying
        \[
            \mathbf{v}
            =
            \begin{pmatrix}
                \mathbf{v}_1 & \cdots & \mathbf{v}_n
            \end{pmatrix}
            [\mathbf{v}]_{\mathcal{B}}.
        \]
        If $T\colon V\to W$ is linear and $\mathcal{C}$ is an ordered basis of $W$,
        then the matrix $[T]_{\mathcal{C}\leftarrow\mathcal{B}}$ is characterized by
        \[
            [T(\mathbf{v})]_{\mathcal{C}}
            =
            [T]_{\mathcal{C}\leftarrow\mathcal{B}}
            [\mathbf{v}]_{\mathcal{B}}.
        \]

        If $\mathcal{B}$ and $\mathcal{B}'$ are ordered bases of the same
        $n$-dimensional space and
        \[
            [\mathbf{v}]_{\mathcal{B}}
            =
            \mathbf{Q}[\mathbf{v}]_{\mathcal{B}'},
        \]
        then $\mathbf{Q}$ is invertible. For a linear operator $T\colon V\to V$,
        \[
            [T]_{\mathcal{B}'}
            =
            \mathbf{Q}^{-1}[T]_{\mathcal{B}}\mathbf{Q}.
        \]
        Thus matrices representing the same linear operator in different bases
        are similar.
    \end{formula}


    \begin{definition}[Column Space, Row Space, Null Space, and Rank]
        Let $\mathbf{A}\in F^{m\times n}$. Its column space and row space are
        the subspaces spanned by its columns and rows:
        \[
            \Col(\mathbf{A})\subseteq F^m,
            \qquad
            \Row(\mathbf{A})\subseteq F^n.
        \]
        Its null space is
        \[
            \Null(\mathbf{A})
            =
            \{\mathbf{x}\in F^n\mid \mathbf{A}\mathbf{x}=\mathbf{0}\}.
        \]
        The rank is
        \[
            \rk(\mathbf{A})=\dim\Col(\mathbf{A}).
        \]
        Equivalently, rank is the number of pivot columns in any row-echelon
        form of $\mathbf{A}$; the Row Rank Equals Column Rank theorem below
        justifies the corresponding row-space interpretation.
    \end{definition}

    \begin{definition}[Elementary Row Operations, Echelon Forms, and Pivots]
        The elementary row operations are: interchange two rows; multiply a row
        by a nonzero scalar; and add a scalar multiple of one row to another.
        A matrix is in \emph{row echelon form} if every nonzero row begins to the
        right of the row above it and every zero row is below every nonzero row.
        The first nonzero entry of a nonzero row is a \emph{pivot}. It is in
        \emph{reduced row echelon form} if, in addition, every pivot is $1$ and
        is the only nonzero entry in its column.
    \end{definition}

    \begin{algorithm}[Gaussian Elimination and Gauss--Jordan Reduction]
        For a linear system $\mathbf{A}\mathbf{x}=\mathbf{b}$, form the
        augmented matrix $[\mathbf{A}\mid\mathbf{b}]$ and apply elementary row
        operations.
        \begin{bookitems}
            \item Reduce to row echelon form to solve by back-substitution, or
                  continue to reduced row echelon form for a direct description
                  of all solutions.
            \item Pivot columns identify basic variables; nonpivot columns
                  identify free variables. The number of pivots is
                  $\rk(\mathbf{A})$.
            \item A row of the form
                  $[\mathbf{0}^{\mathsf T}\mid c]$ with $c\neq0$ means that the
                  system is inconsistent. Otherwise, there is a unique solution
                  exactly when every variable column is a pivot column; free
                  variables produce infinitely many solutions over an infinite
                  field.
            \item To find $\mathbf{A}^{-1}$ for a square matrix, row-reduce
                  $[\mathbf{A}\mid\mathbf{I}_n]$. The matrix is invertible
                  exactly when the left block reduces to $\mathbf{I}_n$, in
                  which case the right block becomes $\mathbf{A}^{-1}$.
        \end{bookitems}
        In a timed problem, exploit zeros, symmetry, triangular or block form,
        and obvious linear relations before carrying out a full reduction.
    \end{algorithm}

    \begin{fact}[Translation Is Affine, Not Linear]
        For a fixed vector $\mathbf{v}\in\R^n$, the translation
        \[
            T_{\mathbf{v}}(\mathbf{x})
            =
            \mathbf{x}+\mathbf{v}
        \]
        is bijective, with inverse $T_{-\mathbf{v}}$. It is a linear map if and
        only if $\mathbf{v}=\mathbf{0}$.
    \end{fact}

    \begin{definition}[Inner Product, Norm, Orthogonality, and Orthonormality]
        For real column vectors, the standard inner product is
        \[
            \langle\mathbf{u},\mathbf{v}\rangle
            =
            \mathbf{u}^{\mathsf T}\mathbf{v},
        \]
        while over $\C$ it is
        \[
            \langle\mathbf{u},\mathbf{v}\rangle
            =
            \mathbf{u}^{\mathsf H}\mathbf{v}.
        \]
        The induced norm is
        \[
            \|\mathbf{v}\|=\sqrt{\langle\mathbf{v},\mathbf{v}\rangle}.
        \]
        Vectors $\mathbf{u}$ and $\mathbf{v}$ are orthogonal if
        $\langle\mathbf{u},\mathbf{v}\rangle=0$. A list is orthonormal if its
        vectors are pairwise orthogonal and each has norm $1$.
    \end{definition}

    \begin{formula}[Adjoint Transfer]
        For a real matrix $\mathbf{A}$ of compatible size,
        \[
            \langle\mathbf{x},\mathbf{A}\mathbf{y}\rangle
            =
            \langle\mathbf{A}^{\mathsf T}\mathbf{x},\mathbf{y}\rangle.
        \]
        Hence, if $\mathbf{A}$ is symmetric,
        \[
            \langle\mathbf{x},\mathbf{A}\mathbf{y}\rangle
            =
            \langle\mathbf{A}\mathbf{x},\mathbf{y}\rangle.
        \]
        Over $\C$, replace $\mathbf{A}^{\mathsf T}$ by
        $\mathbf{A}^{\mathsf H}$.
    \end{formula}

    \begin{theorem}[Rank--Nullity Theorem]
        Let $T\colon V\to W$ be a linear map between finite-dimensional vector spaces.
        Then
        \[
            \dim V
            =
            \dim\ker T+\dim\image T.
        \]
        For a matrix $\mathbf{A}\in F^{m\times n}$, this becomes
        \[
            n=\rk(\mathbf{A})+\dim\Null(\mathbf{A}).
        \]
    \end{theorem}

    \begin{strategy}
        Rank--Nullity connects variables, pivots, and free variables. It is the
        linear algebra version of accounting: nothing disappears; dimensions
        are redistributed.
    \end{strategy}

    \begin{theorem}[Row Rank Equals Column Rank]
        For every matrix $\mathbf{A}$,
        \[
            \dim\Row(\mathbf{A})
            =
            \dim\Col(\mathbf{A}).
        \]
        This common dimension is called the rank of $\mathbf{A}$ and is denoted
        by $\rk(\mathbf{A})$.
    \end{theorem}

    \begin{definition}[Determinant]
        The determinant is a scalar associated to a square matrix
        $\mathbf{A}\in F^{n\times n}$, denoted by $\det(\mathbf{A})$.
        Geometrically, for real matrices, $|\det(\mathbf{A})|$ is the volume
        scaling factor of the linear transformation
        $\mathbf{x}\mapsto\mathbf{A}\mathbf{x}$.
    \end{definition}

    \begin{definition}[Minor, Cofactor, and Adjugate]
        Let $\mathbf{A}=(a_{ij})\in F^{n\times n}$. The minor $M_{ij}$ is the
        determinant of the $(n-1)\times(n-1)$ matrix obtained by deleting row
        $i$ and column $j$. The corresponding cofactor is
        \[
            C_{ij}=(-1)^{i+j}M_{ij}.
        \]
        The cofactor matrix is $(C_{ij})$, and the adjugate is its transpose:
        \[
            \adj(\mathbf{A})=(C_{ij})^{\mathsf T}.
        \]
        Cofactors are the coefficients in Laplace expansion and are therefore
        the bridge between determinants and the adjugate formula for an inverse.
    \end{definition}

    \begin{theorem}[Invertible Matrix Theorem]
        For a square matrix $\mathbf{A}\in F^{n\times n}$, the following
        statements are equivalent:
        \[
            \mathbf{A}\text{ is invertible},
            \qquad
            \det(\mathbf{A})\neq 0,
            \qquad
            \rk(\mathbf{A})=n,
        \]
        \[
            \Null(\mathbf{A})=\{\mathbf{0}\},
        \]
        and
        \[
            \mathbf{A}\mathbf{x}=\mathbf{b}
        \]
        has a unique solution for every $\mathbf{b}\in F^n$.
    \end{theorem}

    \begin{remark}
        This theorem is a dictionary. It translates between determinants, rank,
        null spaces, and solvability of linear systems.
    \end{remark}

    \begin{proposition}[One-Sided Inverse for Square Matrices]
        If $\mathbf{A},\mathbf{B}\in F^{n\times n}$ satisfy
        \[
            \mathbf{A}\mathbf{B}=\mathbf{I}_n,
        \]
        then both matrices are invertible and
        \[
            \mathbf{B}\mathbf{A}=\mathbf{I}_n.
        \]
        Thus, for square matrices, a right inverse is automatically a left
        inverse, and conversely.
    \end{proposition}

    \begin{formula}[Inverse of a $2\times 2$ Matrix]
        If
        \[
            \mathbf{A}
            =
            \begin{pmatrix}
                a & b \\
                c & d
            \end{pmatrix}
        \]
        and $ad-bc\neq 0$, then
        \[
            \mathbf{A}^{-1}
            =
            \frac{1}{ad-bc}
            \begin{pmatrix}
                d & -b \\
                -c & a
            \end{pmatrix}.
        \]
    \end{formula}

    \begin{theorem}[Cramer's Rule]
        Let $\mathbf{A}\in F^{n\times n}$ be invertible. The unique solution of
        $\mathbf{A}\mathbf{x}=\mathbf{b}$ is given by
        \[
            x_i
            =
            \frac{\det(\mathbf{A}_i)}{\det(\mathbf{A})},
        \]
        where $\mathbf{A}_i$ is obtained from $\mathbf{A}$ by replacing its
        $i$th column with $\mathbf{b}$.
    \end{theorem}

    \begin{formula}[Adjugate Formula for the Inverse]
        For every square matrix $\mathbf{A}$,
        \[
            \mathbf{A}\adj(\mathbf{A})
            =
            \adj(\mathbf{A})\mathbf{A}
            =
            \det(\mathbf{A})\mathbf{I}_n.
        \]
        Hence, if $\det(\mathbf{A})\neq 0$, then
        \[
            \mathbf{A}^{-1}
            =
            \frac{1}{\det(\mathbf{A})}\adj(\mathbf{A}).
        \]
    \end{formula}

    \begin{formula}[Determinant as Volume]
        If $\mathbf{v}_1,\dots,\mathbf{v}_n\in\R^n$, then the $n$-dimensional
        volume of the parallelepiped spanned by these vectors is
        \[
            \left|
            \det
            \begin{pmatrix}
                \mathbf{v}_1 & \mathbf{v}_2 & \cdots & \mathbf{v}_n
            \end{pmatrix}
            \right|.
        \]
        In $\R^3$, this is the absolute value of the scalar triple product
        \[
            |(\mathbf{a}\times\mathbf{b})\cdot\mathbf{c}|.
        \]
    \end{formula}

    \begin{formula}[Scalar and Vector Triple Products]
        For $\mathbf{a},\mathbf{b},\mathbf{c}\in\R^3$,
        \[
            \mathbf{a}\cdot(\mathbf{b}\times\mathbf{c})
            =
            \mathbf{b}\cdot(\mathbf{c}\times\mathbf{a})
            =
            \mathbf{c}\cdot(\mathbf{a}\times\mathbf{b})
            =
            \det(\mathbf{a},\mathbf{b},\mathbf{c}).
        \]
        The vector triple-product identities are
        \[
            \mathbf{a}\times(\mathbf{b}\times\mathbf{c})
            =
            (\mathbf{a}\cdot\mathbf{c})\mathbf{b}
            -
            (\mathbf{a}\cdot\mathbf{b})\mathbf{c},
        \]
        \[
            (\mathbf{a}\times\mathbf{b})\times\mathbf{c}
            =
            (\mathbf{a}\cdot\mathbf{c})\mathbf{b}
            -
            (\mathbf{b}\cdot\mathbf{c})\mathbf{a}.
        \]
    \end{formula}

    \begin{formula}[Cross-Product Determinant Identity]
        For $\mathbf{a},\mathbf{b},\mathbf{c}\in\R^3$,
        \[
            \det(
                \mathbf{a}\times\mathbf{b},
                \mathbf{b}\times\mathbf{c},
                \mathbf{c}\times\mathbf{a}
            )
            =
            \det(\mathbf{a},\mathbf{b},\mathbf{c})^2.
        \]
        Therefore, if $\mathbf{a},\mathbf{b},\mathbf{c}$ are linearly
        independent, then
        \[
            \mathbf{a}\times\mathbf{b},
            \qquad
            \mathbf{b}\times\mathbf{c},
            \qquad
            \mathbf{c}\times\mathbf{a}
        \]
        are also linearly independent.
    \end{formula}

    \begin{proposition}[Basic Properties of Determinants]
        For square matrices $\mathbf{A},\mathbf{B}\in F^{n\times n}$, we have
        \[
            \det(\mathbf{A}\mathbf{B})
            =
            \det(\mathbf{A})\det(\mathbf{B}),
        \]
        \[
            \det(\mathbf{A}^{\mathsf T})=\det(\mathbf{A}),
        \]
        and, if $\mathbf{A}$ is invertible,
        \[
            \det(\mathbf{A}^{-1})
            =
            \frac{1}{\det(\mathbf{A})}.
        \]
    \end{proposition}

    \begin{formula}[Determinant of a Triangular Matrix]
        If $\mathbf{A}$ is upper triangular or lower triangular with diagonal
        entries $a_{11},a_{22},\dots,a_{nn}$, then
        \[
            \det(\mathbf{A})
            =
            a_{11}a_{22}\cdots a_{nn}.
        \]
    \end{formula}

    \begin{formula}[Determinant of a Block Triangular Matrix]
        If the diagonal blocks are square, then
        \[
            \det
            \begin{pmatrix}
                \mathbf{A} & \mathbf{B} \\
                \mathbf{0} & \mathbf{D}
            \end{pmatrix}
            =
            \det(\mathbf{A})\det(\mathbf{D}),
        \]
        and the same formula holds for a lower block triangular matrix.
    \end{formula}

    \begin{proposition}[Multilinearity and Alternation of the Determinant]
        The determinant is linear in each row separately and also linear in each
        column separately. If two rows or two columns are interchanged, the sign
        of the determinant changes. In particular, if the rows or columns of
        $\mathbf{A}$ are linearly dependent, then
        \[
            \det(\mathbf{A})=0.
        \]
    \end{proposition}

    \begin{strategy}
        In short determinant problems, first look for dependent rows or
        columns, row or column sums, repeated patterns, triangular form, and
        easy row operations. Direct expansion should usually be the last resort,
        except for $2\times 2$ and $3\times 3$ matrices.
    \end{strategy}

    \begin{proposition}[Characterization of the Determinant]
        Let
        \[
            f\colon (F^n)^n\to F
        \]
        be alternating and multilinear. Then
        \[
            f(\mathbf{v}_1,\dots,\mathbf{v}_n)
            =
            f(\mathbf{e}_1,\dots,\mathbf{e}_n)
            \det
            \begin{pmatrix}
                \mathbf{v}_1 & \cdots & \mathbf{v}_n
            \end{pmatrix}.
        \]
        In particular, the determinant is the unique alternating multilinear
        function satisfying
        \[
            \det(\mathbf{I}_n)=1.
        \]
    \end{proposition}

    \begin{formula}[Laplace Expansion; Cofactor Expansion]
        Let $\mathbf{A}=(a_{ij})\in F^{n\times n}$, and let $C_{ij}$ be the
        cofactor of $a_{ij}$. Expanding along row $i$ gives
        \[
            \det(\mathbf{A})
            =
            \sum_{j=1}^{n}a_{ij}C_{ij}.
        \]
        Expanding along column $j$ gives
        \[
            \det(\mathbf{A})
            =
            \sum_{i=1}^{n}a_{ij}C_{ij}.
        \]
    \end{formula}

    \begin{formula}[Sarrus' Rule]
        For a $3\times 3$ matrix,
        \[
            \det
            \begin{pmatrix}
                a & b & c \\
                d & e & f \\
                g & h & i
            \end{pmatrix}
            =
            aei+bfg+cdh-ceg-bdi-afh.
        \]
    \end{formula}

    \begin{formula}[Determinant and Elementary Row Operations]
        The determinant changes under elementary row operations as follows.
        \begin{itemize}
            \item Swapping two rows multiplies the determinant by $-1$.
            \item Multiplying one row by $c$ multiplies the determinant by $c$.
            \item Adding a multiple of one row to another row does not change
                  the determinant.
        \end{itemize}
        The same rules hold for columns.
    \end{formula}


    \begin{algorithm}[Computing a Determinant]
        For a square matrix $\mathbf{A}$, choose the shortest structure-aware
        route rather than expanding automatically.
        \begin{bookitems}
            \item First test for an immediate value: a zero row or column, or
                  linearly dependent rows or columns, gives determinant $0$;
                  triangular and block-triangular matrices reduce to products of
                  diagonal determinants.
            \item Use row or column operations to create zeros and approach
                  triangular form. Track row swaps and row scalings exactly;
                  adding a multiple of one row to another leaves the determinant
                  unchanged.
            \item If one row or column is sparse, Laplace expansion along that
                  row or column can be faster than elimination. For a
                  $3\times3$ matrix, Sarrus' rule is available when it is
                  genuinely shorter.
            \item Before doing arithmetic, recognize standard matrices and
                  low-rank structure. Vandermonde, Cauchy, Pascal, block,
                  permutation, and rank-one perturbation formulas often turn a
                  long calculation into one line.
            \item Check the result against structural facts: a singular matrix
                  must have determinant $0$, and
                  $\det(\mathbf{A}\mathbf{B})=\det(\mathbf{A})\det(\mathbf{B})$.
        \end{bookitems}
    \end{algorithm}

    \begin{formula}[Matrix Determinant Lemma]
        If $\mathbf{A}\in F^{n\times n}$ is invertible and
        $\mathbf{u},\mathbf{v}\in F^n$, then
        \[
            \det(\mathbf{A}+\mathbf{u}\mathbf{v}^{\mathsf T})
            =
            \det(\mathbf{A})
            \left(1+\mathbf{v}^{\mathsf T}\mathbf{A}^{-1}\mathbf{u}\right).
        \]
        In particular,
        \[
            \det(\mathbf{I}_n+\mathbf{u}\mathbf{v}^{\mathsf T})
            =
            1+\mathbf{v}^{\mathsf T}\mathbf{u}.
        \]
    \end{formula}

    \begin{theorem}[Sylvester's Determinant Theorem]
        If $\mathbf{A}\in F^{m\times n}$ and
        $\mathbf{B}\in F^{n\times m}$, then
        \[
            \det(\mathbf{I}_m+\mathbf{A}\mathbf{B})
            =
            \det(\mathbf{I}_n+\mathbf{B}\mathbf{A}).
        \]
        This allows a determinant involving a large product to be replaced by a
        determinant of the smaller dimension.
    \end{theorem}

    \begin{formula}[Sherman--Morrison Formula]
        If $\mathbf{A}$ is invertible and
        $1+\mathbf{v}^{\mathsf T}\mathbf{A}^{-1}\mathbf{u}\neq0$, then
        \[
            (\mathbf{A}+\mathbf{u}\mathbf{v}^{\mathsf T})^{-1}
            =
            \mathbf{A}^{-1}
            -
            \frac{
                \mathbf{A}^{-1}\mathbf{u}\mathbf{v}^{\mathsf T}\mathbf{A}^{-1}
            }{
                1+\mathbf{v}^{\mathsf T}\mathbf{A}^{-1}\mathbf{u}
            }.
        \]
        It is the inverse counterpart of the Matrix Determinant Lemma.
    \end{formula}

    \begin{formula}[Schur Complement Determinant Formula]
        If $\mathbf{A}$ is invertible, then
        \[
            \det
            \begin{pmatrix}
                \mathbf{A} & \mathbf{B} \\
                \mathbf{C} & \mathbf{D}
            \end{pmatrix}
            =
            \det(\mathbf{A})
            \det(\mathbf{D}-\mathbf{C}\mathbf{A}^{-1}\mathbf{B}).
        \]
        The matrix $\mathbf{D}-\mathbf{C}\mathbf{A}^{-1}\mathbf{B}$ is the
        Schur complement of $\mathbf{A}$.
    \end{formula}

    \begin{formula}[Symmetric Block Determinant]
        If $\mathbf{A}$ and $\mathbf{B}$ are square matrices of the same size,
        then
        \[
            \det
            \begin{pmatrix}
                \mathbf{A} & \mathbf{B} \\
                \mathbf{B} & \mathbf{A}
            \end{pmatrix}
            =
            \det(\mathbf{A}+\mathbf{B})
            \det(\mathbf{A}-\mathbf{B}).
        \]
        This follows by changing to the invariant subspaces consisting of
        vectors of the form $(\mathbf{x},\mathbf{x})$ and
        $(\mathbf{x},-\mathbf{x})$.
    \end{formula}

    \begin{definition}[Eigenvalue, Eigenvector, Eigenspace, and Characteristic Polynomial]
        Let $\mathbf{A}\in F^{n\times n}$. A scalar $\lambda\in F$ is an
        eigenvalue of $\mathbf{A}$ if there exists a nonzero vector
        $\mathbf{v}\in F^n$ s.t.
        \[
            \mathbf{A}\mathbf{v}=\lambda\mathbf{v}.
        \]
        Such a nonzero vector $\mathbf{v}$ is an eigenvector associated with
        $\lambda$. The eigenspace is
        \[
            E_\lambda
            =
            \Null(\mathbf{A}-\lambda\mathbf{I}_n).
        \]
        The characteristic polynomial of $\mathbf{A}$ is
        \[
            p_{\mathbf{A}}(\lambda)
            =
            \det(\lambda\mathbf{I}_n-\mathbf{A}).
        \]
        Regarding the eigenvalues over $\C$, the spectral radius is
        \[
            \rho(\mathbf{A})
            =
            \max\{\lvert\lambda\rvert:\lambda\text{ is an eigenvalue of }\mathbf{A}\}.
        \]
    \end{definition}


    \begin{algorithm}[Computing Eigenvalues and Eigenvectors]
        Let $\mathbf{A}\in F^{n\times n}$.
        \begin{bookitems}
            \item Inspect the matrix first. For triangular or block-triangular
                  matrices, the eigenvalues come from the diagonal blocks;
                  symmetry, all-ones structure, permutation structure, and
                  low-rank formulas can reveal eigenvectors without computing a
                  full characteristic polynomial.
            \item Otherwise compute
                  \[
                      p_{\mathbf{A}}(\lambda)
                      =\det(\lambda\mathbf{I}_n-\mathbf{A})
                  \]
                  and solve $p_{\mathbf{A}}(\lambda)=0$. Over $\R$, remember
                  that nonreal roots may occur; over $\C$, the characteristic
                  polynomial splits completely.
            \item For each eigenvalue $\lambda$, row-reduce
                  $\mathbf{A}-\lambda\mathbf{I}_n$ and find a basis of
                  $\Null(\mathbf{A}-\lambda\mathbf{I}_n)$. Every nonzero vector
                  in that null space is an eigenvector for $\lambda$.
            \item Use
                  \[
                      \sum_i\lambda_i=\tr(\mathbf{A}),
                      \qquad
                      \prod_i\lambda_i=\det(\mathbf{A})
                  \]
                  with algebraic multiplicity as fast checks. Compare algebraic
                  and geometric multiplicities when deciding whether
                  $\mathbf{A}$ is diagonalizable.
        \end{bookitems}
    \end{algorithm}

    \begin{definition}[Algebraic and Geometric Multiplicity]
        Let $\lambda$ be an eigenvalue of $\mathbf{A}\in F^{n\times n}$.
        Its algebraic multiplicity $m_a(\lambda)$ is its multiplicity as a root
        of $p_{\mathbf{A}}(t)$, while its geometric multiplicity is
        \[
            m_g(\lambda)
            =
            \dim\Null(\mathbf{A}-\lambda\mathbf{I}_n).
        \]
        For every eigenvalue,
        \[
            1\leq m_g(\lambda)\leq m_a(\lambda).
        \]
    \end{definition}

    \begin{definition}[Diagonalizable Matrix]
        A matrix $\mathbf{A}\in F^{n\times n}$ is diagonalizable over $F$ if
        there exist an invertible matrix $\mathbf{P}\in F^{n\times n}$ and a
        diagonal matrix $\mathbf{D}$ s.t.
        \[
            \mathbf{A}=\mathbf{P}\mathbf{D}\mathbf{P}^{-1}.
        \]
        Equivalently, the columns of $\mathbf{P}$ can be chosen as a basis of
        eigenvectors of $\mathbf{A}$. Diagonalization is useful because powers
        and polynomial functions reduce to entrywise calculations on
        $\mathbf{D}$:
        \[
            \mathbf{A}^k=\mathbf{P}\mathbf{D}^k\mathbf{P}^{-1}.
        \]
    \end{definition}

    \begin{definition}[Conjugate Transpose; Hermitian Matrix]
        For a complex matrix $\mathbf{A}=(a_{ij})$, its conjugate transpose, or
        Hermitian transpose, is
        \[
            \mathbf{A}^{\mathsf H}
            =
            \overline{\mathbf{A}}^{\mathsf T}
            =
            (\overline{a_{ji}}).
        \]
        A complex square matrix is Hermitian if
        \[
            \mathbf{A}^{\mathsf H}=\mathbf{A}.
        \]
        For real matrices, the Hermitian transpose is the ordinary transpose.
    \end{definition}

    \begin{proposition}[Basic Rules for the Hermitian Transpose]
        For complex matrices of compatible sizes,
        \[
            (\mathbf{A}+\mathbf{B})^{\mathsf H}
            =
            \mathbf{A}^{\mathsf H}+\mathbf{B}^{\mathsf H},
            \qquad
            (c\mathbf{A})^{\mathsf H}
            =
            \overline{c}\,\mathbf{A}^{\mathsf H},
        \]
        \[
            (\mathbf{A}\mathbf{B})^{\mathsf H}
            =
            \mathbf{B}^{\mathsf H}\mathbf{A}^{\mathsf H},
            \qquad
            (\mathbf{A}^{\mathsf H})^{\mathsf H}
            =
            \mathbf{A}.
        \]
    \end{proposition}

    \begin{definition}[Orthogonal and Unitary Matrices]
        A real square matrix $\mathbf{Q}$ is orthogonal if
        \[
            \mathbf{Q}^{\mathsf T}\mathbf{Q}=\mathbf{I}_n.
        \]
        A complex square matrix $\mathbf{U}$ is unitary if
        \[
            \mathbf{U}^{\mathsf H}\mathbf{U}=\mathbf{I}_n.
        \]
        Equivalently, their columns form orthonormal bases over $\R$ and $\C$,
        respectively.
    \end{definition}

    \begin{formula}[Special Matrix Recognition Table]
        For real or complex square matrices, the following implications permit
        rapid determinant, inverse, and eigenvalue calculations.

        {\small
        \renewcommand{\arraystretch}{1.18}
        \begin{longtable}{@{}L{0.25\textwidth}L{0.28\textwidth}L{0.39\textwidth}@{}}
            \toprule
            Matrix type & Defining relation & Immediate consequences \\
            \midrule
            \endfirsthead

            \toprule
            Matrix type & Defining relation & Immediate consequences \\
            \midrule
            \endhead

            \bottomrule
            \endlastfoot

            Diagonal &
            $\mathbf{D}=\diag(d_1,\dots,d_n)$ &
            eigenvalues $d_i$; $\det(\mathbf{D})=\prod_i d_i$ \\
            Triangular &
            $a_{ij}=0$ for $i>j$ or for $i<j$ &
            eigenvalues are diagonal entries; determinant is their product \\
            Orthogonal &
            $\mathbf{Q}^{\mathsf T}\mathbf{Q}=\mathbf{I}_n$ &
            $\mathbf{Q}^{-1}=\mathbf{Q}^{\mathsf T}$; $|\det(\mathbf{Q})|=1$ \\
            Unitary &
            $\mathbf{U}^{\mathsf H}\mathbf{U}=\mathbf{I}_n$ &
            $\mathbf{U}^{-1}=\mathbf{U}^{\mathsf H}$; every eigenvalue has modulus $1$ \\
            Idempotent &
            $\mathbf{P}^2=\mathbf{P}$ &
            eigenvalues belong to $\{0,1\}$; $\rk(\mathbf{P})=\tr(\mathbf{P})$ \\
            Involutory &
            $\mathbf{A}^2=\mathbf{I}_n$ &
            $\mathbf{A}^{-1}=\mathbf{A}$; eigenvalues belong to $\{-1,1\}$ \\
            Nilpotent &
            $\mathbf{N}^k=\mathbf{0}$ for some $k$ &
            every eigenvalue is $0$; $\det(\mathbf{N})=0$ \\
            Permutation &
            one $1$ in each row and column &
            $\mathbf{P}^{-1}=\mathbf{P}^{\mathsf T}$; determinant is $\pm1$ \\
            Orthogonal projection &
            $\mathbf{P}^2=\mathbf{P}=\mathbf{P}^{\mathsf T}$ &
            $\R^n=\Col(\mathbf{P})\oplus\Null(\mathbf{P})$ \\
        \end{longtable}
        }
    \end{formula}

    \begin{formula}[Vandermonde Determinant]
        For scalars $x_1,x_2,\dots,x_n$,
        \[
            \det
            \begin{pmatrix}
                1 & x_1 & x_1^2 & \cdots & x_1^{n-1} \\
                1 & x_2 & x_2^2 & \cdots & x_2^{n-1} \\
                \vdots & \vdots & \vdots & \ddots & \vdots \\
                1 & x_n & x_n^2 & \cdots & x_n^{n-1}
            \end{pmatrix}
            =
            \prod_{1\leq i<j\leq n}(x_j-x_i).
        \]
    \end{formula}

    \begin{fact}[Pascal Matrices]
        Define the lower triangular Pascal matrix by
        \[
            \mathbf{P}_n=(p_{ij})_{0\leq i,j\leq n},
            \qquad
            p_{ij}
            =
            \begin{cases}
                \binom{i}{j}, & j\leq i, \\
                0, & j>i.
            \end{cases}
        \]
        Since all diagonal entries are equal to $1$, we have
        \[
            \det(\mathbf{P}_n)=1.
        \]
        Define the symmetric Pascal matrix by
        \[
            \mathbf{S}_n=(s_{ij})_{0\leq i,j\leq n},
            \qquad
            s_{ij}=\binom{i+j}{i}.
        \]
        This matrix also has determinant $1$.
    \end{fact}

    \begin{remark}
        Pascal matrices have combinatorial entries while their determinants admit
        simple closed forms.
    \end{remark}

    \begin{formula}[Cauchy Determinant]
        Let
        \[
            c_{ij}=\frac{1}{x_i+y_j},
        \]
        where every denominator is nonzero and the $x_i$ and $y_j$ are pairwise
        distinct within their respective families. Then
        \[
            \det(c_{ij})_{1\leq i,j\leq n}
            =
            \frac{
                \displaystyle
                \prod_{1\leq i<j\leq n}(x_j-x_i)(y_j-y_i)
            }{
                \displaystyle
                \prod_{i=1}^{n}\prod_{j=1}^{n}(x_i+y_j)
            }.
        \]
        The Hilbert matrix is a famous special case of a Cauchy matrix.
    \end{formula}

    \begin{formula}[Hilbert Matrix Determinant]
        The $n\times n$ Hilbert matrix is
        \[
            \mathbf{H}_n
            =
            \left(\frac{1}{i+j-1}\right)_{1\leq i,j\leq n}.
        \]
        As a special Cauchy matrix, it satisfies
        \[
            \det(\mathbf{H}_n)
            =
            \frac{\displaystyle\prod_{k=1}^{n-1}(k!)^4}
                 {\displaystyle\prod_{k=1}^{2n-1}k!}.
        \]
    \end{formula}

    \begin{definition}[Hadamard Matrix]
        A Hadamard matrix of order $n$ is an $n\times n$ matrix
        $\mathbf{H}$ whose entries are all $+1$ or $-1$ and whose rows are
        mutually orthogonal. Equivalently,
        \[
            \mathbf{H}\mathbf{H}^{\mathsf T}=n\mathbf{I}_n.
        \]
    \end{definition}

    \begin{formula}[Hadamard Determinant]
        If $\mathbf{H}$ is a Hadamard matrix of order $n$, then
        \[
            |\det(\mathbf{H})|
            =
            n^{n/2}.
        \]
    \end{formula}

    \begin{formula}[Permutation Matrices]
        A permutation matrix $\mathbf{P}_\sigma$ has exactly one entry equal
        to $1$ in each row and each column. It satisfies
        \[
            \mathbf{P}_\sigma^{-1}
            =
            \mathbf{P}_\sigma^{\mathsf T},
            \qquad
            \det(\mathbf{P}_\sigma)=\sgn(\sigma).
        \]
        Its powers are determined by the cycle structure of $\sigma$.
    \end{formula}

    \begin{formula}[Eigenvalues of a Circulant Matrix]
        Let $\mathbf{C}$ be the circulant matrix whose first row is
        \[
            (c_0,c_1,\dots,c_{n-1}).
        \]
        With $\omega_n=e^{-2\pi i/n}$, its eigenvalues are
        \[
            \lambda_j
            =
            \sum_{k=0}^{n-1}c_k\omega_n^{jk},
            \qquad
            j=0,1,\dots,n-1.
        \]
        Hence $\det(\mathbf{C})=\prod_{j=0}^{n-1}\lambda_j$.
    \end{formula}

    \begin{formula}[Discrete Fourier Matrix]
        Let $\omega_n=e^{2\pi i/n}$. The normalized discrete Fourier matrix is
        \[
            \mathbf{F}_n
            =
            \frac{1}{\sqrt{n}}
            \left(\omega_n^{-jk}\right)_{0\leq j,k\leq n-1}.
        \]
        It is unitary:
        \[
            \mathbf{F}_n^{\mathsf H}\mathbf{F}_n=\mathbf{I}_n.
        \]
        Every circulant matrix is diagonalized by $\mathbf{F}_n$.
    \end{formula}

    \begin{formula}[Symmetric Tridiagonal Toeplitz Matrix]
        For
        \[
            \mathbf{T}_n
            =
            \begin{pmatrix}
                a & b & 0 & \cdots & 0 \\
                b & a & b & \ddots & \vdots \\
                0 & b & a & \ddots & 0 \\
                \vdots & \ddots & \ddots & \ddots & b \\
                0 & \cdots & 0 & b & a
            \end{pmatrix},
        \]
        the eigenvalues are
        \[
            a+2b\cos\frac{k\pi}{n+1},
            \qquad
            k=1,2,\dots,n.
        \]
    \end{formula}

    \begin{formula}[Continuant Recurrence for Tridiagonal Determinants]
        Let
        \[
            \mathbf{T}_n
            =
            \begin{pmatrix}
                a_1 & b_1 & 0 & \cdots & 0 \\
                c_1 & a_2 & b_2 & \ddots & \vdots \\
                0 & c_2 & a_3 & \ddots & 0 \\
                \vdots & \ddots & \ddots & \ddots & b_{n-1} \\
                0 & \cdots & 0 & c_{n-1} & a_n
            \end{pmatrix}.
        \]
        If $D_n=\det(\mathbf{T}_n)$, then
        \[
            D_0=1,
            \qquad
            D_1=a_1,
            \qquad
            D_n=a_nD_{n-1}-b_{n-1}c_{n-1}D_{n-2}.
        \]
    \end{formula}

    \begin{formula}[Householder Reflection]
        For a nonzero vector $\mathbf{v}$,
        \[
            \mathbf{H}
            =
            \mathbf{I}_n
            -2\frac{\mathbf{v}\mathbf{v}^{\mathsf T}}
                     {\mathbf{v}^{\mathsf T}\mathbf{v}}
        \]
        is symmetric, orthogonal, and involutory:
        \[
            \mathbf{H}^{\mathsf T}=\mathbf{H},
            \qquad
            \mathbf{H}^{\mathsf T}\mathbf{H}=\mathbf{I}_n,
            \qquad
            \mathbf{H}^2=\mathbf{I}_n.
        \]
        It reflects vectors across the hyperplane perpendicular to
        $\mathbf{v}$.
    \end{formula}

    \begin{formula}[Rotation and Reflection Matrices]
        Rotation by angle $\theta$ in the plane is represented by
        \[
            \begin{pmatrix}
                \cos\theta & -\sin\theta \\
                \sin\theta & \cos\theta
            \end{pmatrix}.
        \]
        Reflection across the line making angle $\theta$ with the positive
        $x$-axis is represented by
        \[
            \begin{pmatrix}
                \cos 2\theta & \sin 2\theta \\
                \sin 2\theta & -\cos 2\theta
            \end{pmatrix}.
        \]
    \end{formula}

    \begin{formula}[Eigenvalues of a Three-Dimensional Rotation Matrix]
        A rotation of $\R^3$ by angle $\theta$ around an axis through the
        origin has eigenvalues
        \[
            1,
            \qquad
            e^{i\theta},
            \qquad
            e^{-i\theta}.
        \]
        The eigenvalue $1$ corresponds to the rotation axis. Therefore, if
        $\theta\not\equiv 0,\pi\pmod{2\pi}$, the rotation matrix has exactly one
        real eigenvalue.
    \end{formula}

    \begin{proposition}[Trace, Determinant, and Eigenvalues]
        Let $\mathbf{A}\in\C^{n\times n}$ have eigenvalues
        $\lambda_1,\dots,\lambda_n$, counted with algebraic multiplicity. Then
        \[
            \tr(\mathbf{A})
            =
            \lambda_1+\cdots+\lambda_n
        \]
        and
        \[
            \det(\mathbf{A})
            =
            \lambda_1\lambda_2\cdots\lambda_n.
        \]
    \end{proposition}

    \begin{proposition}[Cyclic Property of the Trace]
        Whenever the products are defined and square,
        \[
            \tr(\mathbf{A}\mathbf{B})
            =
            \tr(\mathbf{B}\mathbf{A}).
        \]
        More generally,
        \[
            \tr(\mathbf{A}\mathbf{B}\mathbf{C})
            =
            \tr(\mathbf{B}\mathbf{C}\mathbf{A})
            =
            \tr(\mathbf{C}\mathbf{A}\mathbf{B}).
        \]
        The factors may be cyclically permuted, but they may not in general be
        rearranged arbitrarily.
    \end{proposition}

    \begin{proposition}[Eigenvalues and the Characteristic Polynomial]
        A scalar $\lambda$ is an eigenvalue of $\mathbf{A}$ if and only if
        \[
            \det(\lambda\mathbf{I}_n-\mathbf{A})=0.
        \]
        Equivalently, the eigenvalues of $\mathbf{A}$ are the roots of the
        characteristic polynomial $p_{\mathbf{A}}(\lambda)$.
    \end{proposition}

    \begin{formula}[All-Ones Matrix]
        Let $\mathbf{J}_n$ be the $n\times n$ matrix whose entries are all $1$.
        Then
        \[
            \mathbf{J}_n\ones=n\ones,
            \qquad
            \mathbf{J}_n\mathbf{v}=\mathbf{0}
            \quad\text{if}\quad
            v_1+\cdots+v_n=0.
        \]
        Hence the eigenvalues of $\mathbf{J}_n$ are $n$ with multiplicity $1$
        and $0$ with multiplicity $n-1$. The matrix $\mathbf{J}_n-\mathbf{I}_n$
        has eigenvalues $n-1$ with multiplicity $1$ and $-1$ with multiplicity
        $n-1$.
    \end{formula}

    \begin{formula}[Constant Diagonal and Off-Diagonal Matrix]
        Let $\mathbf{A}\in F^{n\times n}$ have diagonal entries $a$ and
        off-diagonal entries $b$. Then
        \[
            \mathbf{A}=(a-b)\mathbf{I}_n+b\mathbf{J}_n.
        \]
        Counting algebraic multiplicity, its eigenvalues consist of
        \[
            a+(n-1)b
        \]
        once and
        \[
            a-b
        \]
        $n-1$ times. If these two values coincide, they together give the single
        eigenvalue of multiplicity $n$.
    \end{formula}

    \begin{formula}[Rank-One Outer Product]
        Let $\mathbf{v},\mathbf{w}\in\R^n$ be nonzero and set
        \[
            \mathbf{A}=\mathbf{v}\mathbf{w}^{\mathsf T}.
        \]
        Then $\rk(\mathbf{A})=1$ and
        \[
            \mathbf{A}\mathbf{v}
            =
            (\mathbf{w}^{\mathsf T}\mathbf{v})\mathbf{v}.
        \]
        The characteristic polynomial is
        \[
            \lambda^{n-1}
            \bigl(\lambda-\mathbf{w}^{\mathsf T}\mathbf{v}\bigr).
        \]
        Thus the eigenvalues, counted with algebraic multiplicity, are
        $\mathbf{w}^{\mathsf T}\mathbf{v}$ and $0$ repeated $n-1$ times. If
        $\mathbf{w}^{\mathsf T}\mathbf{v}=0$, every eigenvalue is $0$ and
        $\mathbf{A}^2=\mathbf{0}$.
    \end{formula}

    \begin{theorem}[Cayley--Hamilton Theorem]
        Every square matrix satisfies its own characteristic polynomial. That
        is, if
        \[
            p_{\mathbf{A}}(\lambda)
            =
            \det(\lambda\mathbf{I}_n-\mathbf{A}),
        \]
        then
        \[
            p_{\mathbf{A}}(\mathbf{A})=\mathbf{0}.
        \]
    \end{theorem}

    \begin{remark}
        Cayley--Hamilton is a famous named theorem. It is especially useful for
        reducing high powers of a matrix to lower powers.
    \end{remark}

    \begin{definition}[Minimal Polynomial]
        The minimal polynomial $m_{\mathbf{A}}(t)$ of a square matrix
        $\mathbf{A}$ is the unique monic polynomial of smallest positive degree
        s.t.
        \[
            m_{\mathbf{A}}(\mathbf{A})=\mathbf{0}.
        \]
        By the Cayley--Hamilton Theorem, it divides the characteristic
        polynomial $p_{\mathbf{A}}(t)$. A short minimal polynomial is useful
        because it reduces high powers of $\mathbf{A}$ to lower powers.
    \end{definition}

    \begin{formula}[Companion Matrix and Linear Recurrences]
        For
        \[
            p(x)=x^m-c_1x^{m-1}-\cdots-c_m,
        \]
        the companion matrix
        \[
            \mathbf{C}
            =
            \begin{pmatrix}
                c_1 & c_2 & \cdots & c_{m-1} & c_m \\
                1 & 0 & \cdots & 0 & 0 \\
                0 & 1 & \ddots & 0 & 0 \\
                \vdots & \ddots & \ddots & \vdots & \vdots \\
                0 & 0 & \cdots & 1 & 0
            \end{pmatrix}
        \]
        has characteristic polynomial $p$. The recurrence
        \[
            a_n=c_1a_{n-1}+\cdots+c_ma_{n-m}
        \]
        is encoded by multiplication by $\mathbf{C}$, so fast matrix
        exponentiation computes distant terms efficiently.
    \end{formula}

    \begin{formula}[Fibonacci $Q$-Matrix]
        For
        \[
            \mathbf{Q}
            =
            \begin{pmatrix}
                1 & 1 \\
                1 & 0
            \end{pmatrix},
        \]
        we have
        \[
            \mathbf{Q}^n
            =
            \begin{pmatrix}
                F_{n+1} & F_n \\
                F_n & F_{n-1}
            \end{pmatrix}
            \qquad(n\geq1).
        \]
        Taking determinants gives Cassini's identity
        \[
            F_{n+1}F_{n-1}-F_n^2=(-1)^n.
        \]
    \end{formula}

    \begin{formula}[Matrix Geometric Series]
        Let $\mathbf{A}$ be an $n\times n$ matrix. If
        $\mathbf{I}_n-\mathbf{A}$ is invertible, then
        \[
            \sum_{k=0}^{N-1}\mathbf{A}^k
            =
            (\mathbf{I}_n-\mathbf{A}^N)(\mathbf{I}_n-\mathbf{A})^{-1}.
        \]
        If the spectral radius of $\mathbf{A}$ is less than $1$, then
        \[
            \sum_{k=0}^{\infty}\mathbf{A}^k
            =
            (\mathbf{I}_n-\mathbf{A})^{-1}.
        \]
    \end{formula}

    \begin{strategy}[Fast Matrix Powers]
        To compute a large power $\mathbf{A}^N$, first look for one of the
        following: diagonalization, a short minimal polynomial, idempotence,
        involution, nilpotence, rank one, a circulant structure, or a decomposition
        into $a\mathbf{I}_n+b\mathbf{J}_n$. Cayley--Hamilton then reduces every
        sufficiently high power to a linear combination of lower powers.
    \end{strategy}



    \begin{theorem}[Diagonalization Criteria]
        Suppose the characteristic polynomial of
        $\mathbf{A}\in F^{n\times n}$ splits over $F$. The following are
        equivalent:
        \[
            \mathbf{A}\text{ is diagonalizable over }F,
        \]
        \[
            \sum_{\lambda}m_g(\lambda)=n,
        \]
        and
        \[
            m_g(\lambda)=m_a(\lambda)
            \quad\text{for every eigenvalue }\lambda.
        \]
        Equivalently, the minimal polynomial of $\mathbf{A}$ splits over $F$
        into distinct linear factors.
    \end{theorem}

    \begin{corollary}[Distinct Eigenvalues Imply Diagonalizability]
        If an $n\times n$ matrix $\mathbf{A}$ has $n$ distinct eigenvalues in
        $F$, then $\mathbf{A}$ is diagonalizable over $F$.
    \end{corollary}

    \begin{remark}
        The field matters. A real matrix can fail to be diagonalizable over
        $\R$ because its characteristic polynomial has nonreal roots, while the
        same matrix may be diagonalizable over $\C$. Over $\C$, every
        characteristic polynomial splits, so only the multiplicity conditions
        remain. Stronger diagonalization results for symmetric, Hermitian, and
        normal matrices appear below.
    \end{remark}

    \begin{proposition}[Polynomials of a Diagonalizable Matrix]
        Let $\mathbf{A}\in F^{n\times n}$ be diagonalizable over $F$, and suppose
        \[
            \mathbf{A}
            =
            \mathbf{S}\mathbf{D}\mathbf{S}^{-1},
            \qquad
            \mathbf{D}
            =
            \diag(\lambda_1,\dots,\lambda_n).
        \]
        For a polynomial $p(x)\in F[x]$, we have
        \[
            p(\mathbf{A})
            =
            \mathbf{S}p(\mathbf{D})\mathbf{S}^{-1},
        \]
        where
        \[
            p(\mathbf{D})
            =
            \diag\bigl(p(\lambda_1),\dots,p(\lambda_n)\bigr).
        \]
        In particular,
        \[
            \det(p(\mathbf{A}))
            =
            p(\lambda_1)p(\lambda_2)\cdots p(\lambda_n).
        \]
    \end{proposition}

    \begin{idea}
        Similar matrices represent the same linear transformation in different
        bases. Therefore, applying a polynomial to a diagonalizable matrix is
        the same as applying that polynomial to each diagonal entry after
        diagonalization.
    \end{idea}

    \begin{strategy}
        Distinct eigenvalues are an easy sufficient condition for
        diagonalizability. They are not necessary: a matrix can be
        diagonalizable even when some eigenvalues are repeated.
    \end{strategy}

    \begin{theorem}[Gershgorin Circle Theorem]
        Let $\mathbf{A}=(a_{ij})\in\C^{n\times n}$. Every eigenvalue of
        $\mathbf{A}$ lies in at least one disk
        \[
            \left\{z\in\C\mid
            |z-a_{ii}|\leq\sum_{j\neq i}|a_{ij}|
            \right\}.
        \]
    \end{theorem}

    \begin{theorem}[Perron--Frobenius Theorem, Positive-Matrix Form]
        If every entry of a real square matrix $\mathbf{A}$ is positive, then
        $\rho(\mathbf{A})$ is a positive simple eigenvalue with a positive
        eigenvector, and every other eigenvalue $\lambda$ satisfies
        \[
            |\lambda|<\rho(\mathbf{A}).
        \]
    \end{theorem}

    \begin{theorem}[Spectral Theorem for Real Symmetric Matrices]
        If $\mathbf{A}\in\R^{n\times n}$ is symmetric, then all eigenvalues of
        $\mathbf{A}$ are real, and there exists an orthogonal matrix
        $\mathbf{Q}$ and a diagonal matrix $\mathbf{D}$ s.t.
        \[
            \mathbf{A}
            =
            \mathbf{Q}\mathbf{D}\mathbf{Q}^{\mathsf T}.
        \]
        Equivalently, $\R^n$ has an orthonormal basis consisting of eigenvectors
        of $\mathbf{A}$.
    \end{theorem}

    \begin{remark}
        A real symmetric matrix has real eigenvalues and admits orthogonal
        diagonalization.
    \end{remark}

    \begin{definition}[Normal Matrix]
        A complex square matrix $\mathbf{A}$ is normal if
        \[
            \mathbf{A}^{\mathsf H}\mathbf{A}
            =
            \mathbf{A}\mathbf{A}^{\mathsf H}.
        \]
        Hermitian, skew-Hermitian, and unitary matrices are normal.
    \end{definition}

    \begin{theorem}[Spectral Theorem for Normal Matrices]
        A complex square matrix is normal if and only if it is unitarily
        diagonalizable. Equivalently, $\mathbf{A}$ is normal if and only if
        there exist a unitary matrix $\mathbf{U}$ and a diagonal matrix
        $\mathbf{D}$ s.t.
        \[
            \mathbf{A}
            =
            \mathbf{U}\mathbf{D}\mathbf{U}^{\mathsf H}.
        \]
    \end{theorem}

    \begin{proposition}[Hermitian Matrix Quick Facts]
        If $\mathbf{A}$ is Hermitian, then
        \[
            \mathbf{x}^{\mathsf H}\mathbf{A}\mathbf{x}\in\R
        \]
        for every $\mathbf{x}\in\C^n$. All eigenvalues of $\mathbf{A}$ are real,
        and eigenvectors corresponding to distinct eigenvalues are orthogonal.
    \end{proposition}

    \begin{theorem}[Spectral Theorem for Hermitian Matrices]
        If $\mathbf{A}\in\C^{n\times n}$ is Hermitian, then there exists a
        unitary matrix $\mathbf{U}$ and a real diagonal matrix $\mathbf{D}$ s.t.
        \[
            \mathbf{A}
            =
            \mathbf{U}\mathbf{D}\mathbf{U}^{\mathsf H}.
        \]
        Equivalently, $\C^n$ has an orthonormal basis consisting of eigenvectors
        of $\mathbf{A}$.
    \end{theorem}

    \begin{remark}
        The real analogue of a Hermitian matrix is a real symmetric matrix.
        The real analogue of a unitary matrix is an orthogonal matrix.
    \end{remark}

    \begin{definition}[Markov Matrix; Stochastic Matrix]
        A column-stochastic matrix is a square matrix with nonnegative entries
        whose columns each sum to $1$. Thus
        \[
            \ones^{\mathsf T}\mathbf{A}=\ones^{\mathsf T}.
        \]
        Consequently, $1$ is an eigenvalue of $\mathbf{A}^{\mathsf T}$ and
        therefore also of $\mathbf{A}$. A probability column vector evolves by
        \[
            \mathbf{p}_{k+1}=\mathbf{A}\mathbf{p}_k.
        \]
        A stationary distribution is a probability vector satisfying
        $\mathbf{A}\mathbf{p}=\mathbf{p}$.
    \end{definition}

    \begin{remark}
        Markov matrices are linear algebra objects behind finite-state Markov
        chains. A stationary distribution is an eigenvector corresponding to the
        eigenvalue $1$.
    \end{remark}

    \begin{theorem}[Cauchy--Schwarz Inequality]
        For vectors $\mathbf{u},\mathbf{v}\in\R^n$, we have
        \[
            |\langle\mathbf{u},\mathbf{v}\rangle|
            \leq
            \|\mathbf{u}\|\,\|\mathbf{v}\|.
        \]
        Equivalently, for real numbers $a_1,\dots,a_n$ and
        $b_1,\dots,b_n$,
        \[
            \left(\sum_{i=1}^{n}a_i b_i\right)^2
            \leq
            \left(\sum_{i=1}^{n}a_i^2\right)
            \left(\sum_{i=1}^{n}b_i^2\right).
        \]
    \end{theorem}

    \begin{strategy}
        The expression $\sum a_i b_i$ is a strong signal for Cauchy--Schwarz.
        The vector form is often the most efficient method to recognize the inequality.
    \end{strategy}

    \begin{corollary}[Triangle Inequality]
        For vectors $\mathbf{u},\mathbf{v}\in\R^n$, we have
        \[
            \|\mathbf{u}+\mathbf{v}\|
            \leq
            \|\mathbf{u}\|+\|\mathbf{v}\|.
        \]
    \end{corollary}

    \begin{proposition}[Properties of Orthogonal Matrices]
        If $\mathbf{Q}$ is orthogonal, then
        \[
            \|\mathbf{Q}\mathbf{v}\|=\|\mathbf{v}\|,
            \qquad
            \langle\mathbf{Q}\mathbf{u},\mathbf{Q}\mathbf{v}\rangle
            =
            \langle\mathbf{u},\mathbf{v}\rangle,
            \qquad
            \det(\mathbf{Q})=\pm 1.
        \]
    \end{proposition}

    \begin{definition}[Positive Definite Matrix]
        A Hermitian matrix $\mathbf{A}\in\C^{n\times n}$ is positive definite if
        \[
            \mathbf{x}^{\mathsf H}\mathbf{A}\mathbf{x}>0
        \]
        for every nonzero $\mathbf{x}\in\C^n$. Over $\R$, this is the familiar
        condition that $\mathbf{A}$ is symmetric and
        \[
            \mathbf{x}^{\mathsf T}\mathbf{A}\mathbf{x}>0
        \]
        for every nonzero $\mathbf{x}\in\R^n$.
    \end{definition}

    \begin{definition}[Singular Values]
        Let $\mathbf{A}\in\C^{m\times n}$. The singular values of
        $\mathbf{A}$ are the nonnegative square roots of the eigenvalues of
        \[
            \mathbf{A}^{\mathsf H}\mathbf{A}.
        \]
        For real matrices, $\mathbf{A}^{\mathsf H}=\mathbf{A}^{\mathsf T}$.
    \end{definition}

    \begin{formula}[Matrix Decomposition Reference Table]
        The most useful matrix decompositions are summarized by their hypotheses
        and the computation they simplify.

        {\small
        \renewcommand{\arraystretch}{1.16}
        \begin{longtable}{@{}L{0.20\textwidth}L{0.26\textwidth}L{0.25\textwidth}L{0.21\textwidth}@{}}
            \toprule
            Decomposition & Form & Typical condition & Main use \\
            \midrule
            \endfirsthead

            \toprule
            Decomposition & Form & Typical condition & Main use \\
            \midrule
            \endhead

            \bottomrule
            \endlastfoot

            Rank factorization & $\mathbf{A}=\mathbf{U}\mathbf{V}$ &
                $\rk(\mathbf{A})=r$, with $\mathbf{U}\in F^{m\times r}$,
                $\mathbf{V}\in F^{r\times n}$ of rank $r$ & expose rank and
                column/row spaces \\
            LU & $\mathbf{A}=\mathbf{L}\mathbf{U}$ & elimination without
                row exchanges; with pivoting, $\mathbf{P}\mathbf{A}=\mathbf{L}\mathbf{U}$ &
                repeated linear solves, determinants \\
            QR & $\mathbf{A}=\mathbf{Q}\mathbf{R}$ & full column rank, with
                orthonormal columns of $\mathbf{Q}$ & least squares,
                orthogonalization \\
            Cholesky & $\mathbf{A}=\mathbf{L}\mathbf{L}^{\mathsf H}$ &
                Hermitian positive definite; over $\R$, this is
                $\mathbf{L}\mathbf{L}^{\mathsf T}$ & fast stable linear solves \\
            Diagonalization & $\mathbf{A}=\mathbf{P}\mathbf{D}\mathbf{P}^{-1}$ &
                a basis of eigenvectors & powers and functions of $\mathbf{A}$ \\
            Orthogonal spectral & $\mathbf{A}=\mathbf{Q}\mathbf{D}\mathbf{Q}^{\mathsf T}$ &
                real symmetric & real orthonormal eigenbasis \\
            Unitary spectral & $\mathbf{A}=\boldsymbol{\Omega}\boldsymbol{\Lambda}
                \boldsymbol{\Omega}^{\mathsf H}$ & complex normal & orthonormal
                eigenbasis over $\C$ \\
            Schur & $\mathbf{A}=\mathbf{Q}\mathbf{T}\mathbf{Q}^{\mathsf H}$ &
                every complex square matrix & triangularize without assuming
                diagonalizability \\
            SVD & $\mathbf{A}=\mathbf{U}\boldsymbol{\Sigma}\mathbf{V}^{\mathsf H}$ &
                every real or complex matrix & rank, least squares,
                conditioning, best low-rank approximation \\
        \end{longtable}
        }
        For real matrices, $\mathbf{V}^{\mathsf H}=\mathbf{V}^{\mathsf T}$.
        The displayed letters are matrices and therefore follow the book's
        bold-uppercase convention.
    \end{formula}

    \begin{theorem}[LU Decomposition]
        When Gaussian elimination can be performed without row exchanges, a
        square matrix $\mathbf{A}$ can be written as
        \[
            \mathbf{A}=\mathbf{L}\mathbf{U},
        \]
        where $\mathbf{L}$ is lower triangular and $\mathbf{U}$ is upper
        triangular.
    \end{theorem}

    \begin{theorem}[Gram--Schmidt Process; QR Decomposition]
        From any linearly independent list
        \[
            \mathbf{v}_1,\dots,\mathbf{v}_k
        \]
        in an inner product space, one can construct an orthonormal list
        \[
            \mathbf{q}_1,\dots,\mathbf{q}_k
        \]
        with the same span. In matrix form, this leads to the QR decomposition
        \[
            \mathbf{A}=\mathbf{Q}\mathbf{R},
        \]
        where the columns of $\mathbf{Q}$ are orthonormal and $\mathbf{R}$ is
        upper triangular.
    \end{theorem}

    \begin{remark}
        Gram--Schmidt is the standard procedure for turning an arbitrary basis
        into an orthonormal basis.
    \end{remark}

    \begin{formula}[Orthogonal Projection Matrices]
        For a nonzero vector $\mathbf{v}\in\R^n$,
        \[
            \proj_{\mathbf{v}}\mathbf{x}
            =
            \frac{\mathbf{v}^{\mathsf T}\mathbf{x}}
                 {\mathbf{v}^{\mathsf T}\mathbf{v}}\mathbf{v},
        \]
        and the projection matrix onto $\Span(\mathbf{v})$ is
        \[
            \mathbf{P}_{\mathbf{v}}
            =
            \frac{\mathbf{v}\mathbf{v}^{\mathsf T}}
                 {\mathbf{v}^{\mathsf T}\mathbf{v}}.
        \]
        If $\|\mathbf{v}\|=1$, then
        \[
            \mathbf{P}_{\mathbf{v}}
            =
            \mathbf{v}\mathbf{v}^{\mathsf T},
            \qquad
            \mathbf{I}_n-\mathbf{v}\mathbf{v}^{\mathsf T}
        \]
        projects onto the hyperplane $\mathbf{v}^{\perp}$.

        If the columns of $\mathbf{Q}\in\R^{n\times k}$ are orthonormal, then
        the orthogonal projection matrix onto $\Col(\mathbf{Q})$ is
        \[
            \mathbf{P}
            =
            \mathbf{Q}\mathbf{Q}^{\mathsf T}.
        \]
        More generally, if the columns of $\mathbf{A}\in\R^{n\times k}$ are
        linearly independent, then
        \[
            \mathbf{P}
            =
            \mathbf{A}(\mathbf{A}^{\mathsf T}\mathbf{A})^{-1}
            \mathbf{A}^{\mathsf T}.
        \]
        Over $\C$, replace every transpose by the Hermitian transpose.
    \end{formula}


    \begin{theorem}[Least-Squares Normal Equations]
        Let $\mathbf{A}\in\R^{m\times n}$ and
        $\mathbf{b}\in\R^m$. A vector $\widehat{\mathbf{x}}$ minimizes
        \[
            \|\mathbf{A}\mathbf{x}-\mathbf{b}\|^2
        \]
        if and only if the residual is orthogonal to
        $\Col(\mathbf{A})$, equivalently,
        \[
            \mathbf{A}^{\mathsf T}\mathbf{A}\widehat{\mathbf{x}}
            =
            \mathbf{A}^{\mathsf T}\mathbf{b}.
        \]
        If the columns of $\mathbf{A}$ are linearly independent, then the
        minimizer is unique and
        \[
            \widehat{\mathbf{x}}
            =
            (\mathbf{A}^{\mathsf T}\mathbf{A})^{-1}
            \mathbf{A}^{\mathsf T}\mathbf{b}.
        \]
        The fitted vector is the orthogonal projection
        \[
            \widehat{\mathbf{b}}
            =
            \mathbf{A}\widehat{\mathbf{x}}
            =
            \mathbf{P}\mathbf{b},
            \qquad
            \mathbf{P}
            =
            \mathbf{A}(\mathbf{A}^{\mathsf T}\mathbf{A})^{-1}
            \mathbf{A}^{\mathsf T}.
        \]
    \end{theorem}

    \begin{proposition}[Projection Matrix Criterion]
        A matrix $\mathbf{P}$ is a projection matrix if
        \[
            \mathbf{P}^2=\mathbf{P}.
        \]
        Over $\R$, it is an orthogonal projection matrix if, in addition,
        \[
            \mathbf{P}^{\mathsf T}=\mathbf{P}.
        \]
        Over $\C$, the corresponding condition is
        \[
            \mathbf{P}^{\mathsf H}=\mathbf{P}.
        \]
    \end{proposition}

    \begin{theorem}[Sylvester's Criterion]
        A real symmetric matrix is positive definite if and only if all of its
        leading principal minors are positive.
    \end{theorem}

    \begin{theorem}[Singular Value Decomposition; SVD]
        Every real or complex matrix $\mathbf{A}\in\C^{m\times n}$ can be
        written as
        \[
            \mathbf{A}
            =
            \mathbf{U}\boldsymbol{\Sigma}\mathbf{V}^{\mathsf H},
        \]
        where $\mathbf{U}$ and $\mathbf{V}$ are unitary matrices and
        $\boldsymbol{\Sigma}$ is a rectangular diagonal matrix whose diagonal
        entries are the singular values of $\mathbf{A}$. For real
        $\mathbf{A}$, the factors may be chosen real and orthogonal, so
        $\mathbf{V}^{\mathsf H}=\mathbf{V}^{\mathsf T}$.
    \end{theorem}

    \begin{remark}
        Singular Value Decomposition applies to every real or complex matrix,
        not only to square or diagonalizable matrices. This is why SVD is one
        of the most useful decompositions in applied linear algebra.
    \end{remark}


    \begin{definition}[Generalized Eigenvectors and Jordan Chains]
        Let $\lambda$ be an eigenvalue of a square matrix $\mathbf{A}$. A
        nonzero vector $\mathbf{v}$ is a generalized eigenvector associated
        with $\lambda$ if
        \[
            (\mathbf{A}-\lambda\mathbf{I}_n)^k\mathbf{v}
            =
            \mathbf{0}
        \]
        for some positive integer $k$. A Jordan chain
        $\mathbf{v}_1,\dots,\mathbf{v}_r$ satisfies
        \[
            (\mathbf{A}-\lambda\mathbf{I}_n)\mathbf{v}_1=\mathbf{0},
            \qquad
            (\mathbf{A}-\lambda\mathbf{I}_n)\mathbf{v}_{j+1}
            =
            \mathbf{v}_j.
        \]
        Ordinary eigenvectors are precisely the generalized eigenvectors of
        rank $1$.
    \end{definition}

    \begin{theorem}[Jordan Canonical Form]
        Over the complex numbers, every square matrix is similar to a Jordan
        matrix. That is, for every $\mathbf{A}\in\C^{n\times n}$, there exists
        an invertible matrix $\mathbf{P}$ s.t.
        \[
            \mathbf{P}^{-1}\mathbf{A}\mathbf{P}
        \]
        is in Jordan canonical form.
    \end{theorem}

    \begin{remark}
        Jordan canonical form describes the structure that remains when a matrix
        cannot be diagonalized.
    \end{remark}

    \chapter{Differential Equations}

    \relatedcourses{(MATH200) Differential Equations}

    \relatedbooks{Elementary Differential Equations and Boundary Value Problems
        (William E. Boyce, Richard C. DiPrima, and Douglas B. Meade)}

    \begin{definition}[Differential Equation, Order, and Initial Value Problem]
        An ordinary differential equation (ODE) relates an unknown function of
        one independent variable to its derivatives. Its order is the highest
        derivative that occurs. An $n$th-order linear ODE has the form
        \[
            a_n(t)y^{(n)}+a_{n-1}(t)y^{(n-1)}+\cdots+a_1(t)y'+a_0(t)y=g(t),
        \]
        with $a_n$ not identically zero. It is homogeneous when $g=0$.
        An initial value problem (IVP) supplements the equation with enough
        values such as $y(t_0),y'(t_0),\dots$ to select a particular solution.
    \end{definition}

    \begin{definition}[Separable, Linear, Exact, Bernoulli, and Autonomous First-Order Equations]
        A first-order ODE is called \emph{separable} if it can be written as
        $y'=g(t)h(y)$, and \emph{linear} if it has the form
        $y'+p(t)y=q(t)$. An equation
        \[
            M(t,y)\dd t+N(t,y)\dd y=0
        \]
        is \emph{exact} on a region if there exists a potential $F$ with
        $F_t=M$ and $F_y=N$. A \emph{Bernoulli equation} has the form
        $y'+p(t)y=q(t)y^n$ with $n\neq0,1$. An \emph{autonomous equation}
        has the form $y'=f(y)$, with no explicit dependence on the independent
        variable. Recognizing the type determines the standard reduction to try.
    \end{definition}

    \begin{formula}[First-Order ODE Recognition Table]
        Before calculating, identify the structural form.

        {\small
        \renewcommand{\arraystretch}{1.16}
        \begin{longtable}{@{}L{0.20\textwidth}L{0.31\textwidth}L{0.41\textwidth}@{}}
            \toprule
            Type & Equation & Reduction or solution \\
            \midrule
            \endfirsthead

            \toprule
            Type & Equation & Reduction or solution \\
            \midrule
            \endhead

            \bottomrule
            \endlastfoot

            Separable & $y'=g(t)h(y)$ &
                $\displaystyle\int\frac{\dd y}{h(y)}=\int g(t)\dd t+C$ on
                intervals where $h(y)\neq0$; also check equilibrium solutions
                $h(y)=0$ \\
            Linear & $y'+p(t)y=q(t)$ & integrating factor
                $\displaystyle\mu(t)=e^{\int p(t)\dd t}$ \\
            Exact & $M(t,y)\dd t+N(t,y)\dd y=0$ & find $F$ with
                $F_t=M$, $F_y=N$, then $F(t,y)=C$ \\
            Bernoulli & $y'+p(t)y=q(t)y^n$, $n\neq0,1$ & substitute
                $u=y^{1-n}$ to obtain a linear equation \\
            Autonomous & $y'=f(y)$ & equilibria satisfy $f(y)=0$; away from
                equilibria the equation is separable \\
        \end{longtable}
        }
    \end{formula}

    \begin{theorem}[Picard--Lindel\"of Existence and Uniqueness Theorem]
        Consider the initial value problem
        \[
            y'=f(t,y),
            \qquad
            y(t_0)=y_0.
        \]
        If $f$ and $\partial f/\partial y$ are continuous near
        $(t_0,y_0)$, then there exists a unique solution near $t=t_0$.
    \end{theorem}

    \begin{remark}
        The decisive conclusion is existence and uniqueness.
        It tells us when an initial condition determines exactly one solution.
    \end{remark}

    \begin{theorem}[Superposition Principle]
        If $y_1,\dots,y_k$ are solutions of a homogeneous linear differential
        equation, then any linear combination
        \[
            c_1y_1+\cdots+c_ky_k
        \]
        is also a solution.
    \end{theorem}

    \begin{remark}
        The Superposition Principle is the differential-equation version of
        linear algebra. It is one of the main reasons linear equations are much
        easier to handle than nonlinear equations.
    \end{remark}

    \begin{proposition}[General Solution of a Nonhomogeneous Linear Equation]
        For a nonhomogeneous linear differential equation, the general solution
        has the form
        \[
            y=y_h+y_p,
        \]
        where $y_h$ is the general solution of the associated homogeneous
        equation and $y_p$ is one particular solution of the nonhomogeneous
        equation.
    \end{proposition}

    \begin{formula}[Separable First-Order Equation]
        A first-order equation of the form
        \[
            \frac{\dd y}{\dd x}=g(x)h(y)
        \]
        is separable. When $h(y)\neq 0$, rewrite it as
        \[
            \frac{1}{h(y)}\dd y=g(x)\dd x
        \]
        and integrate both sides.
    \end{formula}

    \begin{formula}[First-Order Linear Equation; Integrating Factor]
        The equation
        \[
            y'+p(t)y=q(t)
        \]
        has integrating factor
        \[
            \mu(t)=e^{\int p(t)\dd t}.
        \]
        Then
        \[
            (\mu y)'=\mu q,
        \]
        so
        \[
            y(t)=\frac{1}{\mu(t)}
            \left(
                \int \mu(t)q(t)\dd t+C
            \right).
        \]
    \end{formula}

    \begin{formula}[Exact Differential Equation]
        On a domain $D\subseteq\R^2$, an equation
        \[
            M(x,y)\dd x+N(x,y)\dd y=0
        \]
        is exact if there is a potential function $F\in C^1(D)$ s.t.
        \[
            F_x=M,
            \qquad
            F_y=N.
        \]
        If $M,N\in C^1(D)$ and $D$ is open and simply connected, then the
        equation is exact if and only if
        \[
            \frac{\partial M}{\partial y}
            =
            \frac{\partial N}{\partial x}.
        \]
        In that case, the solution is given implicitly by
        \[
            F(x,y)=C.
        \]
    \end{formula}

    \begin{formula}[Bernoulli Equation]
        The Bernoulli equation
        \[
            y'+p(t)y=q(t)y^n
            \qquad(n\neq 0,1)
        \]
        becomes a first-order linear equation after the substitution
        \[
            u=y^{1-n}.
        \]
    \end{formula}

    \begin{tactic}
        For a first-order ODE, first check whether it is separable, linear,
        exact, or Bernoulli. These four patterns cover many quick undergraduate
        differential-equation questions.
    \end{tactic}

    \begin{formula}[Second-Order Constant-Coefficient Equation]
        Consider the homogeneous equation
        \[
            ay''+by'+cy=0,
            \qquad
            a\neq 0.
        \]
        Its characteristic equation is
        \[
            ar^2+br+c=0.
        \]
        If the roots are distinct real numbers $r_1,r_2$, then
        \[
            y_h=c_1e^{r_1t}+c_2e^{r_2t}.
        \]
        If the equation has a repeated real root $r$, then
        \[
            y_h=(c_1+c_2t)e^{rt}.
        \]
        If the roots are $\alpha\pm i\beta$ with $\beta\neq 0$, then
        \[
            y_h=e^{\alpha t}
            \left(
                c_1\cos \beta t+c_2\sin \beta t
            \right).
        \]
    \end{formula}

    \begin{remark}
        The characteristic equation converts a constant-coefficient linear ODE
        into an algebraic equation. This is the key trick.
    \end{remark}

    \begin{definition}[Wronskian]
        For two differentiable functions $y_1$ and $y_2$, their Wronskian is
        \[
            W(y_1,y_2)(t)
            =
            \begin{vmatrix}
                y_1(t) & y_2(t) \\
                y_1'(t) & y_2'(t)
            \end{vmatrix}
            =
            y_1(t)y_2'(t)-y_1'(t)y_2(t).
        \]
    \end{definition}

    \begin{proposition}[Wronskian and Linear Independence]
        If
        \[
            W(y_1,y_2)(t_0)\neq 0
        \]
        for some $t_0$, then $y_1$ and $y_2$ are linearly independent.
    \end{proposition}

    \begin{formula}[Abel's Identity]
        For
        \[
            y''+p(t)y'+q(t)y=0,
        \]
        the Wronskian of two solutions satisfies
        \[
            W(t)=Ce^{-\int p(t)\dd t}.
        \]
        In particular, the Wronskian is either always zero or never zero on an
        interval where $p$ is continuous.
    \end{formula}

    \begin{formula}[Euler--Cauchy Equation]
        On an interval with $x>0$, the equation
        \[
            ax^2y''+bxy'+cy=0,
            \qquad a\neq0,
        \]
        suggests the trial solution $y=x^r$, giving the indicial equation
        \[
            ar(r-1)+br+c=0.
        \]
        If its roots are distinct real numbers $r_1,r_2$, then
        \[
            y=c_1x^{r_1}+c_2x^{r_2}.
        \]
        For a repeated root $r$,
        \[
            y=x^r(c_1+c_2\ln x).
        \]
        For roots $\alpha\pm i\beta$ with $\beta\neq0$,
        \[
            y=x^\alpha\bigl(c_1\cos(\beta\ln x)+c_2\sin(\beta\ln x)\bigr).
        \]
        On intervals with $x<0$, the analogous real formulas use $|x|$.
    \end{formula}

    \begin{algorithm}[Undetermined Coefficients]
        For a constant-coefficient equation with forcing term made from
        polynomials, exponentials, sines, or cosines, guess a particular solution
        of the same type. If the guess overlaps with $y_h$, multiply the guess
        by the smallest power of $t$ that removes the overlap.
    \end{algorithm}

    \begin{formula}[Variation of Parameters]
        For
        \[
            y''+p(t)y'+q(t)y=g(t)
        \]
        with fundamental solutions $y_1,y_2$ of the homogeneous equation, a
        particular solution is
        \[
            y_p
            =
            -y_1\int\frac{y_2g}{W}\dd t
            +
            y_2\int\frac{y_1g}{W}\dd t,
        \]
        where
        \[
            W=y_1y_2'-y_1'y_2.
        \]
    \end{formula}

    \begin{formula}[Exponential Growth and Decay]
        The equation
        \[
            y'=ky
        \]
        has solution
        \[
            y(t)=Ce^{kt}.
        \]
        If $k>0$, this models exponential growth. If $k<0$, this models
        exponential decay.
    \end{formula}

    \begin{formula}[Logistic Equation]
        The logistic equation is
        \[
            y'=ky\left(1-\frac{y}{K}\right),
            \qquad k>0,\quad K>0.
        \]
        Its equilibrium solutions are $y=0$ and $y=K$. Every non-equilibrium
        solution on an interval on which it is defined has the form
        \[
            y(t)=\frac{K}{1+Ce^{-kt}},
        \]
        where $C$ is determined by the initial condition. For
        $0<y(0)<K$, one has $C>0$ and $y(t)\to K$ as $t\to\infty$.
    \end{formula}

    \begin{remark}
        The logistic equation is the standard model for population growth with
        limited resources. The parameter $K$ represents the limiting population.
    \end{remark}

    \begin{formula}[Simple Harmonic Oscillator]
        The equation
        \[
            x''+\omega^2x=0
        \]
        has general solution
        \[
            x(t)=c_1\cos \omega t+c_2\sin \omega t.
        \]
        Equivalently,
        \[
            x(t)=A\cos(\omega t-\delta)
        \]
        for suitable constants $A$ and $\delta$.
    \end{formula}

    \begin{remark}
        This is the equation behind ideal mass--spring motion and many small
        oscillation models in physics.
    \end{remark}

    \begin{formula}[Damped Harmonic Oscillator]
        For $m>0$, $c\geq0$, and $k>0$, the equation
        \[
            mx''+cx'+kx=0
        \]
        has characteristic equation $mr^2+cr+k=0$. Put
        \[
            \alpha=\frac{c}{2m}.
        \]
        The three regimes and solution forms are
        \[
            c^2<4mk:
            \quad
            x(t)=e^{-\alpha t}
            \bigl(c_1\cos(\omega_dt)+c_2\sin(\omega_dt)\bigr),
            \quad
            \omega_d=\sqrt{\frac{k}{m}-\alpha^2},
        \]
        \[
            c^2=4mk:
            \quad
            x(t)=(c_1+c_2t)e^{-\alpha t},
        \]
        and, if $c^2>4mk$,
        \[
            x(t)=c_1e^{r_1t}+c_2e^{r_2t},
            \qquad
            r_{1,2}=\frac{-c\pm\sqrt{c^2-4mk}}{2m}.
        \]
        These are called underdamped, critically damped, and overdamped,
        respectively.
    \end{formula}

    \begin{definition}[Laplace Transform]
        The Laplace transform of a function $f(t)$ is
        \[
            \mathcal{L}\{f(t)\}(s)
            =
            \int_0^\infty e^{-st}f(t)\dd t.
        \]
    \end{definition}

    \begin{formula}[Laplace Transform Table]
        Let $F(s)=\mathcal{L}\{f(t)\}(s)$. The basic transforms are

        {\small
        \renewcommand{\arraystretch}{1.18}
        \begin{longtable}{@{}L{0.38\textwidth}L{0.42\textwidth}L{0.12\textwidth}@{}}
            \toprule
            $f(t)$ & $\mathcal{L}\{f(t)\}(s)$ & Conditions \\
            \midrule
            \endfirsthead

            \toprule
            $f(t)$ & $\mathcal{L}\{f(t)\}(s)$ & Conditions \\
            \midrule
            \endhead

            \bottomrule
            \endlastfoot

            $1$ & $\displaystyle\frac{1}{s}$ & $\Real(s)>0$ \\
            $t^n$ & $\displaystyle\frac{n!}{s^{n+1}}$ & $n\in\Nzero$ \\
            $e^{at}$ & $\displaystyle\frac{1}{s-a}$ & $\Real(s)>a$ \\
            $t^ne^{at}$ & $\displaystyle\frac{n!}{(s-a)^{n+1}}$ & $n\in\Nzero,\ s>a$ \\
            $\sin at$ & $\displaystyle\frac{a}{s^2+a^2}$ & $\Real(s)>0$ \\
            $\cos at$ & $\displaystyle\frac{s}{s^2+a^2}$ & $\Real(s)>0$ \\
            $e^{at}\sin bt$ & $\displaystyle\frac{b}{(s-a)^2+b^2}$ & $\Real(s)>a$ \\
            $e^{at}\cos bt$ & $\displaystyle\frac{s-a}{(s-a)^2+b^2}$ & $\Real(s)>a$ \\
            $\sinh at$ & $\displaystyle\frac{a}{s^2-a^2}$ & $\Real(s)>|a|$ \\
            $\cosh at$ & $\displaystyle\frac{s}{s^2-a^2}$ & $\Real(s)>|a|$ \\
            $f'(t)$ & $sF(s)-f(0)$ & \\
            $f''(t)$ & $s^2F(s)-sf(0)-f'(0)$ & \\
            $t f(t)$ & $-F'(s)$ & \\
            $\displaystyle\int_0^t f(\tau)\dd\tau$ & $\displaystyle\frac{F(s)}{s}$ & \\
            $u(t-a)f(t-a)$ & $e^{-as}F(s)$ & $a>0$ \\
            $f(t+T)=f(t)$ &
            $\displaystyle\frac{\int_0^T e^{-st}f(t)\dd t}{1-e^{-sT}}$ &
            $T>0,\ s>0$ \\
        \end{longtable}
        }
    \end{formula}

    \begin{proposition}[Laplace Shift and Convolution]
        If $F(s)=\mathcal{L}\{f(t)\}(s)$, then
        \[
            \mathcal{L}\{e^{at}f(t)\}(s)=F(s-a).
        \]
        If
        \[
            (f*g)(t)=\int_0^t f(\tau)g(t-\tau)\dd\tau,
        \]
        then
        \[
            \mathcal{L}\{f*g\}=\mathcal{L}\{f\}\mathcal{L}\{g\}.
        \]
    \end{proposition}

    \begin{formula}[Fourier Series]
        For a piecewise smooth function $f$ on $[-L,L]$,
        \[
            f(x)
            \sim
            \frac{a_0}{2}
            +\sum_{n=1}^{\infty}
            \left(
                a_n\cos\frac{n\pi x}{L}
                +b_n\sin\frac{n\pi x}{L}
            \right),
        \]
        where
        \[
            a_n
            =
            \frac{1}{L}\int_{-L}^{L}
            f(x)\cos\frac{n\pi x}{L}\dd x,
            \qquad
            b_n
            =
            \frac{1}{L}\int_{-L}^{L}
            f(x)\sin\frac{n\pi x}{L}\dd x.
        \]
        If $f$ is even, then $b_n=0$; if $f$ is odd, then $a_n=0$.
        At a jump discontinuity, the series converges to the midpoint of the
        left and right limits.
    \end{formula}

    \begin{formula}[Classical Fourier Series]
        For $-\pi<x<\pi$,
        \[
            x
            =
            2\sum_{n=1}^{\infty}
            \frac{(-1)^{n+1}}{n}\sin(nx),
        \]
        \[
            x^2
            =
            \frac{\pi^2}{3}
            +4\sum_{n=1}^{\infty}
            \frac{(-1)^n}{n^2}\cos(nx),
        \]
        and, for $0<x<\pi$,
        \[
            1
            =
            \frac{4}{\pi}
            \sum_{k=0}^{\infty}
            \frac{\sin((2k+1)x)}{2k+1}.
        \]
    \end{formula}

    \begin{definition}[Fourier Transform]
        For an integrable function $f\colon \R\to\C$, define
        \[
            \widehat f(\xi)
            =
            \int_{-\infty}^{\infty}f(t)e^{-2\pi i t\xi}\,\dd t.
        \]
        Under suitable hypotheses, the inversion formula is
        \[
            f(t)
            =
            \int_{-\infty}^{\infty}\widehat f(\xi)e^{2\pi i t\xi}\,\dd\xi.
        \]
    \end{definition}

    \begin{formula}[Fourier Transform Table]
        With the convention above, the standard identities are

        {\small
        \renewcommand{\arraystretch}{1.18}
        \begin{longtable}{@{}L{0.42\textwidth}L{0.46\textwidth}@{}}
            \toprule
            Function or operation & Fourier transform \\
            \midrule
            \endfirsthead

            \toprule
            Function or operation & Fourier transform \\
            \midrule
            \endhead

            \bottomrule
            \endlastfoot

            $f'(t)$ & $2\pi i\xi\widehat f(\xi)$ \\
            $f(t-a)$ & $e^{-2\pi i a\xi}\widehat f(\xi)$ \\
            $e^{2\pi iat}f(t)$ & $\widehat f(\xi-a)$ \\
            $f(at)$, $a\neq0$ & $\displaystyle\frac{1}{|a|}\widehat f\left(\frac{\xi}{a}\right)$ \\
            $t f(t)$ & $\displaystyle\frac{i}{2\pi}\frac{\dd}{\dd\xi}\widehat f(\xi)$ \\
            $(f*g)(t)$ & $\widehat f(\xi)\widehat g(\xi)$ \\
            $e^{-at^2}$, $a>0$ & $\displaystyle\sqrt{\frac{\pi}{a}}e^{-\pi^2\xi^2/a}$ \\
        \end{longtable}
        }
    \end{formula}

    \begin{formula}[Linear Systems and Matrix Exponential]
        The homogeneous linear system
        \[
            \mathbf{x}'=\mathbf{A}\mathbf{x}
        \]
        has solution
        \[
            \mathbf{x}(t)=e^{t\mathbf{A}}\mathbf{x}(0).
        \]
        If $\mathbf{A}$ is diagonalizable as
        $\mathbf{A}=\mathbf{S}\mathbf{D}\mathbf{S}^{-1}$, then
        \[
            e^{t\mathbf{A}}
            =
            \mathbf{S}e^{t\mathbf{D}}\mathbf{S}^{-1},
            \qquad
            e^{t\mathbf{D}}
            =
            \diag(e^{\lambda_1t},\dots,e^{\lambda_nt}).
        \]
    \end{formula}

    \begin{strategy}
        For a constant-coefficient linear system, the eigenvalues of
        $\mathbf{A}$ control the qualitative behavior. Negative real parts
        indicate decay, positive real parts indicate growth, and nonzero
        imaginary parts indicate rotation or oscillation.
    \end{strategy}

    \begin{remark}
        The Laplace transform turns many differential equations into algebraic
        equations. This is its main computational advantage.
    \end{remark}

    \begin{definition}[Heat Equation]
        The heat equation is the partial differential equation
        \[
            u_t=\alpha u_{xx},
            \qquad
            \alpha>0.
        \]
        It models diffusion, especially the flow of heat in a one-dimensional
        medium.
    \end{definition}

    \begin{definition}[Wave Equation]
        The one-dimensional wave equation is
        \[
            u_{tt}=c^2u_{xx}.
        \]
        It models waves traveling with speed $c$.
    \end{definition}

    \begin{theorem}[d'Alembert's Formula]
        Let $c>0$, let $f\in C^2(\R)$, and let $g\in C^1(\R)$. The classical
        solution of
        \[
            u_{tt}=c^2u_{xx},
            \qquad
            u(x,0)=f(x),
            \qquad
            u_t(x,0)=g(x),
        \]
        is
        \[
            u(x,t)
            =
            \frac{f(x-ct)+f(x+ct)}{2}
            +
            \frac{1}{2c}\int_{x-ct}^{x+ct}g(s)\dd s.
        \]
    \end{theorem}

    \begin{remark}
        d'Alembert's formula shows that waves move along the characteristic
        lines $x-ct=\text{constant}$ and $x+ct=\text{constant}$.
    \end{remark}

    \begin{definition}[Laplace Equation]
        Laplace's equation is
        \[
            \Delta u=0.
        \]
        In two variables, this means
        \[
            u_{xx}+u_{yy}=0.
        \]
        A function satisfying Laplace's equation is called harmonic.
    \end{definition}

    \begin{definition}[Poisson Equation]
        Poisson's equation is
        \[
            \Delta u=f.
        \]
        It is an inhomogeneous version of Laplace's equation.
    \end{definition}

    \begin{remark}
        The heat equation, wave equation, and Laplace equation are three fundamental
        partial differential equations with distinct physical interpretations.
    \end{remark}

    \chapter{Combinatorics and Discrete Mathematics}

    \relatedcourses{(MATH261) Discrete Mathematics}

    \relatedbooks{Introductory Combinatorics (Richard A. Brualdi),
        Invitation to Discrete Mathematics (Ji\v{r}\'i Matou\v{s}ek and Jaroslav Ne\v{s}et\v{r}il)}

    \begin{definition}[Permutation, Combination, and Multiset Selection]
        An ordered selection of $k$ distinct objects from $n$ distinct objects
        is a $k$-permutation, counted by
        \[
            P(n,k)=\frac{n!}{(n-k)!}.
        \]
        An unordered $k$-element selection is a $k$-combination, counted by
        \[
            \binom{n}{k}.
        \]
        If repetition is allowed and order is ignored, the number of size-$k$
        multisets from $n$ types is
        \[
            \multichoose{n}{k}=\binom{n+k-1}{k}.
        \]
        These three models should be distinguished before applying a formula.
    \end{definition}

    \begin{formula}[Binomial Theorem; Newton's Binomial Formula]
        For every nonnegative integer $n$,
        \[
            (x+y)^n
            =
            \sum_{k=0}^{n}\binom{n}{k}x^{n-k}y^k.
        \]
        In particular,
        \[
            (1+x)^n
            =
            \sum_{k=0}^{n}\binom{n}{k}x^k.
        \]
    \end{formula}

    \begin{remark}
        The coefficients $\binom{n}{k}$ are the entries of Pascal's triangle.
        This is one of the most recognizable bridges between algebra and
        counting.
    \end{remark}

    \begin{strategy}[Bijection, Double Counting, and Complementary Counting]
        Before expanding a complicated sum, try to count the same objects in a
        second way. Bijections, double counting, and counting the complement
        are among the fastest and most reusable contest methods.
    \end{strategy}

    \begin{formula}[Binomial Inversion]
        The relations
        \[
            b_n=\sum_{k=0}^{n}\binom{n}{k}a_k
        \]
        and
        \[
            a_n=\sum_{k=0}^{n}(-1)^{n-k}\binom{n}{k}b_k
        \]
        are equivalent.
    \end{formula}

    \begin{formula}[Pascal's Identity and Vandermonde's Identity]
        Pascal's identity is
        \[
            \binom{n}{k}
            =
            \binom{n-1}{k-1}+\binom{n-1}{k}.
        \]
        Vandermonde's identity is
        \[
            \sum_{k=0}^{r}\binom{m}{k}\binom{n}{r-k}
            =
            \binom{m+n}{r}.
        \]
    \end{formula}

    \begin{formula}[Multinomial Theorem]
        For nonnegative integers $k_1,\dots,k_m$ with
        $k_1+\cdots+k_m=n$,
        \[
            (x_1+\cdots+x_m)^n
            =
            \sum_{k_1+\cdots+k_m=n}
            \frac{n!}{k_1!\cdots k_m!}
            x_1^{k_1}\cdots x_m^{k_m}.
        \]
        The coefficient
        \[
            \frac{n!}{k_1!\cdots k_m!}
        \]
        is called a multinomial coefficient.
    \end{formula}

    \begin{formula}[Stars and Bars]
        The number of nonnegative integer solutions of
        \[
            x_1+x_2+\cdots+x_n=k
        \]
        is
        \[
            \binom{n+k-1}{k}
            =
            \multichoose{n}{k}.
        \]
        Equivalently, the expansion of $(x_1+x_2+\cdots+x_n)^k$ has
        \[
            \binom{n+k-1}{k}
        \]
        distinct monomial terms.
    \end{formula}

    \begin{formula}[Lattice Path Counting]
        The number of shortest lattice paths from $(0,0)$ to $(m,n)$ using only
        unit steps to the right and upward is
        \[
            \binom{m+n}{m}
            =
            \binom{m+n}{n}.
        \]
    \end{formula}

    \begin{theorem}[Reflection Principle and Ballot-Path Formula]
        The reflection principle converts a path that first crosses a boundary
        into an unrestricted reflected path. In particular, for integers
        $a\geq b\geq0$, the number of paths from $(0,0)$ to $(a,b)$ using right
        and up steps and never entering the region $y>x$ is
        \[
            \binom{a+b}{b}-\binom{a+b}{b-1}
            =
            \frac{a-b+1}{a+1}\binom{a+b}{b}.
        \]
        When $a=b=n$, this becomes the Catalan number
        \[
            \frac{1}{n+1}\binom{2n}{n}.
        \]
    \end{theorem}

    \begin{strategy}[Tail Switching for Intersecting Paths]
        When two directed lattice paths meet, exchange their tails after the
        first intersection. This produces a bijection with a pair of paths whose
        endpoints have been crossed. The method is a standard way to count or
        cancel intersecting path pairs.
    \end{strategy}

    \begin{formula}[Ordinary Generating Function Table]
        For a sequence $(a_n)_{n\geq 0}$, its ordinary generating function is
        \[
            A(x)=\sum_{n=0}^{\infty}a_nx^n.
        \]
        The basic examples are

        {\small
        \renewcommand{\arraystretch}{1.18}
        \begin{longtable}{@{}L{0.31\textwidth}L{0.43\textwidth}L{0.18\textwidth}@{}}
            \toprule
            Sequence $a_n$ & Ordinary generating function $A(x)$ & Conditions \\
            \midrule
            \endfirsthead

            \toprule
            Sequence $a_n$ & Ordinary generating function $A(x)$ & Conditions \\
            \midrule
            \endhead

            \bottomrule
            \endlastfoot

            $1$ & $\displaystyle\frac{1}{1-x}$ & $|x|<1$ \\
            $r^n$ & $\displaystyle\frac{1}{1-rx}$ & $|rx|<1$ \\
            $n$ & $\displaystyle\frac{x}{(1-x)^2}$ & $|x|<1$ \\
            $n^2$ & $\displaystyle\frac{x(1+x)}{(1-x)^3}$ & $|x|<1$ \\
            $\displaystyle\binom{n}{k}$ &
            $\displaystyle\frac{x^k}{(1-x)^{k+1}}$ &
            $k\in\Nzero,\ |x|<1$ \\
            $\displaystyle\binom{n+r}{r}$ &
            $\displaystyle\frac{1}{(1-x)^{r+1}}$ &
            $r\in\Nzero,\ |x|<1$ \\
            $F_n$, where $F_0=0$ and $F_1=1$ &
            $\displaystyle\frac{x}{1-x-x^2}$ &
            $|x|<\frac{\sqrt{5}-1}{2}$ \\
            $\displaystyle\frac{1}{n+1}\binom{2n}{n}$ &
            $\displaystyle\frac{1-\sqrt{1-4x}}{2x}$ &
            $|x|<\frac14$ \\
        \end{longtable}
        }
    \end{formula}

    \begin{tactic}
        Generating functions are useful when a recurrence or counting problem has a
        repeated structure; a known generating function can replace a long
        derivation.
    \end{tactic}

    \begin{formula}[Stirling Numbers and Bell Numbers]
        The Stirling number of the second kind
        \[
            \stirlingsecond{n}{k}
        \]
        counts partitions of an $n$-element set into $k$ nonempty blocks and
        satisfies
        \[
            \stirlingsecond{n}{k}
            =
            k\stirlingsecond{n-1}{k}
            +
            \stirlingsecond{n-1}{k-1}.
        \]
        The Bell number
        \[
            B_n=\sum_{k=0}^{n}\stirlingsecond{n}{k}
        \]
        counts all set partitions of an $n$-element set.
    \end{formula}

    \begin{formula}[Fibonacci Numbers; Binet's Formula]
        If $F_0=0$, $F_1=1$, and $F_{n+1}=F_n+F_{n-1}$, then
        \[
            F_n
            =
            \frac{\varphi^n-\psi^n}{\sqrt{5}},
            \qquad
            \varphi=\frac{1+\sqrt{5}}{2},
            \qquad
            \psi=\frac{1-\sqrt{5}}{2}.
        \]
        In particular,
        \[
            \lim\limits_{n\to\infty}\frac{F_{n+1}}{F_n}=\varphi.
        \]
    \end{formula}

    \begin{formula}[Roots-of-Unity Filter]
        Let $\omega_m=e^{2\pi i/m}$. Then
        \[
            \frac{1}{m}\sum_{j=0}^{m-1}\omega_m^{j(n-r)}
            =
            \begin{cases}
                1, & n\equiv r\pmod m, \\
                0, & n\not\equiv r\pmod m.
            \end{cases}
        \]
        Consequently, if $A(x)=\sum_{n\geq0}a_nx^n$, then the terms with
        indices congruent to $r$ modulo $m$ are extracted by
        \[
            \sum_{n\equiv r\, (\mathrm{mod}\,m)}a_nx^n
            =
            \frac1m\sum_{j=0}^{m-1}\omega_m^{-rj}A(\omega_m^j x).
        \]
    \end{formula}

    \begin{formula}[Tower of Hanoi]
        Let $T_n$ be the minimum number of moves needed to move $n$ disks in the
        Tower of Hanoi puzzle. Then
        \[
            T_n=2T_{n-1}+1,
            \qquad
            T_1=1,
        \]
        and hence
        \[
            T_n=2^n-1.
        \]
    \end{formula}

    \begin{remark}
        The recurrence comes from moving the top $n-1$ disks away, moving the
        largest disk once, and then moving the $n-1$ disks back on top of it.
    \end{remark}

    \begin{strategy}[State Compression and Transfer Matrices]
        When a counting process depends only on finitely many recent states,
        place the transition counts in a matrix $\mathbf{T}$. Counts of length
        $n$ then have the form
        \[
            \mathbf{u}^{\mathsf T}\mathbf{T}^n\mathbf{v}.
        \]
        This is effective for forbidden substrings, tilings, walks, and finite
        automata.
    \end{strategy}

    \begin{theorem}[Pigeonhole Principle; Dirichlet's Box Principle]
        If $n+1$ objects are placed into $n$ boxes, then at least one box
        contains at least two objects.
        More generally, if $N$ objects are placed into $k$ boxes, then at least
        one box contains at least
        \[
            \left\lceil\frac{N}{k}\right\rceil
        \]
        objects.
    \end{theorem}

    \begin{strategy}
        The Pigeonhole Principle is often hidden in problems asking us to prove
        that two objects must share a property. The hard part is usually
        choosing the right ``boxes.''
    \end{strategy}

    \begin{theorem}[Consecutive-Sum Divisibility Lemma]
        For any integers $a_1,\dots,a_n$, there exists a nonempty consecutive
        block whose sum is divisible by $n$. Define prefix sums
        \[
            S_0=0,
            \qquad
            S_k=\sum_{i=1}^{k}a_i.
        \]
        If some $S_k\equiv0\pmod n$, use the first $k$ terms. Otherwise, two
        of $S_1,\dots,S_n$ have the same residue modulo $n$, and their difference
        is the required block sum.
    \end{theorem}

    \begin{definition}[Graphs, Walks, Paths, Trails, Cycles, and Trees]
        Unless stated otherwise, a graph $G=(V,E)$ here is a finite undirected
        graph with vertex set $V$ and edge set $E$. Two vertices joined by an
        edge are adjacent, and the degree $\deg(v)$ is the number of edges
        incident to $v$. A walk is a sequence of adjacent vertices. A trail
        repeats no edge, while a path repeats no vertex. A closed trail is a
        circuit, and a cycle is a closed path apart from its repeated initial
        and terminal vertex. A graph is connected if every two vertices are
        joined by a path. A tree is a connected graph containing no cycle.

        An Eulerian trail uses every edge exactly once; an Eulerian circuit is
        a closed Eulerian trail. A Hamiltonian path visits every vertex exactly
        once; a Hamiltonian cycle returns to its starting vertex. A graph is
        bipartite if its vertices can be divided into two classes with every
        edge joining different classes. A planar graph can be drawn in the
        plane without edge crossings except at common endpoints. The complete
        graph $K_n$ contains every edge between distinct vertices, and the
        complete bipartite graph $K_{m,n}$ contains every edge between its two
        vertex classes. Thus a one-stroke drawing problem is an Eulerian, not a
        Hamiltonian, question.
    \end{definition}

    \begin{strategy}[Coloring, Parity, and Bipartite Obstructions]
        Colorings convert geometric or movement constraints into parity data.
        On a bipartite graph, every path alternates colors. Thus a Hamiltonian
        path can exist only if the two color classes differ in size by at most
        $1$, and a Hamiltonian cycle requires the two classes to have equal size.
        Checkerboard colorings are especially effective for tiling and movement
        problems.
    \end{strategy}

    \begin{theorem}[Inclusion--Exclusion Principle; Sieve Formula]
        For finite sets $A_1,\dots,A_n$,
        \[
            \left|\bigcup_{i=1}^{n}A_i\right|
            =
            \sum_i |A_i|
            -
            \sum_{i<j}|A_i\cap A_j|
            +
            \sum_{i<j<k}|A_i\cap A_j\cap A_k|
            -\cdots
            +(-1)^{n+1}|A_1\cap\cdots\cap A_n|.
        \]
    \end{theorem}

    \begin{remark}
        Inclusion--Exclusion is the standard correction mechanism for
        overcounting. It is also historically connected to sieve methods.
    \end{remark}

    \begin{formula}[Derangement Formula and Recurrence]
        Let $D_n$ be the number of derangements of $n$ objects, that is,
        permutations with no fixed points. Then
        \[
            D_n=n!\sum_{k=0}^{n}\frac{(-1)^k}{k!}
            =\left\lfloor \frac{n!}{e}+\frac12\right\rfloor.
        \]
        The characteristic recurrence is
        \[
            D_n=(n-1)(D_{n-1}+D_{n-2})
            \qquad(n\geq2),
        \]
        with $D_0=1$ and $D_1=0$. Equivalently,
        \[
            D_n=nD_{n-1}+(-1)^n.
        \]
        The first recurrence comes from choosing the image of a distinguished
        object and separating the cases in which a resulting transposition is
        completed or not.
    \end{formula}

    \begin{remark}
        The approximation $D_n\approx n!/e$ is very memorable. In probability
        language, a random permutation has no fixed point with probability
        approximately $1/e$.
    \end{remark}

    \begin{formula}[Catalan Numbers]
        The $n$th Catalan number is
        \[
            C_n
            =
            \frac{1}{n+1}\binom{2n}{n}.
        \]
        It begins as
        \[
            1,\ 1,\ 2,\ 5,\ 14,\ 42,\ 132,\dots.
        \]
    \end{formula}

    \begin{remark}
        Catalan numbers count many different objects, including correct
        parenthesizations, Dyck paths, triangulations of a polygon, and rooted
        full binary trees. The number $42$ is a standard recognition item.
    \end{remark}


    \begin{formula}[Cayley's Formula]
        The number of labeled trees on $n$ vertices is
        \[
            n^{n-2}
            \qquad(n\geq2).
        \]
    \end{formula}

    \begin{remark}
        Cayley's formula gives a compact and surprising enumeration of labeled
        trees.
    \end{remark}

    \begin{definition}[Finite Symmetry Action, Orbit, and Fixed Point]
        For counting purposes, let $G$ be a finite collection of symmetries that
        contains the identity and is closed under composition and inverses. An
        action of $G$ on a finite set $X$ assigns to each $g\in G$ a permutation
        of $X$ compatibly with composition. The orbit of $x\in X$ is the set of
        objects obtainable from $x$ by these symmetries, and
        \[
            \Fix(g)=\{x\in X\mid g\cdot x=x\}
        \]
        is the fixed-point set of $g$. Chapter 9 gives the full abstract group
        definition and the general orbit--stabilizer framework.
    \end{definition}

    \begin{theorem}[Burnside's Lemma; Cauchy--Frobenius--Burnside Lemma]
        Suppose a finite group $G$ acts on a finite set $X$, meaning that each
        $g\in G$ acts as a permutation of $X$ and the action is compatible with
        group multiplication. The number of orbits is
        \[
            \frac{1}{|G|}
            \sum_{g\in G}|\Fix(g)|,
        \]
        where
        \[
            \Fix(g)
            =
            \{x\in X\mid g\cdot x=x\}.
        \]
    \end{theorem}

    \begin{idea}
        Burnside's Lemma says that the number of essentially different objects
        is the average number of objects fixed by a symmetry.
    \end{idea}

    \begin{theorem}[Equivalent Characterizations of a Tree]
        For a finite graph $G$ with $n$ vertices, any two of the following
        conditions imply the third, and in fact all three are equivalent:
        \[
            G\text{ is connected and acyclic},
        \]
        \[
            G\text{ is connected and }|E|=n-1,
        \]
        and
        \[
            \text{every two vertices are joined by a unique path}.
        \]
        Equivalently, a tree is minimally connected and maximally acyclic.
    \end{theorem}

    \begin{theorem}[Bipartite Graph Criterion]
        A graph is bipartite if and only if it contains no odd cycle. Therefore
        a two-coloring obstruction and an odd-cycle obstruction are the same
        phenomenon.
    \end{theorem}

    \begin{lemma}[Handshaking Lemma]
        For a finite undirected graph $G=(V,E)$,
        \[
            \sum_{v\in V}\deg(v)=2|E|.
        \]
        In particular, the number of vertices of odd degree is even.
    \end{lemma}

    \begin{remark}
        The name comes from the idea that every handshake contributes $2$ to
        the total count of hands shaken.
    \end{remark}

    \begin{theorem}[Euler's Formula for Planar Graphs]
        If a connected planar graph has $V$ vertices, $E$ edges, and $F$ faces,
        then
        \[
            V-E+F=2.
        \]
    \end{theorem}

    \begin{remark}
        This is the graph-theoretic version of Euler's polyhedron formula.
        It is one of the most famous formulas involving planar graphs.
    \end{remark}


    \begin{theorem}[Kuratowski's Theorem]
        A finite graph is planar if and only if it contains no subgraph that is
        a subdivision of $K_5$ or $K_{3,3}$. Thus the two classical complete
        graphs $K_5$ and $K_{3,3}$ are the fundamental obstructions to
        planarity in the sense of subdivisions.
    \end{theorem}

    \begin{formula}[Regions Cut by Hyperplanes in General Position]
        The maximum number of regions into which $n$ hyperplanes in general
        position divide $d$-dimensional space is
        \[
            R_d(n)=\sum_{k=0}^{d}\binom{n}{k}.
        \]
        In particular, $n$ lines divide the plane into at most
        \[
            1+n+\binom{n}{2}
        \]
        regions, and $n$ planes divide $\R^3$ into at most
        \[
            1+n+\binom{n}{2}+\binom{n}{3}
        \]
        regions.
    \end{formula}

    \begin{formula}[Chords and Diagonals in Convex Position]
        If $n$ points lie on a circle and no three interior chords are concurrent,
        drawing every chord divides the disk into
        \[
            1+\binom{n}{2}+\binom{n}{4}
        \]
        regions. A convex $n$-gon has at most
        \[
            \binom{n}{4}
        \]
        interior intersection points of diagonals.
    \end{formula}

    \begin{formula}[Triangulation with Interior Vertices]
        Suppose a polygonal disk is triangulated using $V$ vertices in total,
        of which $b$ lie on the boundary. Then the number of edges and triangular
        faces are
        \[
            E=3V-b-3,
            \qquad
            F_{\triangle}=2V-b-2.
        \]
        These values follow from Euler's formula and double counting the
        triangle-edge incidences.
    \end{formula}

    \begin{strategy}[Double Counting with Euler's Formula]
        In a planar subdivision, count edge-face incidences first and then use
        $V-E+F=2$. Interior edges are counted twice and boundary edges once.
        This is often faster and safer than trying to enumerate faces directly.
    \end{strategy}

    \begin{theorem}[Eulerian Trail Criterion]
        A connected finite graph has an Eulerian circuit if and only if every
        vertex has even degree. It has an Eulerian trail but not an Eulerian
        circuit if and only if exactly two vertices have odd degree.
    \end{theorem}

    \begin{remark}
        This theorem is historically connected to Euler's solution of the
        Seven Bridges of K\"onigsberg problem, which marks the beginning of
        graph theory and topology. The historical city K\"onigsberg is now
        Kaliningrad, Russia.
    \end{remark}

    \begin{definition}[Matching, Perfect Matching, Vertex Cover, and Neighborhood]
        In a graph, a \emph{matching} is a set of edges no two of which share a
        vertex. A matching \emph{covers} a vertex if one of its edges is incident
        to that vertex, and it is \emph{perfect} if it covers every vertex. A
        \emph{vertex cover} is a set of vertices meeting every edge. For
        $S\subseteq V$, its neighborhood is
        \[
            N(S)=\{v\in V\mid v\text{ is adjacent to some }s\in S\}.
        \]
        These notions are the basic language of assignment and pairing problems.
    \end{definition}

    \begin{theorem}[Hall's Marriage Theorem]
        Let $G=(X\cup Y,E)$ be a finite bipartite graph. There is a matching
        that covers every vertex of $X$ if and only if, for every subset
        $S\subseteq X$,
        \[
            |N(S)|\geq |S|,
        \]
        where $N(S)$ is the set of neighbors of vertices in $S$.
    \end{theorem}

    \begin{strategy}
        Hall's Marriage Theorem is the standard named theorem for proving that
        a perfect assignment or matching exists.
    \end{strategy}

    \begin{theorem}[K\H{o}nig's Theorem]
        In every finite bipartite graph, the size of a maximum matching equals
        the size of a minimum vertex cover.
    \end{theorem}

    \begin{definition}[Capacitated Network, Flow, and Cut]
        A capacitated network is a directed graph with a distinguished source
        $s$, sink $t$, and nonnegative capacity $c(e)$ on each edge. An
        $s$--$t$ flow assigns values $f(e)$ satisfying
        $0\leq f(e)\leq c(e)$ and flow conservation at every vertex other than
        $s$ and $t$. An $s$--$t$ cut is a partition $V=S\sqcup T$ with
        $s\in S$ and $t\in T$; its capacity is the sum of capacities of edges
        directed from $S$ to $T$.
    \end{definition}

    \begin{theorem}[Max-Flow Min-Cut Theorem]
        In a finite capacitated network, the maximum value of a source--sink
        flow equals the minimum capacity of a source--sink cut.
    \end{theorem}

    \begin{theorem}[Menger's Theorem, Vertex Form]
        In a finite graph, for nonadjacent vertices $u$ and $v$, the maximum
        number of pairwise internally vertex-disjoint $u$--$v$ paths equals the
        minimum number of vertices whose deletion separates $u$ from $v$.
    \end{theorem}

    \begin{theorem}[Kirchhoff's Matrix--Tree Theorem]
        Let
        \[
            \mathbf{L}=\mathbf{D}-\mathbf{A}
        \]
        be the Laplacian matrix of a finite graph. The number of spanning trees
        equals any cofactor of $\mathbf{L}$ obtained by deleting the same row
        and column.
    \end{theorem}

    \begin{definition}[Partially Ordered Set, Chain, and Antichain]
        A partially ordered set, or poset, is a set $P$ equipped with a relation
        $\preceq$ that is reflexive, antisymmetric, and transitive. A
        \emph{chain} is a subset whose elements are pairwise comparable, while
        an \emph{antichain} is a subset in which no two distinct elements are
        comparable. The power set $\mathcal{P}([n])$ ordered by inclusion is the
        standard finite example.
    \end{definition}

    \begin{theorem}[Sperner's Theorem]
        The largest family of subsets of $[n]$ in which no member contains
        another has size
        \[
            \binom{n}{\lfloor n/2\rfloor}.
        \]
    \end{theorem}

    \begin{theorem}[Dilworth's Theorem]
        In a finite partially ordered set, the maximum size of an antichain
        equals the minimum number of chains covering the poset.
    \end{theorem}

    \begin{theorem}[Tur\'an's Theorem]
        Among all $n$-vertex graphs containing no $K_{r+1}$, the maximum number
        of edges is attained by the complete $r$-partite graph whose part sizes
        differ by at most $1$.
    \end{theorem}

    \begin{theorem}[Dirac's and Ore's Hamiltonian Criteria]
        Let $G$ be a simple graph on $n\geq3$ vertices. If every vertex has
        degree at least $n/2$, then $G$ is Hamiltonian. More generally, it is
        enough that
        \[
            \deg(u)+\deg(v)\geq n
        \]
        for every pair of nonadjacent vertices $u,v$.
    \end{theorem}

    \begin{theorem}[Ramsey's Theorem]
        For any positive integers $r$ and $s$, there exists an integer
        $R(r,s)$ s.t. every red-blue coloring of the edges of a sufficiently
        large complete graph contains either a red $K_r$ or a blue $K_s$.
    \end{theorem}

    \begin{formula}[The Party Problem]
        The classical Ramsey number
        \[
            R(3,3)=6
        \]
        says that among any six people, there are either three mutual
        acquaintances or three mutual strangers.
    \end{formula}

    \begin{remark}
        Ramsey theory is the mathematics of unavoidable structure. Its slogan
        is that complete disorder is impossible in a large enough system.
    \end{remark}

    \begin{theorem}[Erd\H{o}s--Szekeres Theorem]
        Any sequence of
        \[
            (r-1)(s-1)+1
        \]
        distinct real numbers contains either an increasing subsequence of
        length $r$ or a decreasing subsequence of length $s$.
    \end{theorem}

    \begin{remark}
        This is a beautiful pigeonhole-style result. It is often remembered as
        the monotone subsequence theorem.
    \end{remark}

    \chapter{Probability and Statistics}

    \relatedcourses{(MATH230) Probability \& Statistics}

    \relatedbooks{Probability \& Statistics for Engineers \& Scientists
        (Ronald E. Walpole, Raymond H. Myers, Sharon L. Myers, and Keying Ye)}

    \begin{definition}[Probability Space, Event, and Conditional Probability]
        A probability space is a triple $(\Omega,\mathcal{F},\Prob)$, where
        $\Omega$ is the sample space, $\mathcal{F}$ is a $\sigma$-algebra of
        events (closed under complements and countable unions), and
        $\Prob\colon \mathcal{F}\to[0,1]$ satisfies
        \[
            \Prob(\Omega)=1,
            \qquad
            \Prob\!\left(\bigsqcup_{i=1}^{\infty}A_i\right)
            =
            \sum_{i=1}^{\infty}\Prob(A_i)
        \]
        for pairwise disjoint events $A_i\in\mathcal{F}$. If
        $\Prob(B)>0$, the conditional probability of $A$ given $B$ is
        \[
            \Prob(A\mid B)
            =
            \frac{\Prob(A\cap B)}{\Prob(B)}.
        \]
    \end{definition}

    \begin{definition}[Random Variable, Distribution, and Distribution Function]
        A random variable is a measurable function $X\colon \Omega\to\R$. Its
        distribution is the probability law of its values. For a discrete
        random variable, the probability mass function is
        \[
            p_X(x)=\Prob(X=x),
            \qquad
            \sum_x p_X(x)=1.
        \]
        For a continuous random variable with density $f_X$,
        \[
            \Prob(a<X\leq b)=\int_a^b f_X(x)\,\dd x,
            \qquad
            \int_{-\infty}^{\infty}f_X(x)\,\dd x=1.
        \]
        In either case the cumulative distribution function is
        \[
            F_X(x)=\Prob(X\leq x).
        \]
    \end{definition}

    \begin{definition}[Expectation and Variance]
        If the relevant sum or integral converges absolutely, then
        \[
            \Exp[X]
            =
            \begin{cases}
                \displaystyle\sum_x x\,\Prob(X=x), & X\text{ discrete},\\[1mm]
                \displaystyle\int_{-\infty}^{\infty}x f_X(x)\,\dd x,
                & X\text{ continuous},
            \end{cases}
        \]
        and, whenever the second moment is finite,
        \[
            \Var(X)=\Exp\!\left[(X-\Exp[X])^2\right].
        \]
    \end{definition}

    \begin{definition}[Joint, Marginal, and Conditional Distributions; Conditional Expectation]
        For discrete random variables $X$ and $Y$, the joint probability mass
        function is
        \[
            p_{X,Y}(x,y)=\Prob(X=x,Y=y),
        \]
        and the marginal mass functions are obtained by summing over the other
        variable. If $\Prob(Y=y)>0$, then
        \[
            \Prob(X=x\mid Y=y)
            =
            \frac{p_{X,Y}(x,y)}{\Prob(Y=y)},
        \]
        and, whenever the sum converges absolutely,
        \[
            \Exp[X\mid Y=y]
            =
            \sum_x x\Prob(X=x\mid Y=y).
        \]

        If $(X,Y)$ has joint density $f_{X,Y}$, then the marginal density of
        $Y$ is
        \[
            f_Y(y)=\int_{-\infty}^{\infty}f_{X,Y}(x,y)\dd x.
        \]
        Whenever $f_Y(y)>0$,
        \[
            f_{X\mid Y}(x\mid y)
            =
            \frac{f_{X,Y}(x,y)}{f_Y(y)},
            \qquad
            \Exp[X\mid Y=y]
            =
            \int_{-\infty}^{\infty}x f_{X\mid Y}(x\mid y)\dd x.
        \]
        The random variable $\Exp[X\mid Y]$ is obtained by evaluating this
        conditional mean at the observed value of $Y$.
    \end{definition}

    \begin{theorem}[Law of Total Probability]
        If $B_1,\dots,B_n$ form a partition of the sample space and
        $\Prob(B_i)>0$ for all $i$, then
        \[
            \Prob(A)
            =
            \sum_{i=1}^{n}\Prob(A\mid B_i)\Prob(B_i).
        \]
    \end{theorem}

    \begin{theorem}[Law of Total Expectation; Tower Property]
        If $X$ is integrable and $Y$ is a random variable, then
        \[
            \Exp[X]=\Exp\!\left[\Exp[X\mid Y]\right].
        \]
        When $Y$ is discrete, this becomes
        \[
            \Exp[X]
            =
            \sum_y\Exp[X\mid Y=y]\Prob(Y=y),
        \]
        where the sum is taken over values $y$ with positive probability.
    \end{theorem}

    \begin{formula}[Law of Total Variance]
        If $X$ has finite second moment, then
        \[
            \Var(X)
            =
            \Exp\!\left[\Var(X\mid Y)\right]
            +
            \Var\!\left(\Exp[X\mid Y]\right).
        \]
    \end{formula}

    \begin{theorem}[Bayes' Theorem]
        If $B_1,\dots,B_n$ form a partition of the sample space and
        $\Prob(A)>0$, then
        \[
            \Prob(B_j\mid A)
            =
            \frac{\Prob(A\mid B_j)\Prob(B_j)}
                 {\sum_{i=1}^{n}\Prob(A\mid B_i)\Prob(B_i)}.
        \]
        In the two-event form,
        \[
            \Prob(B\mid A)
            =
            \frac{\Prob(A\mid B)\Prob(B)}{\Prob(A)}.
        \]
    \end{theorem}

    \begin{remark}
        Bayes' Theorem is the standard formula for reversing conditional
        probability. It updates a prior probability after observing evidence.
    \end{remark}

    \begin{remark}
        In general,
        \[
            \Prob(A\mid B)\neq\Prob(B\mid A).
        \]
        The two quantities reverse the roles of evidence and hypothesis.
        Bayes' Theorem relates them; it does not permit the conditioning bar to
        be reversed without the corresponding prior probabilities.
    \end{remark}

    \begin{definition}[Independence]
        Events $A$ and $B$ are independent if
        \[
            \Prob(A\cap B)=\Prob(A)\Prob(B).
        \]
        Equivalently, if $\Prob(B)>0$, then
        \[
            \Prob(A\mid B)=\Prob(A).
        \]
        Random variables $X$ and $Y$ are independent if events determined by
        $X$ and events determined by $Y$ are independent.
        A family of events is mutually independent if every finite subfamily
        has intersection probability equal to the product of its individual
        probabilities. In general,
        \[
            \text{pairwise independence}
            \nRightarrow
            \text{mutual independence}.
        \]
    \end{definition}

    \begin{tactic}
        ``Disjoint'' and ``independent'' are different. If two nonempty events
        are disjoint, then knowing that one occurred makes the other impossible;
        they are usually not independent.
    \end{tactic}

    \begin{formula}[Complement Rule and Inclusion--Exclusion for Probability]
        The complement rule is
        \[
            \Prob(A^c)=1-\Prob(A).
        \]
        For two events,
        \[
            \Prob(A\cup B)
            =
            \Prob(A)+\Prob(B)-\Prob(A\cap B).
        \]
        For three events,
        \[
            \begin{aligned}
                \Prob(A\cup B\cup C)
                &=
                \Prob(A)+\Prob(B)+\Prob(C)\\
                &\quad
                -\Prob(A\cap B)-\Prob(B\cap C)-\Prob(C\cap A)\\
                &\quad
                +\Prob(A\cap B\cap C).
            \end{aligned}
        \]
    \end{formula}

    \begin{proposition}[Linearity of Expectation]
        If $X_1,\dots,X_n$ are integrable random variables and
        $a_1,\dots,a_n$ are scalars, then
        \[
            \Exp\left[\sum_{i=1}^{n}a_iX_i\right]
            =
            \sum_{i=1}^{n}a_i\Exp[X_i].
        \]
        Independence is not required.
    \end{proposition}

    \begin{strategy}
        Linearity of expectation is often the most efficient method for counting
        an expected number of objects. Introduce indicator random variables and add their
        expectations.
    \end{strategy}

    \begin{proposition}[Expectation of a Product of Independent Random Variables]
        If $X_1,\dots,X_n$ are independent integrable random variables, then
        their product is integrable and
        \[
            \Exp[X_1\cdots X_n]
            =
            \Exp[X_1]\cdots\Exp[X_n].
        \]
        In particular, if they are also identically distributed, the right-hand
        side is $(\Exp[X_1])^n$.
    \end{proposition}

    \begin{formula}[Common Probability Distributions]
        The most frequently used distributions are summarized below. Unless
        otherwise indicated, $0<p<1$, $n,r\in\N$, $\lambda>0$,
        $\alpha,\beta>0$, $a<b$, $\sigma>0$, and $\nu\in\N$. For the
        negative binomial row, $X$ is the trial number on which the $r$th
        success occurs; the Gamma distribution uses the rate parameter
        $\lambda$.

        {\small
        \renewcommand{\arraystretch}{1.18}
        \begin{longtable}{@{}L{0.19\textwidth}L{0.35\textwidth}L{0.17\textwidth}L{0.17\textwidth}@{}}
            \toprule
            Distribution & Mass or density function & Mean & Variance \\
            \midrule
            \endfirsthead

            \toprule
            Distribution & Mass or density function & Mean & Variance \\
            \midrule
            \endhead

            \bottomrule
            \endlastfoot

            Bernoulli$(p)$ &
            $\Prob(X=x)=p^x(1-p)^{1-x}$, $x\in\{0,1\}$ &
            $p$ & $p(1-p)$ \\
            $\Bin(n,p)$ &
            $\displaystyle\Prob(X=k)=\binom{n}{k}p^k(1-p)^{n-k}$ &
            $np$ & $np(1-p)$ \\
            Geometric$(p)$ &
            $\Prob(X=k)=(1-p)^{k-1}p$, $k\geq1$ &
            $\displaystyle\frac1p$ & $\displaystyle\frac{1-p}{p^2}$ \\
            Negative binomial$(r,p)$ &
            $\displaystyle\Prob(X=k)=\binom{k-1}{r-1}p^r(1-p)^{k-r}$, $k\geq r$ &
            $\displaystyle\frac{r}{p}$ & $\displaystyle\frac{r(1-p)}{p^2}$ \\
            $\Poi(\lambda)$ &
            $\displaystyle\Prob(X=k)=e^{-\lambda}\frac{\lambda^k}{k!}$ &
            $\lambda$ & $\lambda$ \\
            Uniform$(a,b)$ &
            $\displaystyle f(x)=\frac{1}{b-a}$, $a\leq x\leq b$ &
            $\displaystyle\frac{a+b}{2}$ & $\displaystyle\frac{(b-a)^2}{12}$ \\
            Exponential$(\lambda)$ &
            $f(x)=\lambda e^{-\lambda x}$, $x\geq0$ &
            $\displaystyle\frac1\lambda$ & $\displaystyle\frac1{\lambda^2}$ \\
            Gamma$(\alpha,\lambda)$ &
            $\displaystyle f(x)=\frac{\lambda^\alpha}{\Gamma(\alpha)}x^{\alpha-1}e^{-\lambda x}$, $x>0$ &
            $\displaystyle\frac{\alpha}{\lambda}$ & $\displaystyle\frac{\alpha}{\lambda^2}$ \\
            Beta$(\alpha,\beta)$ &
            $\displaystyle f(x)=\frac{x^{\alpha-1}(1-x)^{\beta-1}}{B(\alpha,\beta)}$, $0<x<1$ &
            $\displaystyle\frac{\alpha}{\alpha+\beta}$ &
            {\scriptsize$\displaystyle\frac{\alpha\beta}{(\alpha+\beta)^2(\alpha+\beta+1)}$} \\
            $\chi^2(\nu)$ &
            $\displaystyle f(x)=\frac{x^{\nu/2-1}e^{-x/2}}{2^{\nu/2}\Gamma(\nu/2)}$, $x>0$ &
            $\nu$ & $2\nu$ \\
            $\Normal(\mu,\sigma^2)$ &
            $\displaystyle f(x)=\frac{1}{\sigma\sqrt{2\pi}}
            e^{-(x-\mu)^2/(2\sigma^2)}$ &
            $\mu$ & $\sigma^2$ \\
        \end{longtable}
        }
    \end{formula}

    \begin{formula}[Classical Sampling-Distribution Constructions]
        Let $Z,Z_1,\dots,Z_\nu\sim\Normal(0,1)$ be independent. Then
        \[
            \sum_{i=1}^{\nu}Z_i^2\sim\chi^2(\nu).
        \]
        If $U\sim\chi^2(\nu)$ is independent of $Z$, then
        \[
            \frac{Z}{\sqrt{U/\nu}}\sim t(\nu).
        \]
        If $U_1\sim\chi^2(d_1)$ and $U_2\sim\chi^2(d_2)$ are independent,
        then
        \[
            \frac{U_1/d_1}{U_2/d_2}\sim F(d_1,d_2).
        \]
        These three constructions are the standard recognition rules for the
        chi-square, Student $t$, and $F$ distributions.
    \end{formula}

    \begin{strategy}
        Waiting-until questions usually start with a geometric distribution. If
        a running sum is also involved, combine the expected number of trials
        with the expected contribution per trial, or condition on the first
        step.
    \end{strategy}

    \begin{formula}[Variance, Covariance, and Correlation]
        For random variables with finite second moments,
        \[
            \Var(X)=\Exp[X^2]-(\Exp[X])^2,
        \]
        \[
            \Cov(X,Y)
            =
            \Exp[XY]-\Exp[X]\Exp[Y],
        \]
        and
        \[
            \Var(X+Y)
            =
            \Var(X)+\Var(Y)+2\Cov(X,Y).
        \]
        If $\Var(X)>0$ and $\Var(Y)>0$, their correlation coefficient is
        \[
            \rho_{X,Y}
            =
            \frac{\Cov(X,Y)}
                 {\sqrt{\Var(X)\Var(Y)}},
            \qquad
            -1\leq\rho_{X,Y}\leq1.
        \]
        If $X$ and $Y$ are independent, then $\Cov(X,Y)=0$, so their variances
        add. In general,
        \[
            \Cov(X,Y)=0
            \nRightarrow
            X\text{ and }Y\text{ are independent}.
        \]
    \end{formula}

    \begin{formula}[Hypergeometric Distribution]
        Suppose a population has $N$ objects, of which $K$ are successes. If
        $n$ objects are chosen without replacement and $X$ is the number of
        successes chosen, then
        \[
            \Prob(X=k)
            =
            \frac{\binom{K}{k}\binom{N-K}{n-k}}{\binom{N}{n}}.
        \]
        Its mean and variance are
        \[
            \Exp[X]=n\frac{K}{N},
        \]
        \[
            \Var(X)
            =
            n\frac{K}{N}
            \left(1-\frac{K}{N}\right)
            \frac{N-n}{N-1}.
        \]
        The final factor is the finite-population correction caused by sampling
        without replacement.
    \end{formula}

    \begin{theorem}[Poisson Limit Theorem; Law of Small Numbers]
        Let $p_n\in[0,1]$ satisfy $np_n\to\lambda>0$. For each fixed
        nonnegative integer $k$,
        \[
            \binom{n}{k}p_n^k(1-p_n)^{n-k}
            \longrightarrow
            e^{-\lambda}\frac{\lambda^k}{k!}.
        \]
        Thus $\Bin(n,p_n)$ converges in distribution to
        $\Poi(\lambda)$.
    \end{theorem}

    \begin{remark}
        This theorem explains why the Poisson distribution models rare events:
        many trials, small success probability, and finite expected count.
    \end{remark}

    \begin{formula}[Empirical Rule; 68--95--99.7 Rule]
        For a normal distribution, approximately
        \[
            68\%
        \]
        of the mass lies within one standard deviation of the mean,
        \[
            95\%
        \]
        lies within two standard deviations, and
        \[
            99.7\%
        \]
        lies within three standard deviations.
    \end{formula}

    \begin{definition}[Modes of Convergence for Random Variables]
        Let $X_n$ and $X$ be random variables on the same probability space.
        The principal modes of convergence used in probability are:
        \begin{conditionlist}
            \item $X_n\asto X$ if
            \[
                \Prob\bigl(\{\omega\mid X_n(\omega)\to X(\omega)\}\bigr)=1;
            \]
            \item $X_n\probto X$ if
            \[
                \forall\varepsilon>0,
                \qquad
                \Prob\bigl(|X_n-X|>\varepsilon\bigr)\to 0;
            \]
            \item $X_n\Lpto{p}X$ for $p\geq 1$ if
            \[
                \Exp\bigl[|X_n-X|^p\bigr]\to 0;
            \]
            \item $X_n\distto X$ if
            \[
                F_{X_n}(x)\to F_X(x)
            \]
            at every continuity point $x$ of the distribution function $F_X$.
        \end{conditionlist}
        Almost-sure convergence implies convergence in probability, and
        convergence in probability implies convergence in distribution.
    \end{definition}

    \begin{theorem}[Borel--Cantelli Lemmas]
        If
        \[
            \sum_{n=1}^{\infty}\Prob(A_n)<\infty,
        \]
        then $\Prob(A_n\text{ infinitely often})=0$. If the $A_n$ are
        independent and the series diverges, then
        \[
            \Prob(A_n\text{ infinitely often})=1.
        \]
    \end{theorem}

    \begin{theorem}[Weak Law of Large Numbers]
        Let $X_1,X_2,\dots$ be independent and identically distributed random
        variables with finite mean $\mu$ and finite variance $\sigma^2$. Then
        \[
            \frac{X_1+\cdots+X_n}{n}
            \xrightarrow{\Prob}
            \mu.
        \]
        This follows directly from Chebyshev's inequality because the variance
        of the sample mean is $\sigma^2/n$.
    \end{theorem}

    \begin{idea}
        The Law of Large Numbers says that the sample average stabilizes near
        the true mean when the number of trials becomes large.
    \end{idea}

    \begin{theorem}[Central Limit Theorem]
        Let $X_1,X_2,\dots$ be independent identically distributed random
        variables with mean $\mu$ and variance $\sigma^2>0$. Then
        \[
            \frac{X_1+\cdots+X_n-n\mu}{\sigma\sqrt n}
            \distto
            Z,
            \qquad
            Z\sim\Normal(0,1).
        \]
        Thus the standardized sum converges in distribution to the standard
        normal distribution.
    \end{theorem}

    \begin{remark}
        The Central Limit Theorem explains why the normal distribution appears
        so often in nature and statistics.
    \end{remark}

    \begin{theorem}[Markov's Inequality]
        If $X$ is a nonnegative random variable and $a>0$, then
        \[
            \Prob(X\geq a)
            \leq
            \frac{\Exp[X]}{a}.
        \]
    \end{theorem}

    \begin{theorem}[Chebyshev's Inequality]
        If $X$ has finite mean $\mu$ and finite variance $\sigma^2>0$, then
        for every $k>0$,
        \[
            \Prob(|X-\mu|\geq k\sigma)
            \leq
            \frac{1}{k^2}.
        \]
        Equivalently, for every $\varepsilon>0$,
        \[
            \Prob(|X-\mu|\geq \varepsilon)
            \leq
            \frac{\Var(X)}{\varepsilon^2}.
        \]
    \end{theorem}

    \begin{remark}
        Markov's inequality uses only the mean, while Chebyshev's inequality
        uses the variance. Both are probability bounds that work without
        assuming a normal distribution.
    \end{remark}

    \begin{theorem}[Chernoff Bound for a Binomial Random Variable]
        If $X\sim\Bin(n,p)$ and $\mu=np$, then for $\delta>0$,
        \[
            \Prob\bigl(X\geq(1+\delta)\mu\bigr)
            \leq
            \left(\frac{e^\delta}{(1+\delta)^{1+\delta}}\right)^\mu.
        \]
        For $0<\delta<1$,
        \[
            \Prob\bigl(|X-\mu|\geq\delta\mu\bigr)
            \leq
            2e^{-\delta^2\mu/3}.
        \]
    \end{theorem}

    \begin{definition}[Convex and Concave Functions]
        Let $I\subseteq\R$ be an interval. A function $\varphi\colon I\to\R$ is
        convex if, for all $x,y\in I$ and $t\in[0,1]$,
        \[
            \varphi((1-t)x+ty)
            \leq
            (1-t)\varphi(x)+t\varphi(y).
        \]
        It is concave if the reverse inequality holds. Equivalently,
        $\varphi$ is concave if $-\varphi$ is convex.
    \end{definition}

    \begin{theorem}[Jensen's Inequality]
        Let $X$ be integrable and take values in an interval $I$, and let
        $\varphi\colon I\to\R$ be convex with $\varphi(X)$ integrable. Then
        \[
            \varphi(\Exp[X])
            \leq
            \Exp[\varphi(X)].
        \]
        If $\varphi$ is concave under the analogous integrability
        assumptions, the inequality is reversed.
    \end{theorem}

    \begin{idea}
        Jensen's inequality says that, for a convex function, applying the
        function after averaging gives a value no larger than averaging after
        applying the function.
    \end{idea}

    \begin{formula}[Exact Birthday Problem]
        Under the standard model of $365$ equally likely birthdays and no leap
        day, the probability that at least two people among $n$ share a birthday
        is
        \[
            1-
            \frac{365\cdot364\cdots(365-n+1)}{365^n}
            \qquad(0\leq n\leq365).
        \]
        For $n\geq366$, the probability is $1$ by the pigeonhole principle.
    \end{formula}

    \begin{formula}[Birthday Problem Approximation]
        In a group of $n$ people, assuming $365$ equally likely birthdays, the
        probability that at least two people share a birthday is approximately
        \[
            1-\exp\left(-\frac{n(n-1)}{730}\right).
        \]
        For $n=23$, this probability is already greater than $1/2$.
    \end{formula}

    \begin{remark}
        The birthday problem is famous because the answer is surprisingly
        small: only $23$ people are needed for a better-than-even chance of a
        shared birthday.
    \end{remark}

    \begin{formula}[Expected Number of Runs in Coin Tosses]
        In $n$ tosses of a fair coin, a run is a maximal consecutive block of
        identical outcomes. If $R_n$ is the number of runs, then
        \[
            \Exp[R_n]
            =
            1+\frac{n-1}{2}
            =
            \frac{n+1}{2}.
        \]
    \end{formula}

    \begin{idea}
        The first toss starts one run. For each of the remaining $n-1$ positions,
        a new run begins exactly when the current toss differs from the previous
        toss, which has probability $1/2$.
    \end{idea}

    \begin{formula}[Monty Hall Problem]
        There are three doors, exactly one of which hides a prize. The contestant
        first chooses a door. A host who knows the prize location must then open
        one of the other doors that hides no prize; whenever two such doors are
        available, the host chooses between them by a rule independent of the
        prize location. The contestant is then offered the unopened door.
        Under this protocol, staying wins with probability $1/3$, whereas
        switching wins with probability $2/3$.
    \end{formula}

    \begin{remark}
        The host's action is not random information: the host always reveals a
        losing door. This is why the original $2/3$ probability of having chosen
        incorrectly transfers to the remaining unopened door.
    \end{remark}

    \begin{formula}[Coupon Collector]
        If there are $n$ equally likely coupon types, the expected number of
        trials needed to see all $n$ types is
        \[
            nH_n
            =
            n\left(1+\frac12+\cdots+\frac1n\right).
        \]
    \end{formula}

    \begin{strategy}
        For a fair die, the expected number of rolls needed to see all six faces
        is $6H_6$. If a question asks about the order in which the first
        appearances occur, ignore repeated rolls and look only at the random
        ordering of the six faces.
    \end{strategy}

    \begin{formula}[Gambler's Ruin]
        Suppose a gambler starts with $i$ dollars and stops when reaching
        either $0$ dollars or $N$ dollars. If the probability of winning each
        round is $p$ and $q=1-p$, then the probability of reaching $N$ before
        $0$ is
        \[
            \frac{1-(q/p)^i}{1-(q/p)^N}
            \qquad
            (p\neq q).
        \]
        In the fair case $p=q=1/2$, the probability is
        \[
            \frac{i}{N}.
        \]
    \end{formula}

    \begin{remark}
        Gambler's ruin is a classic random walk result. It is a good example
        where the fair case has a very simple answer.
    \end{remark}

    \chapter{Geometry}

    \relatedcourses{}

    \relatedbooks{}

    \section*{Synthetic Geometry}

    \begin{fact}[Euclid's \textit{Elements} and the Axiomatic Method]
        Euclid's \textit{Elements} consists of thirteen books and organizes
        geometry and number theory through definitions, common notions,
        postulates, propositions, and proofs. Book I begins with five postulates;
        the fifth is the parallel postulate.
    \end{fact}

    \begin{fact}[Euclid's Fifth Postulate and Playfair's Axiom]
        Euclid's fifth postulate states that if a transversal makes the two
        interior angles on one side sum to less than two right angles, then the
        two lines meet on that side when extended. In the usual axiomatic setting
        for Euclidean plane geometry, this is equivalent to Playfair's axiom:
        through a point not on a given line, there is exactly one line parallel
        to the given line.
    \end{fact}

    \begin{definition}[Euclidean, Hyperbolic, and Elliptic Geometry]
        The three constant-curvature geometries are distinguished by their
        parallel behavior and triangle angle sums.
        \begin{longtable}{@{}L{0.22\textwidth}L{0.34\textwidth}L{0.30\textwidth}@{}}
            \toprule
            Geometry & Lines through a point off a given line & Triangle angle sum \\
            \midrule
            \endfirsthead

            \toprule
            Geometry & Lines through a point off a given line & Triangle angle sum \\
            \midrule
            \endhead

            \bottomrule
            \endlastfoot

            Euclidean & Exactly one parallel & Equal to $\pi$ \\
            Hyperbolic & Infinitely many disjoint lines & Less than $\pi$ \\
            Elliptic / spherical & No disjoint geodesic line & Greater than $\pi$ \\
        \end{longtable}
    \end{definition}

    \begin{theorem}[Parallel Postulate and Triangle Angle Sum]
        Within standard neutral geometry, the Euclidean parallel postulate is
        equivalent to the assertion that every triangle has angle sum
        \[
            \pi.
        \]
        Consequently, the familiar triangle angle-sum theorem is not a
        consequence of Euclid's first four postulates alone.
    \end{theorem}

    \begin{formula}[Geodesic Triangle Angle Sum on Constant Curvature]
        Let a geodesic triangle lie on a surface of constant Gaussian curvature
        $K$. If its angles are $\alpha,\beta,\gamma$, then the Gauss--Bonnet
        formula gives
        \[
            \alpha+\beta+\gamma-\pi
            =
            K\area(\triangle ABC).
        \]
        Thus positive curvature produces angle excess, zero curvature gives the
        Euclidean angle sum, and negative curvature produces angle defect.
    \end{formula}

    \begin{definition}[Convex Set, Convex Combination, and Convex Hull]
        A set $C\subseteq\R^n$ is convex if for every $\mathbf{x},\mathbf{y}\in C$
        and every $t\in[0,1]$,
        \[
            (1-t)\mathbf{x}+t\mathbf{y}\in C.
        \]
        More generally, a convex combination of points
        $\mathbf{x}_1,\dots,\mathbf{x}_m$ is a point
        \[
            \sum_{i=1}^{m}\lambda_i\mathbf{x}_i,
            \qquad
            \lambda_i\geq0,
            \qquad
            \sum_{i=1}^{m}\lambda_i=1.
        \]
        The convex hull $\conv(S)$ of a set $S\subseteq\R^n$ is the smallest
        convex set containing $S$; equivalently, it is the set of all finite
        convex combinations of points of $S$.
    \end{definition}

    \begin{formula}[Basic Plane and Solid Geometry Table]
        The basic measurement formulas are

        {\small
        \renewcommand{\arraystretch}{1.18}
        \begin{longtable}{@{}L{0.27\textwidth}L{0.28\textwidth}L{0.35\textwidth}@{}}
            \toprule
            Figure & Perimeter or surface area & Area or volume \\
            \midrule
            \endfirsthead

            \toprule
            Figure & Perimeter or surface area & Area or volume \\
            \midrule
            \endhead

            \bottomrule
            \endlastfoot

            Triangle & $a+b+c$ & $\displaystyle\frac12 bh$ \\
            Parallelogram & $2(a+b)$ & $bh$ \\
            Trapezoid & $a+b+c+d$ & $\displaystyle\frac12(a+c)h$ \\
            Regular $n$-gon & $na$ & $\displaystyle\frac12(na)r=\frac{na^2}{4\tan(\pi/n)}$ \\
            Circle & $2\pi r$ & $\pi r^2$ \\
            Circular sector & $r\theta$ & $\displaystyle\frac12r^2\theta$ \\
            Rectangular box & $2(ab+bc+ca)$ & $abc$ \\
            Right circular cylinder & $2\pi r(r+h)$ & $\pi r^2h$ \\
            Right circular cone & $\pi r(r+\ell)$ & $\displaystyle\frac13\pi r^2h$ \\
            Sphere & $4\pi r^2$ & $\displaystyle\frac43\pi r^3$ \\
        \end{longtable}
        }
        In the regular-polygon row, $r$ is the apothem; in the cone row,
        $\ell=\sqrt{r^2+h^2}$ is the slant height. Sector angles are measured in
        radians.
    \end{formula}

    \begin{theorem}[Triangle Angle Sum and Exterior Angle Theorem]
        For every triangle $\triangle ABC$,
        \[
            \angle A+\angle B+\angle C=\pi.
        \]
        An exterior angle equals the sum of the two nonadjacent interior angles.
    \end{theorem}

    \begin{proposition}[Triangle Congruence and Similarity Criteria]
        Two triangles are congruent under any of the criteria
        \[
            \mathrm{SSS},\qquad \mathrm{SAS},\qquad \mathrm{ASA},
            \qquad \mathrm{AAS},\qquad \mathrm{RHS}.
        \]
        They are similar under any of the criteria
        \[
            \mathrm{AA},\qquad \mathrm{SAS},\qquad \mathrm{SSS},
        \]
        where the latter two use proportional corresponding side lengths.
    \end{proposition}


    \begin{remark}
        For right triangles, RHS (right angle--hypotenuse--side), also called
        HL in some conventions, is a valid congruence criterion. A condition
        sometimes called RHA (right angle--hypotenuse--acute angle) is also
        sufficient, but it is already an ASA/AAS case because the two acute
        angles of a right triangle are complementary. By contrast, SSA is not
        a congruence criterion in general; the right-triangle hypotheses are
        essential in the familiar hypotenuse-based special cases.
    \end{remark}

    \begin{theorem}[Midpoint Theorem and Centroid Ratio]
        The segment joining the midpoints of two sides of a triangle is parallel
        to the third side and has half its length. The three medians are
        concurrent at the centroid $G$, and $G$ divides every median in the ratio
        \[
            2:1
        \]
        measured from the vertex.
    \end{theorem}

    \begin{strategy}[Reverse Engineering in Geometric Constructions]
        First draw a completed configuration satisfying the desired conditions.
        Then identify which points, perpendiculars, parallels, angle bisectors,
        circles, or reflections in the completed picture can be constructed from
        the original data. Familiar configurations involving triangle centers,
        similar triangles, and power of a point are especially useful targets.
    \end{strategy}


    \begin{fact}[Huzita--Hatori Origami Axioms]
        The standard single-fold model of mathematical origami is described by
        seven folding axioms commonly called the Huzita--Hatori axioms. They
        specify folds determined by prescribed incidences among points and
        lines. Unlike straightedge-and-compass constructions, single-fold
        origami can solve certain cubic equations; in particular, it can carry
        out classical constructions such as general angle trisection and
        doubling the cube.
    \end{fact}

    \begin{definition}[Classical Triangle Centers]
        For a triangle $\triangle ABC$, the most important classical centers
        are:
        \[
            O=\text{circumcenter},
            \qquad
            I=\text{incenter},
            \qquad
            H=\text{orthocenter},
            \qquad
            G=\text{centroid}.
        \]
        Here $O$ is the center of the circumcircle, $I$ is the center of the
        incircle, $H$ is the intersection of the altitudes, and $G$ is the
        intersection of the medians.
    \end{definition}

    \begin{definition}[Incircle, Excircles, Inradius, and Exradii]
        The incircle of a triangle is the circle tangent to all three sides; its
        center is the incenter $I$, and its radius is the inradius $r$. An
        excircle is tangent to one side of the triangle and to the extensions of
        the other two sides. The corresponding excenters and exradii are denoted
        by $I_a,I_b,I_c$ and $r_a,r_b,r_c$, respectively, according to the
        opposite vertex.
    \end{definition}

    \begin{theorem}[Gergonne Point and Nagel Point]
        The cevians joining the vertices of a triangle to the contact points of
        its incircle with the opposite sides are concurrent at the Gergonne point.
        The cevians joining the vertices to the contact points of the opposite
        excircles are concurrent at the Nagel point.
    \end{theorem}

    \begin{definition}[Isogonal Lines and Symmedians]
        Two lines through a vertex of a triangle are isogonal if they are
        reflections of one another across the internal angle bisector at that
        vertex. The $A$-symmedian of $\triangle ABC$ is the line isogonal to the
        $A$-median; the other symmedians are defined cyclically.
    \end{definition}

    \begin{theorem}[Symmedian Lemma and Lemoine Point]
        Let the $A$-symmedian of $\triangle ABC$ meet $BC$ at $K$. Then
        \[
            \frac{BK}{KC}=\frac{AB^2}{AC^2}.
        \]
        The three symmedians are concurrent at the Lemoine point, also called
        the symmedian point.
    \end{theorem}

    \begin{theorem}[Fermat Point; Torricelli Point]
        If every angle of $\triangle ABC$ is less than $120^\circ$, there is a
        unique point $P$ minimizing
        \[
            PA+PB+PC.
        \]
        It satisfies
        \[
            \angle APB=\angle BPC=\angle CPA=120^\circ.
        \]
        If one angle is at least $120^\circ$, the minimizing point is that vertex.
    \end{theorem}

    \begin{formula}[Extended Law of Sines]
        For a triangle $\triangle ABC$ with side lengths
        \[
            a=BC,\qquad b=CA,\qquad c=AB,
        \]
        and circumradius $R$, we have
        \[
            \frac{a}{\sin A}
            =
            \frac{b}{\sin B}
            =
            \frac{c}{\sin C}
            =
            2R.
        \]
    \end{formula}

    \begin{formula}[Law of Cosines]
        For a triangle $\triangle ABC$ with side lengths
        \[
            a=BC,\qquad b=CA,\qquad c=AB,
        \]
        we have
        \[
            c^2=a^2+b^2-2ab\cos C.
        \]
        The analogous formulas hold after cyclic permutation of $a,b,c$.
    \end{formula}

    \begin{formula}[Heron's Formula]
        Let $\triangle ABC$ have side lengths $a,b,c$ and semiperimeter
        \[
            s=\frac{a+b+c}{2}.
        \]
        Then its area $[ABC]$ is
        \[
            [ABC]
            =
            \sqrt{s(s-a)(s-b)(s-c)}.
        \]
        Also,
        \[
            [ABC]=rs
            \qquad\text{and}\qquad
            [ABC]=\frac{abc}{4R},
        \]
        where $r$ is the inradius and $R$ is the circumradius.
    \end{formula}

    \begin{strategy}
        The formulas
        \[
            [ABC]=rs,
            \qquad
            [ABC]=\frac{abc}{4R}
        \]
        are often more useful than Heron's formula when incircles or
        circumcircles appear.
    \end{strategy}

    \begin{formula}[Triangle Area Formulas]
        For a triangle with side lengths $a,b,c$, angles $A,B,C$, circumradius
        $R$, and inradius $r$,
        \[
            [ABC]
            =
            \frac12bc\sin A
            =
            \frac12ca\sin B
            =
            \frac12ab\sin C,
        \]
        and
        \[
            [ABC]=rs=\frac{abc}{4R}.
        \]
    \end{formula}

    \begin{theorem}[Weitzenb\"ock's Inequality]
        For a triangle with side lengths $a,b,c$ and area $[ABC]$,
        \[
            a^2+b^2+c^2\geq4\sqrt{3}\,[ABC],
        \]
        with equality if and only if the triangle is equilateral.
    \end{theorem}

    \begin{formula}[Inradius and Exradii]
        For a triangle $\triangle ABC$ with semiperimeter $s$, the inradius is
        \[
            r=\frac{[ABC]}{s}.
        \]
        If $r_a,r_b,r_c$ are the exradii opposite $A,B,C$, respectively, then
        \[
            r_a=\frac{[ABC]}{s-a},
            \qquad
            r_b=\frac{[ABC]}{s-b},
            \qquad
            r_c=\frac{[ABC]}{s-c}.
        \]
    \end{formula}

    \begin{formula}[Median Length Formula; Apollonius' Theorem]
        If $m_a$ is the median from $A$ to side $a=BC$, then
        \[
            m_a^2
            =
            \frac{2b^2+2c^2-a^2}{4}.
        \]
        Equivalently, if $M$ is the midpoint of $BC$, then
        \[
            AB^2+AC^2=2(AM^2+BM^2).
        \]
    \end{formula}

    \begin{theorem}[Stewart's Theorem]
        Let $D$ lie on side $BC$ of a triangle $\triangle ABC$. Write
        \[
            a=BC,
            \qquad
            b=CA,
            \qquad
            c=AB,
            \qquad
            m=BD,
            \qquad
            n=DC,
            \qquad
            d=AD,
        \]
        so that $a=m+n$. Then
        \[
            b^2m+c^2n
            =
            a(d^2+mn).
        \]
        Equivalently,
        \[
            d^2
            =
            \frac{b^2m+c^2n}{a}-mn.
        \]
    \end{theorem}

    \begin{remark}
        Stewart's Theorem is the standard general formula for the length of a
        cevian. The median-length formula is the special case $m=n=a/2$.
    \end{remark}

    \begin{formula}[Internal Angle-Bisector Length]
        Let the internal angle bisector from $A$ meet $BC$ at $D$, and denote
        its length by $\ell_a=AD$. For
        \[
            a=BC,
            \qquad
            b=CA,
            \qquad
            c=AB,
        \]
        we have
        \[
            \ell_a^2
            =
            bc\left(1-\frac{a^2}{(b+c)^2}\right)
            =
            \frac{4bc\,s(s-a)}{(b+c)^2},
        \]
        where $s=(a+b+c)/2$.
    \end{formula}

    \begin{theorem}[Angle Bisector Theorem]
        In a triangle $\triangle ABC$, suppose the internal angle bisector of
        $\angle A$ meets $BC$ at $D$. Then
        \[
            \frac{BD}{DC}
            =
            \frac{AB}{AC}.
        \]
    \end{theorem}

    \begin{theorem}[Ceva's Theorem]
        Let $D,E,F$ lie on the lines $BC,CA,AB$, respectively. Then the cevians
        $AD,BE,CF$ are concurrent if and only if
        \[
            \frac{BD}{DC}
            \cdot
            \frac{CE}{EA}
            \cdot
            \frac{AF}{FB}
            =
            1,
        \]
        using directed lengths.
    \end{theorem}

    \begin{remark}
        Ceva's Theorem is the standard theorem for proving concurrence of three
        lines in a triangle.
    \end{remark}

    \begin{theorem}[Trigonometric Form of Ceva's Theorem]
        Let $D,E,F$ lie on the lines $BC,CA,AB$, respectively. The cevians
        $AD,BE,CF$ are concurrent if and only if
        \[
            \frac{\sin\angle BAD}{\sin\angle DAC}
            \cdot
            \frac{\sin\angle CBE}{\sin\angle EBA}
            \cdot
            \frac{\sin\angle ACF}{\sin\angle FCB}
            =1,
        \]
        with the usual directed-angle convention when points lie on extended
        sides.
    \end{theorem}

    \begin{strategy}[Mass Points]
        In a cevian configuration, assign masses to vertices so that a point
        dividing a segment balances the endpoint masses inversely to the segment
        lengths. For example, if $D\in BC$, then
        \[
            \frac{BD}{DC}=\frac{m_C}{m_B}.
        \]
        Consistent masses convert several ratio computations into weighted
        averages and often recover Ceva-type relations immediately.
    \end{strategy}

    \begin{theorem}[Menelaus' Theorem]
        Let $D,E,F$ lie on the lines $BC,CA,AB$, respectively. Then the points
        $D,E,F$ are collinear if and only if
        \[
            \frac{BD}{DC}
            \cdot
            \frac{CE}{EA}
            \cdot
            \frac{AF}{FB}
            =
            -1,
        \]
        using directed lengths.
    \end{theorem}

    \begin{remark}
        Ceva is for concurrence, while Menelaus is for collinearity. This pair
        is one of the most important toolkits in olympiad geometry.
    \end{remark}

    \begin{theorem}[Euler Line Theorem]
        In any non-equilateral triangle, the circumcenter $O$, centroid $G$,
        and orthocenter $H$ are collinear. Moreover,
        \[
            OG:GH=1:2.
        \]
        The line passing through these points is called the Euler line.
    \end{theorem}

    \begin{theorem}[Nine-Point Circle Theorem; Euler Circle]
        In any triangle, the following nine points lie on one circle:
        the three side midpoints, the three feet of the altitudes, and the
        three midpoints between the orthocenter and the vertices. This circle
        is called the nine-point circle.
    \end{theorem}

    \begin{remark}
        The center of the nine-point circle is the midpoint of the segment
        joining the circumcenter $O$ and the orthocenter $H$.
    \end{remark}

    \begin{theorem}[Feuerbach's Theorem]
        The nine-point circle of a triangle is internally tangent to the incircle
        and externally tangent to each of the three excircles.
    \end{theorem}

    \begin{theorem}[Euler's Formula for the Incenter and Circumcenter]
        For a triangle with circumradius $R$, inradius $r$, circumcenter $O$,
        and incenter $I$, we have
        \[
            OI^2=R^2-2Rr.
        \]
        In particular,
        \[
            R\geq 2r.
        \]
    \end{theorem}

    \begin{remark}
        The inequality $R\geq 2r$ is called Euler's inequality for triangles.
        Equality holds exactly for an equilateral triangle.
    \end{remark}

    \begin{strategy}[Spiral Similarity]
        When two segments must be compared through both an angle and a ratio,
        look for a center of spiral similarity. It often reveals the required
        similar triangles immediately.
    \end{strategy}

    \begin{definition}[Directed Angles Modulo $\pi$]
        When directed angles modulo $\pi$ are used, $\angle ABC$ denotes the
        directed angle from the line $BA$ to the line $BC$, taken modulo $\pi$.
        Reversing the orientation of either line changes an ordinary oriented
        angle by $\pi$ and therefore does not change its class modulo $\pi$.
        This convention avoids separate case distinctions for many collinearity
        and cyclicity arguments.
    \end{definition}

    \begin{theorem}[Inscribed Angle Theorem]
        An inscribed angle equals half the central angle subtending the same arc.
        Consequently, angles subtending the same chord are equal.
    \end{theorem}

    \begin{theorem}[Cyclic Quadrilateral Criteria]
        Let $A,B,C,D$ be four distinct points with no three collinear. Using
        directed angles modulo $\pi$, the points are concyclic if and only if
        \[
            \angle BAC\equiv\angle BDC\pmod\pi.
        \]
        If $ABCD$ is a convex quadrilateral in cyclic order, this is equivalent
        to
        \[
            \angle ABC+\angle ADC=\pi.
        \]
    \end{theorem}

    \begin{theorem}[Tangent--Chord Theorem; Alternate Segment Theorem]
        The angle between a tangent to a circle and a chord through the point of
        tangency equals the inscribed angle subtending that chord in the opposite
        arc.
    \end{theorem}

    \begin{theorem}[Power of a Point Theorem]
        Let $P$ be a point and let a line through $P$ meet a circle at
        $A$ and $B$. The product
        \[
            PA\cdot PB
        \]
        is independent of the choice of the line through $P$. It is called the
        power of $P$ with respect to the circle.
    \end{theorem}

    \begin{formula}[Tangent--Secant Theorem]
        If $PT$ is tangent to a circle at $T$ and a secant through $P$ meets
        the circle at $A$ and $B$, then
        \[
            PT^2=PA\cdot PB.
        \]
    \end{formula}

    \begin{theorem}[Radical Axis Theorem]
        For two nonconcentric circles, the set of points having equal powers
        with respect to the two circles is a line. This line is called the
        radical axis of the two circles.
    \end{theorem}

    \begin{remark}
        Radical axes are powerful because they turn circle configurations into
        collinearity statements.
    \end{remark}

    \begin{theorem}[Ptolemy's Theorem]
        Let $A,B,C,D$ be the vertices of a convex cyclic quadrilateral in this
        cyclic order. Then
        \[
            AC\cdot BD
            =
            AB\cdot CD+BC\cdot DA.
        \]
        Conversely, equality in Ptolemy's inequality for a convex
        quadrilateral characterizes cyclicity.
    \end{theorem}

    \begin{theorem}[Ptolemy's Inequality]
        For any four points $A,B,C,D$ in the Euclidean plane,
        \[
            AC\cdot BD
            \leq
            AB\cdot CD+BC\cdot DA.
        \]
        Equality holds in the cyclic convex case and in the corresponding
        degenerate collinear case.
    \end{theorem}

    \begin{formula}[Brahmagupta's Formula]
        If a cyclic quadrilateral has side lengths $a,b,c,d$ and semiperimeter
        \[
            s=\frac{a+b+c+d}{2},
        \]
        then its area is
        \[
            \area(ABCD)
            =
            \sqrt{(s-a)(s-b)(s-c)(s-d)}.
        \]
        Heron's formula is the limiting triangular case.
    \end{formula}

    \begin{theorem}[Varignon's Theorem]
        The midpoints of the four sides of any quadrilateral form a
        parallelogram. Its area is half the area of the original quadrilateral.
    \end{theorem}

    \begin{theorem}[Pitot's Theorem]
        If a convex quadrilateral with consecutive side lengths $a,b,c,d$ has an
        incircle tangent to all four sides, then
        \[
            a+c=b+d.
        \]
    \end{theorem}

    \begin{theorem}[Euler's Quadrilateral Theorem]
        If $M$ and $N$ are the midpoints of the diagonals of a quadrilateral with
        side lengths $a,b,c,d$ and diagonal lengths $p,q$, then
        \[
            a^2+b^2+c^2+d^2=p^2+q^2+4MN^2.
        \]
        The parallelogram law is the case $M=N$.
    \end{theorem}

    \begin{theorem}[Simson Line Theorem]
        Let $P$ be a point on the circumcircle of $\triangle ABC$. The feet of
        the perpendiculars from $P$ to the three lines $AB,BC,CA$ are collinear.
        This line is called the Simson line of $P$ with respect to
        $\triangle ABC$.
    \end{theorem}

    \begin{theorem}[Miquel's Theorem]
        In a complete quadrilateral, the circumcircles of the four triangles
        formed by choosing three of the four lines pass through a common point.
        This point is called the Miquel point.
    \end{theorem}




    \begin{theorem}[British Flag Theorem]
        For any point $P$ in the plane of a rectangle $ABCD$,
        \[
            PA^2+PC^2=PB^2+PD^2.
        \]
    \end{theorem}

    \begin{theorem}[Van Aubel's Theorem]
        Construct squares externally on the four sides of a quadrilateral. The
        segments joining the centers of opposite squares are equal in length
        and perpendicular.
    \end{theorem}

    \begin{theorem}[Viviani's Theorem]
        For any point $P$ inside an equilateral triangle of altitude $h$, the sum
        of the perpendicular distances from $P$ to the three sides is $h$.
    \end{theorem}

    \begin{theorem}[Desargues' Theorem]
        If two triangles are perspective from a point, then the three
        intersections of corresponding sides are collinear. Conversely, in the
        projective plane, collinearity of those intersections implies concurrence
        of the three lines joining corresponding vertices.
    \end{theorem}

    \begin{theorem}[Pappus' Hexagon Theorem]
        Let $A,B,C$ lie on one line and $A',B',C'$ lie on another. Then the three
        points
        \[
            AB'\cap A'B,
            \qquad
            AC'\cap A'C,
            \qquad
            BC'\cap B'C
        \]
        are collinear.
    \end{theorem}

    \begin{theorem}[Homothety Centers of Two Circles]
        For two nonconcentric circles, the two external common tangents meet at
        the external homothety center, and the two internal common tangents meet
        at the internal homothety center, when these intersections exist. Each
        homothety center lies on the line joining the two circle centers.
    \end{theorem}

    \begin{formula}[Circle Inversion]
        Under inversion with center $O$ and radius $\rho$, a point $P\neq O$
        is sent to the point $P'$ on the ray $OP$ satisfying
        \[
            OP\cdot OP'=\rho^2.
        \]
        Inversion preserves angles in magnitude, sends circles through $O$ to
        lines not through $O$, and sends such lines back to circles through $O$.
    \end{formula}

    \begin{formula}[Regular Polygon Tiling Angle Condition]
        If $q$ regular $p$-gons meet edge-to-edge at each vertex of a Euclidean
        tiling, then
        \[
            q\left(1-\frac{2}{p}\right)\pi=2\pi,
        \]
        equivalently,
        \[
            (p-2)(q-2)=4.
        \]
        Thus the only tilings by one type of regular polygon are
        \[
            (p,q)=(3,6),(4,4),(6,3).
        \]
    \end{formula}

    \begin{theorem}[Classification of Platonic Solids]
        There are exactly five Platonic solids:
        \[
            \text{tetrahedron},\quad
            \text{cube},\quad
            \text{octahedron},\quad
            \text{dodecahedron},\quad
            \text{icosahedron}.
        \]
    \end{theorem}

    \section*{Analytic Geometry}

    \begin{formula}[Coordinate Geometry Table]
        For points $A=(x_1,y_1)$ and $B=(x_2,y_2)$,

        {\small
        \renewcommand{\arraystretch}{1.18}
        \begin{longtable}{@{}L{0.27\textwidth}L{0.63\textwidth}@{}}
            \toprule
            Quantity & Formula \\
            \midrule
            \endfirsthead

            \toprule
            Quantity & Formula \\
            \midrule
            \endhead

            \bottomrule
            \endlastfoot

            Distance & $\displaystyle AB=\sqrt{(x_2-x_1)^2+(y_2-y_1)^2}$ \\
            Midpoint & $\displaystyle\left(\frac{x_1+x_2}{2},\frac{y_1+y_2}{2}\right)$ \\
            Internal division $AP:PB=m:n$ &
            $\displaystyle\left(\frac{nx_1+mx_2}{m+n},
            \frac{ny_1+my_2}{m+n}\right)$ \\
            Slope & $\displaystyle\frac{y_2-y_1}{x_2-x_1}$, when $x_1\neq x_2$ \\
            Point-slope line & $y-y_1=m(x-x_1)$ \\
            General line & $ax+by+c=0$ \\
            Distance from $(x_0,y_0)$ to $ax+by+c=0$ &
            $\displaystyle\frac{|ax_0+by_0+c|}{\sqrt{a^2+b^2}}$ \\
        \end{longtable}
        }
    \end{formula}

    \begin{formula}[Triangle Centers in Coordinates]
        For $A=(x_1,y_1)$, $B=(x_2,y_2)$, and $C=(x_3,y_3)$, the centroid is
        \[
            G
            =
            \left(
                \frac{x_1+x_2+x_3}{3},
                \frac{y_1+y_2+y_3}{3}
            \right).
        \]
        If the side lengths opposite $A,B,C$ are $a,b,c$, respectively, then the
        incenter is
        \[
            I
            =
            \frac{aA+bB+cC}{a+b+c}.
        \]
    \end{formula}

    \begin{formula}[Barycentric Coordinates]
        Let $A,B,C$ be noncollinear. Every point $P$ in their affine plane has
        unique coordinates $(\alpha,\beta,\gamma)$ satisfying
        \[
            \alpha+\beta+\gamma=1,
            \qquad
            \mathbf{p}
            =
            \alpha\mathbf{a}+\beta\mathbf{b}+\gamma\mathbf{c}.
        \]
        For $P$ inside $\triangle ABC$,
        \[
            \alpha=\frac{[PBC]}{[ABC]},
            \qquad
            \beta=\frac{[PCA]}{[ABC]},
            \qquad
            \gamma=\frac{[PAB]}{[ABC]}.
        \]
    \end{formula}

    \begin{formula}[Shoelace Formula; Gauss's Area Formula]
        For a polygon with vertices
        \[
            (x_1,y_1),(x_2,y_2),\dots,(x_n,y_n)
        \]
        in cyclic order, its area is
        \[
            \frac12
            \left|
            \sum_{i=1}^{n}
            (x_i y_{i+1}-x_{i+1}y_i)
            \right|,
        \]
        where
        \[
            (x_{n+1},y_{n+1})=(x_1,y_1).
        \]
    \end{formula}

    \begin{remark}
        The Shoelace Formula is one of the cleanest examples of analytic
        geometry: a geometric area becomes a determinant-like coordinate
        calculation.
    \end{remark}

    \begin{formula}[Oriented Area and Collinearity Test]
        For points $A=(x_1,y_1)$, $B=(x_2,y_2)$, and $C=(x_3,y_3)$, the
        oriented double area is
        \[
            2[ABC]_{\mathrm{or}}
            =
            \det
            \begin{pmatrix}
                x_1 & y_1 & 1 \\
                x_2 & y_2 & 1 \\
                x_3 & y_3 & 1
            \end{pmatrix}.
        \]
        Hence $A,B,C$ are collinear if and only if this determinant is $0$.
    \end{formula}

    \begin{formula}[Plane Rotation]
        Counterclockwise rotation through angle $\theta$ about the origin sends
        \[
            \begin{pmatrix}x\\y\end{pmatrix}
            \longmapsto
            \begin{pmatrix}
                \cos\theta & -\sin\theta \\
                \sin\theta & \cos\theta
            \end{pmatrix}
            \begin{pmatrix}x\\y\end{pmatrix}.
        \]
        In complex notation, the same transformation is simply
        \[
            z\longmapsto e^{i\theta}z.
        \]
    \end{formula}

    \begin{formula}[Reflection across a Line through the Origin]
        If a line through the origin makes angle $\theta$ with the positive
        $x$-axis, reflection across that line is
        \[
            z\longmapsto e^{2i\theta}\overline{z}
        \]
        in complex notation.
    \end{formula}

    \begin{strategy}[Affine Normalization]
        Affine transformations preserve collinearity, parallelism, ratios along
        one line, concurrence, and ratios of areas. Therefore, when a problem
        depends only on these properties, one may send a triangle to convenient
        coordinates such as
        \[
            A=(0,0),
            \qquad
            B=(1,0),
            \qquad
            C=(0,1).
        \]
        Angles, lengths, and circles are not generally preserved.
    \end{strategy}

    \begin{theorem}[Cavalieri's Principle]
        If two plane regions lie between the same pair of parallel lines and
        every line parallel to them cuts the regions in segments whose lengths
        have a constant ratio $c$, then their areas have ratio $c$. The
        three-dimensional analogue replaces segment lengths by cross-sectional
        areas and concludes the same ratio for volumes.
    \end{theorem}

    \begin{strategy}[Unfolding Method for Reflection Problems]
        A broken path reflecting from a line or face becomes a straight segment
        after reflecting the region across each boundary in succession. Thus a
        shortest reflected path is found by measuring a straight-line distance in
        the unfolded picture and then folding the path back.
    \end{strategy}

    \begin{formula}[Lattice Segment and Cell Counts]
        A segment from $(0,0)$ to $(m,n)$, where $m,n\in\N$, passes through
        \[
            m+n-\gcd(m,n)
        \]
        unit squares. The number of lattice points on the segment, including its
        endpoints, is
        \[
            \gcd(m,n)+1.
        \]
        In three dimensions, a segment from $(0,0,0)$ to $(a,b,c)$ passes through
        \[
            a+b+c
            -\gcd(a,b)-\gcd(a,c)-\gcd(b,c)
            +\gcd(a,b,c)
        \]
        unit cubes.
    \end{formula}

    \begin{formula}[Vector Test for Parallel Segments]
        Directed segments $AB$ and $CD$ are parallel if and only if
        \[
            \mathbf{b}-\mathbf{a}
            =
            \lambda(\mathbf{d}-\mathbf{c})
        \]
        for some scalar $\lambda$. Equality with $\lambda=1$ identifies
        equal directed segments and is often the most efficient method to detect a hidden
        parallelogram.
    \end{formula}

    \begin{formula}[Vector Projection]
        If $\mathbf{u}\neq\mathbf{0}$, then the projection of $\mathbf{v}$ onto
        the line spanned by $\mathbf{u}$ is
        \[
            \proj_{\mathbf{u}}\mathbf{v}
            =
            \frac{\mathbf{v}\cdot\mathbf{u}}{\mathbf{u}\cdot\mathbf{u}}
            \mathbf{u}.
        \]
        The scalar component of $\mathbf{v}$ in the direction of $\mathbf{u}$ is
        \[
            \frac{\mathbf{v}\cdot\mathbf{u}}{\|\mathbf{u}\|}.
        \]
    \end{formula}

    \begin{formula}[Vector Equations of Lines and Planes]
        A line through a point $\mathbf{p}$ with direction vector
        $\mathbf{v}\neq\mathbf{0}$ is
        \[
            \mathbf{x}=\mathbf{p}+t\mathbf{v}
            \qquad (t\in\R).
        \]
        A plane through a point $\mathbf{p}$ with normal vector
        $\mathbf{n}\neq\mathbf{0}$ is
        \[
            (\mathbf{x}-\mathbf{p})\cdot\mathbf{n}=0.
        \]
    \end{formula}


    \begin{formula}[Angles and Volumes in $\R^3$]
        If nonzero vectors $\mathbf{u},\mathbf{v}$ are directions of two lines,
        their acute angle $\theta$ satisfies
        \[
            \cos\theta
            =
            \frac{|\mathbf{u}\cdot\mathbf{v}|}
                 {\|\mathbf{u}\|\,\|\mathbf{v}\|}.
        \]
        If $\mathbf{n}_1,\mathbf{n}_2$ are normals to two planes, the acute
        angle between the planes is given by the same formula with the normals.
        If $\mathbf{v}$ is a line direction and $\mathbf{n}$ is a plane normal,
        the acute line--plane angle $\alpha$ satisfies
        \[
            \sin\alpha
            =
            \frac{|\mathbf{v}\cdot\mathbf{n}|}
                 {\|\mathbf{v}\|\,\|\mathbf{n}\|}.
        \]
        The volume of the parallelepiped spanned by
        $\mathbf{a},\mathbf{b},\mathbf{c}$ is
        \[
            \bigl|\mathbf{a}\cdot(\mathbf{b}\times\mathbf{c})\bigr|,
        \]
        and the tetrahedron with the same three edge vectors from one vertex
        has one sixth of this volume.
    \end{formula}

    \begin{formula}[Distance between Skew Lines]
        Let
        \[
            L_1:\mathbf{x}=\mathbf{p}_1+t\mathbf{u},
            \qquad
            L_2:\mathbf{x}=\mathbf{p}_2+s\mathbf{v},
        \]
        with $\mathbf{u}\times\mathbf{v}\neq\mathbf{0}$. Then
        \[
            \dist(L_1,L_2)
            =
            \frac{|(\mathbf{p}_2-\mathbf{p}_1)\cdot
                (\mathbf{u}\times\mathbf{v})|}
                 {\|\mathbf{u}\times\mathbf{v}\|}.
        \]
    \end{formula}

    \begin{formula}[Distance from a Point to a Line]
        In $\R^3$, the distance from a point $\mathbf{p}_0$ to the line
        \[
            \mathbf{x}=\mathbf{p}+t\mathbf{v}
        \]
        is
        \[
            \frac{\|(\mathbf{p}_0-\mathbf{p})\times\mathbf{v}\|}{\|\mathbf{v}\|}.
        \]
        In $\R^2$, the distance from $(x_0,y_0)$ to the line $ax+by+c=0$ is
        \[
            \frac{|ax_0+by_0+c|}{\sqrt{a^2+b^2}}.
        \]
    \end{formula}

    \begin{formula}[Distance from a Point to a Plane]
        The distance from a point $\mathbf{p}_0=(x_0,y_0,z_0)$ to the plane
        \[
            ax+by+cz+d=0
        \]
        is
        \[
            \frac{|ax_0+by_0+cz_0+d|}{\sqrt{a^2+b^2+c^2}}.
        \]
    \end{formula}

    \begin{formula}[Sphere Equation]
        A sphere with center $\mathbf{c}=(a,b,c)$ and radius $r$ has equation
        \[
            (x-a)^2+(y-b)^2+(z-c)^2=r^2.
        \]
        Equivalently,
        \[
            \|\mathbf{x}-\mathbf{c}\|=r.
        \]
    \end{formula}


    \begin{formula}[Standard Forms of Common Quadrics in $\R^3$]
        With $a,b,c>0$, the principal-axis standard forms include:

        {\small
        \renewcommand{\arraystretch}{1.16}
        \begin{longtable}{@{}L{0.23\textwidth}L{0.67\textwidth}@{}}
            \toprule
            Surface & Standard form \\
            \midrule
            \endfirsthead

            \toprule
            Surface & Standard form \\
            \midrule
            \endhead

            \bottomrule
            \endlastfoot

            Ellipsoid &
            $\displaystyle \frac{x^2}{a^2}+\frac{y^2}{b^2}+\frac{z^2}{c^2}=1$ \\
            Elliptic paraboloid &
            $\displaystyle z=\frac{x^2}{a^2}+\frac{y^2}{b^2}$ \\
            Hyperbolic paraboloid &
            $\displaystyle z=\frac{x^2}{a^2}-\frac{y^2}{b^2}$ \\
            One-sheet hyperboloid &
            $\displaystyle \frac{x^2}{a^2}+\frac{y^2}{b^2}-\frac{z^2}{c^2}=1$ \\
            Two-sheet hyperboloid &
            $\displaystyle -\frac{x^2}{a^2}-\frac{y^2}{b^2}+\frac{z^2}{c^2}=1$ \\
            Elliptic cone &
            $\displaystyle \frac{x^2}{a^2}+\frac{y^2}{b^2}-\frac{z^2}{c^2}=0$ \\
            Elliptic cylinder &
            $\displaystyle \frac{x^2}{a^2}+\frac{y^2}{b^2}=1$ \\
        \end{longtable}
        }
        More generally, a quadric can be written
        \[
            \mathbf{x}^{\mathsf T}\mathbf{A}\mathbf{x}
            +\mathbf{b}^{\mathsf T}\mathbf{x}+c_0=0,
        \]
        with $\mathbf{A}$ real symmetric. Translation and orthogonal
        diagonalization of $\mathbf{A}$ reduce the quadratic part to principal
        axes, linking the geometry directly to the spectral theorem.
    \end{formula}

    \begin{strategy}[Line--Circle Intersection by a Discriminant]
        Substitute a parametric line into a circle equation. The resulting
        quadratic equation has two, one, or no real solutions according as its
        discriminant is positive, zero, or negative. The zero-discriminant case
        is the tangency condition.
    \end{strategy}

    \begin{formula}[Circle Equation]
        The circle with center $(a,b)$ and radius $r$ has equation
        \[
            (x-a)^2+(y-b)^2=r^2.
        \]
        The general equation
        \[
            x^2+y^2+Dx+Ey+F=0
        \]
        can be converted to center-radius form by completing the square.
    \end{formula}

    \begin{theorem}[Three Perpendiculars Theorem]
        Let $\Pi$ be a plane, let $P\notin\Pi$, and let $H$ be the orthogonal
        projection of $P$ onto $\Pi$. Let $Q\in\Pi$ with $Q\neq H$, and let
        $\ell\subseteq\Pi$ be a line through $Q$. If $HQ\perp\ell$, then
        \[
            PQ\perp\ell.
        \]
        Equivalently, a line in a plane that is perpendicular to the orthogonal
        projection of an oblique line is perpendicular to the oblique line
        itself.
    \end{theorem}

    \begin{remark}
        In coordinate geometry, the three perpendiculars theorem is often
        replaced by dot products and projections. Still, it is a standard
        spatial-geometry fact from the classical school curriculum.
    \end{remark}

    \begin{definition}[Conic Section]
        A conic section is a curve obtained by intersecting a plane with a
        double cone. The nondegenerate conics are:
        \[
            \text{ellipse},\qquad
            \text{parabola},\qquad
            \text{hyperbola}.
        \]
        Equivalently, a conic can be described using a focus, a directrix, and
        an eccentricity $e$.
    \end{definition}

    \begin{formula}[Focus--Directrix Definition of Conics]
        A conic is the set of points $P$ satisfying
        \[
            \frac{\dist(P,\text{focus})}
                 {\dist(P,\text{directrix})}
            =
            e.
        \]
        It is an ellipse if
        \[
            0<e<1,
        \]
        a parabola if
        \[
            e=1,
        \]
        and a hyperbola if
        \[
            e>1.
        \]
    \end{formula}

    \begin{formula}[Standard Forms of Conics]
        After translation and, when necessary, rotation of coordinates, every
        nondegenerate real conic can be put into one of the following standard
        forms.

        {\small
        \renewcommand{\arraystretch}{1.18}
        \begin{longtable}{@{}L{0.16\textwidth}L{0.46\textwidth}L{0.30\textwidth}@{}}
            \toprule
            Conic & Standard form & Basic condition and feature \\
            \midrule
            \endfirsthead

            \toprule
            Conic & Standard form & Basic condition and feature \\
            \midrule
            \endhead

            \bottomrule
            \endlastfoot

            Circle &
            $\displaystyle (x-h)^2+(y-k)^2=r^2$ &
            center $(h,k)$, radius $r>0$ \\
            Ellipse &
            $\displaystyle \frac{(x-h)^2}{a^2}+\frac{(y-k)^2}{b^2}=1$,
            $a\geq b>0$ &
            $c^2=a^2-b^2$, foci at distance $c$ from the center \\
            Parabola &
            $\displaystyle (y-k)^2=4p(x-h)$ or
            $\displaystyle (x-h)^2=4p(y-k)$ &
            focus--directrix distance parameter $p\neq0$ \\
            Hyperbola &
            $\displaystyle \frac{(x-h)^2}{a^2}-\frac{(y-k)^2}{b^2}=1$
            or with the signs reversed &
            $c^2=a^2+b^2$, foci at distance $c$ from the center \\
        \end{longtable}
        }
        The general second-degree equation
        \[
            Ax^2+Bxy+Cy^2+Dx+Ey+F=0
        \]
        is classified, in the nondegenerate case, by the quadratic part; in
        particular $B^2-4AC<0$, $=0$, and $>0$ correspond to ellipse-type,
        parabola-type, and hyperbola-type conics, respectively.
    \end{formula}

    \begin{formula}[Eccentricity of an Ellipse and a Hyperbola]
        For an ellipse
        \[
            \frac{x^2}{a^2}+\frac{y^2}{b^2}=1
            \qquad (a\geq b>0),
        \]
        the eccentricity is
        \[
            e=\frac{c}{a},
            \qquad
            c^2=a^2-b^2.
        \]
        For a hyperbola
        \[
            \frac{x^2}{a^2}-\frac{y^2}{b^2}=1,
        \]
        the eccentricity is
        \[
            e=\frac{c}{a},
            \qquad
            c^2=a^2+b^2.
        \]
    \end{formula}

    \begin{formula}[Tangent Lines to Standard Conics]
        At a point $(x_1,y_1)$ on the indicated conic, the tangent line is
        \[
            xx_1+yy_1=r^2
            \qquad
            \text{for }x^2+y^2=r^2,
        \]
        \[
            yy_1=2a(x+x_1)
            \qquad
            \text{for }y^2=4ax,
        \]
        \[
            \frac{xx_1}{a^2}+\frac{yy_1}{b^2}=1
            \qquad
            \text{for }\frac{x^2}{a^2}+\frac{y^2}{b^2}=1,
        \]
        and
        \[
            \frac{xx_1}{a^2}-\frac{yy_1}{b^2}=1
            \qquad
            \text{for }\frac{x^2}{a^2}-\frac{y^2}{b^2}=1.
        \]
    \end{formula}

    \begin{theorem}[Reflection Property of Conics]
        At a point of a nondegenerate conic, the tangent makes equal angles with
        the two relevant focal directions. Consequently, a ray emitted from one
        focus of an ellipse reflects toward the other focus; a ray parallel to
        the axis of a parabola reflects through its focus; and a ray directed
        toward one focus of a hyperbola reflects along a ray whose backward
        extension passes through the other focus.
    \end{theorem}

    \begin{remark}
        The reflection property gives a geometric reason why conics appear in
        optics, satellite dishes, whispering galleries, and orbital mechanics.
    \end{remark}

    \begin{theorem}[Quadratic Classification of Conics]
        Consider a real quadratic equation
        \[
            Ax^2+Bxy+Cy^2+Dx+Ey+F=0.
        \]
        The discriminant
        \[
            B^2-4AC
        \]
        determines the type of the nondegenerate conic:
        \[
            B^2-4AC<0
            \quad\Rightarrow\quad
            \text{ellipse-type},
        \]
        \[
            B^2-4AC=0
            \quad\Rightarrow\quad
            \text{parabola-type},
        \]
        and
        \[
            B^2-4AC>0
            \quad\Rightarrow\quad
            \text{hyperbola-type}.
        \]
    \end{theorem}

    \begin{remark}
        This is the analytic-geometry counterpart of classifying conic sections
        by their shapes.
    \end{remark}

    \begin{formula}[Polar Equation of a Conic]
        With a focus at the pole, a conic has polar equation
        \[
            r
            =
            \frac{\ell}{1+e\cos\theta}
        \]
        after choosing a suitable axis. Here $e$ is the eccentricity and
        $\ell$ is the semi-latus rectum.
    \end{formula}

    \begin{remark}
        The same equation describes ellipses, parabolas, and hyperbolas
        depending on whether $e<1$, $e=1$, or $e>1$.
    \end{remark}

    \begin{definition}[Apollonius Circle]
        Given two fixed points $A$ and $B$ and a positive constant
        $k\neq 1$, the locus of points $P$ satisfying
        \[
            \frac{PA}{PB}=k
        \]
        is a circle. This circle is called an Apollonius circle.
    \end{definition}

    \begin{remark}
        Apollonius circles turn a ratio of distances into a concrete circle,
        which is a common hidden structure in geometry problems.
    \end{remark}

    \begin{theorem}[Descartes' Circle Theorem; Soddy Circle Theorem]
        Suppose four circles are mutually tangent, and let their curvatures be
        \[
            k_1,k_2,k_3,k_4,
        \]
        where curvature means the reciprocal of the radius, with sign convention
        for an enclosing circle. Then
        \[
            (k_1+k_2+k_3+k_4)^2
            =
            2(k_1^2+k_2^2+k_3^2+k_4^2).
        \]
    \end{theorem}

    \begin{remark}
        Descartes' Circle Theorem is a striking formula: tangency of four
        circles becomes one quadratic equation in their curvatures.
    \end{remark}

    \begin{formula}[Boundary Lattice Points on a Segment]
        If a lattice segment has displacement $(\Delta x,\Delta y)$, then the
        number of lattice points on it, including both endpoints, is
        \[
            \gcd(|\Delta x|,|\Delta y|)+1.
        \]
        Summing the gcd contributions over polygon edges gives the number $B$ of
        boundary lattice points without double-counting vertices.
    \end{formula}

    \begin{theorem}[Pick's Theorem]
        Let $P$ be a simple polygon whose vertices lie on lattice points in the
        plane. If $I$ is the number of lattice points strictly inside $P$ and
        $B$ is the number of lattice points on the boundary of $P$, then
        \[
            \area(P)
            =
            I+\frac{B}{2}-1.
        \]
    \end{theorem}

    \begin{remark}
        Pick's Theorem is a beautiful bridge between geometry and arithmetic:
        area is determined by counting lattice points.
    \end{remark}

    \chapter{Number Theory}

    \relatedcourses{(MATH304) Introduction to Number Theory}

    \relatedbooks{Introduction to Number Theory
        (William W. Adams and Larry Joel Goldstein)}

    \begin{definition}[Divisibility, Prime, Greatest Common Divisor, and Congruence]
        For $a,b\in\Z$, write $a\mid b$ if
        \[
            \exists k\in\Z,
            \qquad
            b=ak.
        \]
        An integer $p>1$ is prime if its only positive divisors are $1$ and
        $p$. For integers $a,b$, not both zero, $\gcd(a,b)$ is their greatest
        positive common divisor. For $n\in\N$,
        \[
            a\equiv b\pmod n
            \quad\Longleftrightarrow\quad
            n\mid(a-b).
        \]
    \end{definition}

    \begin{theorem}[Fundamental Theorem of Arithmetic]
        Every integer $n>1$ can be written as a product of prime numbers.
        Moreover, this factorization is unique up to the order of the factors.
        Equivalently,
        \[
            n=p_1^{\alpha_1}p_2^{\alpha_2}\cdots p_k^{\alpha_k},
        \]
        where $p_1,\dots,p_k$ are distinct primes and
        $\alpha_1,\dots,\alpha_k$ are positive integers, and this expression is
        unique up to reordering.
    \end{theorem}

    \begin{remark}
        This theorem is often called unique prime factorization. It is the
        reason prime numbers are treated as the basic building blocks of
        integers.
    \end{remark}

    \begin{theorem}[Euclid's Theorem; Infinitude of Primes]
        There are infinitely many prime numbers.
    \end{theorem}

    \begin{idea}
        If $p_1,\dots,p_n$ are finitely many primes, consider
        \[
            p_1p_2\cdots p_n+1.
        \]
        This number is not divisible by any prime on the list.
    \end{idea}

    \begin{theorem}[Brun's Theorem]
        The reciprocal sum over twin-prime pairs converges:
        \[
            \sum_{\substack{p\text{ prime}\\p+2\text{ prime}}}
            \left(\frac{1}{p}+\frac{1}{p+2}\right)
            <\infty.
        \]
        Its value is called Brun's constant. This contrasts with Euler's
        theorem that the sum of the reciprocals of all primes diverges.
    \end{theorem}

    \begin{theorem}[B\'ezout's Identity; B\'ezout's Lemma]
        Let $a,b\in\Z$, not both zero. Then there exist integers $x,y\in\Z$
        s.t.
        \[
            ax+by=\gcd(a,b).
        \]
        In particular, if $\gcd(a,b)=1$, then there exist integers $x,y\in\Z$
        s.t.
        \[
            ax+by=1.
        \]
    \end{theorem}

    \begin{remark}
        B\'ezout's identity is the algebraic heart of the Euclidean algorithm and
        modular inverses.
    \end{remark}

    \begin{algorithm}[Euclidean Algorithm]
        To compute $\gcd(a,b)$ for integers $a,b\in\Z$, repeatedly divide the
        larger number by the smaller number and replace the pair by
        \[
            (b,r),
        \]
        where $r$ is the remainder. Continue until the remainder becomes $0$.
        The last nonzero remainder is $\gcd(a,b)$.
    \end{algorithm}

    \begin{definition}[Gaussian Integers]
        The Gaussian integers are
        \[
            \Z[i]
            =
            \{a+bi\mid a,b\in\Z\}.
        \]
        Their norm is
        \[
            N(a+bi)=a^2+b^2,
        \]
        and it is multiplicative:
        \[
            N(zw)=N(z)N(w).
        \]
        With the usual addition and multiplication, $\Z[i]$ inherits its
        arithmetic from $\C$. In the terminology of the Algebra chapter,
        $\Z[i]$ is a Euclidean domain with respect to this norm, so a Euclidean
        algorithm, greatest common divisors, and unique factorization are
        available in $\Z[i]$ just as they are in $\Z$.
    \end{definition}

    \begin{proposition}[Basic Rules of Congruences]
        If
        \[
            a\equiv b\pmod n,
            \qquad
            c\equiv d\pmod n,
        \]
        then
        \[
            a+c\equiv b+d\pmod n,
            \qquad
            ac\equiv bd\pmod n.
        \]
        If $\gcd(c,n)=1$, then
        \[
            ac\equiv bc\pmod n
            \quad\Longrightarrow\quad
            a\equiv b\pmod n.
        \]
    \end{proposition}


    \begin{theorem}[Linear Congruence Criterion]
        The congruence
        \[
            ax\equiv b\pmod n
        \]
        has a solution if and only if
        \[
            \gcd(a,n)\mid b.
        \]
        If $d=\gcd(a,n)$ and $d\mid b$, then there are exactly $d$ solutions
        modulo $n$.
    \end{theorem}

    \begin{formula}[Least Common Multiple from Prime Factorization]
        If
        \[
            a=\prod_p p^{\alpha_p},
            \qquad
            b=\prod_p p^{\beta_p},
        \]
        then
        \[
            \gcd(a,b)=\prod_p p^{\min(\alpha_p,\beta_p)},
            \qquad
            \lcm(a,b)=\prod_p p^{\max(\alpha_p,\beta_p)}.
        \]
        More generally, the least common multiple of several integers is found
        by taking the largest exponent of each prime appearing in the list.
    \end{formula}

    \begin{formula}[Basic Divisibility Tests]
        For a decimal integer $N$, the following tests are useful.
        \begin{itemize}
            \item Divisibility by $2$, $4$, and $8$ depends on the last $1$, $2$,
                  and $3$ digits, respectively.
            \item Divisibility by $3$ and $9$ depends on the digit sum.
            \item Divisibility by $5$ depends on the last digit.
            \item Divisibility by $11$ depends on the alternating sum of the
                  digits.
        \end{itemize}
    \end{formula}

    \begin{strategy}
        Digit problems often combine an lcm condition, a digit-sum condition,
        and a last-digits condition. First minimize the number of digits, then
        arrange the final few digits to satisfy the strongest divisibility
        conditions.
    \end{strategy}

    \begin{definition}[Complete Residue System]
        A set of $n$ integers is a complete residue system modulo $n$ if its
        elements represent every residue class modulo $n$ exactly once. To prove
        that $a_0,\dots,a_{n-1}$ form such a system, it is enough to prove
        \[
            a_i\equiv a_j\pmod n
            \quad\Longrightarrow\quad
            i=j.
        \]
    \end{definition}

    \begin{strategy}[Modular Casework and Base Representation]
        When a problem repeats division by $N$, rewrite the integer in base $N$.
        The quotient deletes the final digit and the remainder reveals it. For
        divisibility or existence problems, first split integers into residue
        classes and look for a complete residue system rather than checking
        values individually.
    \end{strategy}

    \begin{strategy}[Infinite Descent]
        Assume a positive-integer solution exists and choose one minimizing a
        natural positive quantity. Construct a strictly smaller solution of the
        same type, contradicting well-ordering.
    \end{strategy}

    \begin{strategy}[Vieta Jumping]
        If an integer equation is quadratic in one variable, regard that
        variable as one root of a monic quadratic. Vi\`ete's formulas give the
        other integral root. Prove that it yields a smaller positive solution
        and apply infinite descent.
    \end{strategy}

    \begin{theorem}[Frobenius Coin Theorem for Two Denominations]
        If $a$ and $b$ are coprime positive integers, the largest integer not
        representable as
        \[
            ax+by,
            \qquad
            x,y\in\Nzero,
        \]
        is $ab-a-b$. Exactly
        \[
            \frac{(a-1)(b-1)}{2}
        \]
        positive integers are not representable.
    \end{theorem}

    \begin{theorem}[Chinese Remainder Theorem; CRT]
        Let $m_1,\dots,m_k$ be pairwise relatively prime positive integers.
        Then the system
        \[
            x\equiv a_1\pmod {m_1},
            \qquad
            x\equiv a_2\pmod {m_2},
            \qquad
            \dots,
            \qquad
            x\equiv a_k\pmod {m_k}
        \]
        has a solution. Moreover, the solution is unique modulo
        \[
            M=m_1m_2\cdots m_k.
        \]
    \end{theorem}

    \begin{strategy}
        CRT is the standard theorem for combining several congruence conditions
        into one congruence.
    \end{strategy}

    \begin{definition}[Arithmetic and Multiplicative Functions]
        An arithmetic function is a function $f\colon \N\to\C$. It is
        multiplicative if
        \[
            \gcd(m,n)=1
            \quad\Longrightarrow\quad
            f(mn)=f(m)f(n).
        \]
        Important multiplicative arithmetic functions include Euler's totient
        $\varphi(n)$, the divisor-counting function $\tau(n)$, the
        divisor-sum function $\sigma(n)$, and the M\"obius function $\mu(n)$.
        Multiplicativity reduces their evaluation to prime powers.
    \end{definition}

    \begin{formula}[Euler's Phi Function Formula]
        If
        \[
            n=p_1^{\alpha_1}p_2^{\alpha_2}\cdots p_k^{\alpha_k}
        \]
        is the prime factorization of $n$, then
        \[
            \varphi(n)
            =
            n\prod_{i=1}^{k}\left(1-\frac{1}{p_i}\right).
        \]
        In particular, for a prime power,
        \[
            \varphi(p^\alpha)=p^\alpha-p^{\alpha-1}
            =
            p^{\alpha-1}(p-1).
        \]
    \end{formula}

    \begin{formula}[Sum-of-Divisors Function]
        If
        \[
            n=p_1^{\alpha_1}\cdots p_k^{\alpha_k},
        \]
        then the number of positive divisors of $n$ is
        \[
            \tau(n)=(\alpha_1+1)\cdots(\alpha_k+1),
        \]
        and the sum of positive divisors of $n$ is
        \[
            \sigma(n)
            =
            \prod_{i=1}^{k}\frac{p_i^{\alpha_i+1}-1}{p_i-1}.
        \]
    \end{formula}

    \begin{definition}[$p$-Adic Valuation, $p$-Adic Absolute Value, and $p$-Adic Numbers]
        Fix a prime $p$. For a nonzero rational number $x$, write uniquely
        \[
            x=p^{\nu_p(x)}\frac{a}{b},
            \qquad
            p\nmid ab,
        \]
        where $\nu_p(x)\in\Z$ is the $p$-adic valuation of $x$. For a nonzero
        integer, $\nu_p(n)$ is therefore the exponent of $p$ in its prime
        factorization. Define
        \[
            |x|_p=p^{-\nu_p(x)},
            \qquad
            |0|_p=0.
        \]
        The $p$-adic absolute value satisfies the strong triangle inequality
        \[
            |x+y|_p\leq\max\{|x|_p,|y|_p\}.
        \]
        The field $\mathbb{Q}_p$ of $p$-adic numbers is the completion of
        $\Q$ with respect to this absolute value, just as $\R$ is a completion
        of $\Q$ with respect to the usual absolute value.
    \end{definition}

    \begin{formula}[$p$-Adic Valuation of a Factorial; Legendre's Formula]
        For a prime $p$ and a positive integer $n$,
        \[
            \nu_p(n!)
            =
            \sum_{k=1}^{\infty}
            \left\lfloor\frac{n}{p^k}\right\rfloor.
        \]
        Only finitely many terms are nonzero. This formula gives the exponent
        of $p$ in $n!$ and is especially useful for trailing-zero and
        divisibility questions.
    \end{formula}

    \begin{definition}[M\"obius Function]
        The M\"obius function is
        \[
            \mu(n)
            =
            \begin{cases}
                1, & n=1,\\
                0, & p^2\mid n\text{ for some prime }p,\\
                (-1)^k, & n\text{ is a product of }k\text{ distinct primes}.
            \end{cases}
        \]
        It satisfies
        \[
            \sum_{d\mid n}\mu(d)
            =
            \begin{cases}
                1, & n=1,\\
                0, & n>1.
            \end{cases}
        \]
    \end{definition}

    \begin{theorem}[M\"obius Inversion Formula]
        If arithmetic functions $f$ and $g$ satisfy
        \[
            g(n)=\sum_{d\mid n}f(d),
        \]
        then
        \[
            f(n)=\sum_{d\mid n}\mu(d)g\!\left(\frac{n}{d}\right).
        \]
    \end{theorem}

    \begin{theorem}[Fermat's Little Theorem]
        Let $p$ be a prime number. If $p\nmid a$, then
        \[
            a^{p-1}\equiv 1\pmod p.
        \]
        Equivalently, for every integer $a$,
        \[
            a^p\equiv a\pmod p.
        \]
    \end{theorem}

    \begin{theorem}[Euler's Theorem; Euler--Fermat Theorem]
        If $\gcd(a,n)=1$, then
        \[
            a^{\varphi(n)}\equiv 1\pmod n.
        \]
    \end{theorem}

    \begin{remark}
        Fermat's Little Theorem is the special case of Euler's Theorem when
        $n=p$ is prime, since $\varphi(p)=p-1$.
    \end{remark}

    \begin{proposition}[Multiplicative Order]
        If $\gcd(a,n)=1$, the multiplicative order $\ord_n(a)$ is the least
        positive integer $t$ s.t.
        \[
            a^t\equiv1\pmod n.
        \]
        The invertible residue classes modulo $n$ form a set of size
        $\varphi(n)$ and, in the group terminology of the Algebra chapter,
        Lagrange's theorem gives
        \[
            \ord_n(a)\mid\varphi(n).
        \]
        Consequently, for every positive integer $k$,
        \[
            a^k\equiv 1\pmod n
            \quad\Longleftrightarrow\quad
            \ord_n(a)\mid k.
        \]
    \end{proposition}

    \begin{definition}[Primitive Root Modulo $n$]
        If $\gcd(g,n)=1$, then $g$ is a primitive root modulo $n$ if
        \[
            \ord_n(g)=\varphi(n).
        \]
        Equivalently, in the group terminology of the Algebra chapter, the
        residue class $[g]_n$ generates $(\Z/n\Z)^\times$.
    \end{definition}

    \begin{theorem}[Primitive Root Theorem; Gauss's Theorem on Primitive Roots]
        A primitive root modulo $n$ exists if and only if
        \[
            n\in\{1,2,4,p^k,2p^k\},
        \]
        where $p$ is an odd prime and $k\geq 1$. Equivalently, in the language
        of the Algebra chapter, $(\Z/n\Z)^\times$ is cyclic exactly for these
        values of $n$.
    \end{theorem}

    \begin{definition}[Finite Field]
        A finite field is a finite set equipped with addition and multiplication
        satisfying the field axioms stated in the Algebra chapter. Its number of
        elements is called its order.
    \end{definition}

    \begin{theorem}[Existence and Uniqueness of Finite Fields]
        Every finite field has $q=p^n$ elements for some prime $p$ and positive
        integer $n$. Conversely, for each prime power $q=p^n$, there exists a
        field $\mathbb{F}_q$ with $q$ elements, unique up to isomorphism. Its
        nonzero elements form a cyclic multiplicative group
        $\mathbb{F}_q^\times$ of order $q-1$.
    \end{theorem}

    \begin{algorithm}[Modular Exponentiation by Period Reduction]
        To compute $a^n\pmod m$, first reduce $a$ modulo $m$. If
        $\gcd(a,m)=1$, determine a convenient exponent $t>0$ s.t.
        \[
            a^t\equiv1\pmod m,
        \]
        for example from the multiplicative order of $a$, Euler's theorem, or
        Fermat's little theorem when applicable. Write
        \[
            n=qt+r,
            \qquad
            0\leq r<t,
        \]
        and reduce
        \[
            a^n\equiv a^r\pmod m.
        \]
        If $a$ is not a unit modulo $m$, use repeated squaring directly or
        analyze prime-power factors separately rather than applying a unit-group
        period without justification.
    \end{algorithm}

    \begin{strategy}
        Last-digit and last-two-digit questions are modular arithmetic problems.
        Work modulo $10$ or $100$, not with the full number.
    \end{strategy}

    \begin{theorem}[Wilson's Theorem]
        A positive integer $p>1$ is prime if and only if
        \[
            (p-1)!\equiv -1\pmod p.
        \]
    \end{theorem}

    \begin{remark}
        Wilson's theorem is famous because it gives a clean factorial
        characterization of primes, even though it is not efficient as a
        primality test.
    \end{remark}

    \begin{theorem}[Lucas's Theorem]
        Let $p$ be prime and write
        \[
            n=\sum_{i=0}^{r}n_ip^i,
            \qquad
            k=\sum_{i=0}^{r}k_ip^i
        \]
        in base $p$. Then
        \[
            \binom{n}{k}
            \equiv
            \prod_{i=0}^{r}\binom{n_i}{k_i}
            \pmod p.
        \]
    \end{theorem}

    \begin{theorem}[Kummer's Theorem]
        The valuation
        \[
            \nu_p\!\left(\binom{n}{k}\right)
        \]
        equals the number of carries when $k$ and $n-k$ are added in base $p$.
    \end{theorem}

    \begin{theorem}[Lifting the Exponent Lemma; LTE]
        Let $p$ be an odd prime, and suppose that $p\mid x-y$ while
        $p\nmid xy$. Then, for every positive integer $n$,
        \[
            \nu_p(x^n-y^n)
            =
            \nu_p(x-y)+\nu_p(n).
        \]
        A standard companion form is: if $n$ is odd, $p\mid x+y$, and
        $p\nmid xy$, then
        \[
            \nu_p(x^n+y^n)
            =
            \nu_p(x+y)+\nu_p(n).
        \]
    \end{theorem}

    \begin{fact}[Small-Modulus Residue Filters]
        Small powers have restricted residues modulo small integers:
        \[
            x^2\equiv 0,1\pmod 4,
            \qquad
            x^2\equiv 0,1,4\pmod 8.
        \]
        Also,
        \[
            x^3\equiv0,\pm1\pmod9,
            \qquad
            x^4\equiv0,1\pmod{16}.
        \]
        In particular, if $x$ is odd, then $x^2\equiv1\pmod8$.
        Hence an integer congruent to $2$ or $3$ modulo $4$ cannot be a
        square, and an integer congruent to $2,3,5,6,7$ modulo $8$ cannot be a
        square. These tests can often rule out a proposed square, cube, or fourth
        power before any large calculation is attempted.
    \end{fact}

    \begin{tactic}
        For Diophantine equations, first reduce modulo $4$, $8$, $3$, $5$, or $9$. Many impossible equations are eliminated by
        the small list of possible square residues.
    \end{tactic}

    \begin{definition}[Legendre Symbol]
        Let $p$ be an odd prime. The Legendre symbol is defined by
        \[
            \left(\frac{a}{p}\right)
            =
            \begin{cases}
                1, & \text{if }a\text{ is a nonzero quadratic residue modulo }p, \\
                -1, & \text{if }a\text{ is a quadratic nonresidue modulo }p, \\
                0, & \text{if }p\mid a.
            \end{cases}
        \]
    \end{definition}

    \begin{theorem}[Euler's Criterion]
        Let $p$ be an odd prime. For any integer $a$,
        \[
            \left(\frac{a}{p}\right)
            \equiv
            a^{(p-1)/2}
            \pmod p.
        \]
        Equivalently, if $p\nmid a$, then
        \[
            a^{(p-1)/2}
            \equiv
            \begin{cases}
                1\pmod p, & \text{if }a\text{ is a quadratic residue modulo }p, \\
                -1\pmod p, & \text{if }a\text{ is a quadratic nonresidue modulo }p.
            \end{cases}
        \]
    \end{theorem}

    \begin{theorem}[Law of Quadratic Reciprocity; Gauss's Theorem of Quadratic Reciprocity]
        Let $p$ and $q$ be distinct odd primes. Then
        \[
            \left(\frac{p}{q}\right)
            \left(\frac{q}{p}\right)
            =
            (-1)^{\frac{p-1}{2}\frac{q-1}{2}}.
        \]
        Equivalently,
        \[
            \left(\frac{p}{q}\right)
            =
            \left(\frac{q}{p}\right)
        \]
        unless both $p$ and $q$ are congruent to $3$ modulo $4$.
    \end{theorem}

    \begin{remark}
        Quadratic reciprocity is one of the crown jewels of elementary number
        theory. Gauss called it the golden theorem.
    \end{remark}

    \begin{theorem}[Fermat's Two-Square Theorem; Fermat's Theorem on Sums of Two Squares]
        An odd prime $p$ can be written as a sum of two squares,
        \[
            p=a^2+b^2
        \]
        for some integers $a,b\in\Z$, if and only if
        \[
            p\equiv 1\pmod 4.
        \]
        Also,
        \[
            2=1^2+1^2.
        \]
    \end{theorem}

    \begin{theorem}[Sum of Two Squares Criterion]
        A positive integer
        \[
            n=2^a\prod_{p\equiv1\ (\mathrm{mod}\ 4)}p^{\alpha_p}
              \prod_{q\equiv3\ (\mathrm{mod}\ 4)}q^{\beta_q}
        \]
        is representable as $n=x^2+y^2$ with integers $x,y$ if and only if every
        exponent $\beta_q$ is even.
    \end{theorem}

    \begin{formula}[Brahmagupta--Fibonacci Identity]
        \[
            (a^2+b^2)(c^2+d^2)
            =
            (ac-bd)^2+(ad+bc)^2
            =
            (ac+bd)^2+(ad-bc)^2.
        \]
    \end{formula}

    \begin{formula}[Sophie Germain's Identity]
        \[
            a^4+4b^4
            =
            (a^2-2ab+2b^2)(a^2+2ab+2b^2).
        \]
    \end{formula}

    \begin{theorem}[Lagrange's Four-Square Theorem]
        Every nonnegative integer can be written as a sum of four integer
        squares. That is, for every $n\in\Nzero$, there exist integers
        $a,b,c,d\in\Z$ s.t.
        \[
            n=a^2+b^2+c^2+d^2.
        \]
    \end{theorem}

    \begin{remark}
        The two-square theorem has a congruence condition. The four-square
        theorem has no exception at all.
    \end{remark}

    \begin{theorem}[Continued-Fraction Best Approximation Theorem]
        Let $p_k/q_k$ be the convergents of the simple continued fraction of an
        irrational number $\alpha$. Then
        \[
            \left|\alpha-\frac{p_k}{q_k}\right|
            <
            \frac{1}{q_kq_{k+1}}.
        \]
        Moreover, the convergents are best approximations of the second kind:
        if $p,q\in\Z$, $0<q<q_{k+1}$, and $p/q\neq p_k/q_k$, then
        \[
            |q_k\alpha-p_k|<|q\alpha-p|.
        \]
    \end{theorem}

    \begin{strategy}
        Problems asking for an unusually accurate fraction with the smallest
        possible denominator often signal a continued-fraction expansion.
    \end{strategy}

    \begin{definition}[Pell Equation]
        A Pell equation has the form
        \[
            x^2-Dy^2=1,
        \]
        where $D$ is a positive nonsquare integer. If $(x_1,y_1)$ is the
        fundamental positive solution, then all positive solutions satisfy
        \[
            x_n+y_n\sqrt{D}
            =
            (x_1+y_1\sqrt{D})^n.
        \]
        Continued fractions of $\sqrt{D}$ produce the fundamental solution.
    \end{definition}

    \begin{theorem}[Dirichlet's Theorem on Arithmetic Progressions]
        If $a$ and $d$ are relatively prime positive integers, then the
        arithmetic progression
        \[
            a,\ a+d,\ a+2d,\ a+3d,\dots
        \]
        contains infinitely many primes.
    \end{theorem}

    \begin{remark}
        Dirichlet's theorem is a major result connecting primes and arithmetic
        progressions. The condition $\gcd(a,d)=1$ is necessary.
    \end{remark}

    \begin{theorem}[Bertrand's Postulate; Chebyshev's Theorem]
        For every integer $n>1$, there exists a prime $p$ s.t.
        \[
            n<p<2n.
        \]
    \end{theorem}

    \begin{theorem}[Prime Number Theorem]
        Let $\pi(x)$ be the number of primes less than or equal to $x$. Then
        \[
            \pi(x)\sim \frac{x}{\ln x}
            \qquad
            \text{as }x\to\infty.
        \]
    \end{theorem}

    \begin{idea}
        The Prime Number Theorem says that the ``probability'' that a large
        integer near $x$ is prime is roughly $1/\ln x$. This is a mental model,
        not a literal probability statement.
    \end{idea}

    \begin{definition}[Riemann Zeta Function]
        For $\Real(s)>1$, the Riemann zeta function is defined by
        \[
            \zeta(s)
            =
            \sum_{n=1}^{\infty}\frac{1}{n^s}.
        \]
        It also has the Euler product
        \[
            \zeta(s)
            =
            \prod_{p\text{ prime}}
            \frac{1}{1-p^{-s}}.
        \]
    \end{definition}

    \begin{remark}
        The Euler product is the bridge between the zeta function and prime
        numbers. The Riemann Hypothesis asserts that the nontrivial zeros of
        $\zeta(s)$ all have real part $1/2$.
    \end{remark}

    \begin{definition}[Perfect Number and Mersenne Prime]
        A positive integer $N$ is perfect if it equals the sum of its positive
        proper divisors, equivalently
        \[
            \sigma(N)=2N.
        \]
        A Mersenne prime is a prime of the form
        \[
            2^p-1.
        \]
        If $2^p-1$ is prime, then $p$ must be prime.
    \end{definition}

    \begin{theorem}[Euclid--Euler Theorem on Perfect Numbers]
        An even positive integer $N$ is perfect if and only if
        \[
            N=2^{p-1}(2^p-1),
        \]
        where $2^p-1$ is a Mersenne prime.
    \end{theorem}

    \begin{theorem}[Fermat's Last Theorem]
        For every integer $n\geq 3$, there do not exist positive integers
        $a,b,c$ s.t.
        \[
            a^n+b^n=c^n.
        \]
    \end{theorem}

    \begin{remark}
        Fermat's Last Theorem was proved by Andrew Wiles. It is one of the most
        famous theorems in the history of mathematics.
    \end{remark}

    \chapter{Algebra}

    \relatedcourses{(MATH301) Modern Algebra I, (MATH302) Modern Algebra II}

    \relatedbooks{Abstract Algebra (David S. Dummit and Richard M. Foote)}

    \begin{definition}[Group]
        A group is a set $G$ with a binary operation $\ast\colon G\times G\to G$ s.t.
        \[
            \forall a,b,c\in G,
            \qquad
            (a\ast b)\ast c=a\ast(b\ast c),
        \]
        \[
            \exists e\in G\ \text{s.t.}\quad
            \forall a\in G,
            \qquad
            e\ast a=a\ast e=a,
        \]
        and
        \[
            \forall a\in G,
            \qquad
            \exists a^{-1}\in G\ \text{s.t.}\quad
            a\ast a^{-1}=a^{-1}\ast a=e.
        \]
        The group is abelian if
        \[
            \forall a,b\in G,
            \qquad
            a\ast b=b\ast a.
        \]
    \end{definition}


    \begin{definition}[Subgroup, Homomorphism, and Kernel]
        A subset $H\subseteq G$ is a subgroup, written $H\leq G$, if it is a
        group under the operation of $G$. A map $\varphi\colon G\to H$ is a group
        homomorphism if
        \[
            \varphi(ab)=\varphi(a)\varphi(b)
            \qquad
            \forall a,b\in G.
        \]
        Its kernel and image are
        \[
            \ker\varphi=\{g\in G\mid\varphi(g)=e_H\},
            \qquad
            \image\varphi=\{\varphi(g)\mid g\in G\}.
        \]
        A bijective homomorphism is an \emph{isomorphism}. An isomorphism
        $G\to G$ is an \emph{automorphism}. Isomorphic groups have the same
        group-theoretic structure up to relabeling.
    \end{definition}

    \begin{definition}[Generated Subgroup, Generating Set, and Finitely Generated Group]
        If $A\subseteq G$, the subgroup generated by $A$, denoted
        $\langle A\rangle$, is the smallest subgroup of $G$ containing $A$.
        Equivalently, it consists of all finite products of elements of $A$ and
        their inverses. The set $A$ is a \emph{generating set} for $G$ if
        $G=\langle A\rangle$. The group $G$ is \emph{finitely generated} if
        it has a finite generating set. For a single element $g$, the notation
        $\langle g\rangle$ denotes the cyclic subgroup generated by $g$.
    \end{definition}

    \begin{proposition}[Countability of Finitely Generated Groups]
        Every finitely generated group is at most countable. Every finite group
        is finitely generated. Consequently, an uncountable group cannot be
        finitely generated.
    \end{proposition}

    \begin{definition}[Normal Subgroup and Simple Group]
        A subgroup $N\leq G$ is normal, written $N\trianglelefteq G$, if
        \[
            gNg^{-1}=N
            \qquad
            \text{for every }g\in G.
        \]
        A nontrivial group $G$ is simple if its only normal subgroups are
        $\{e\}$ and $G$.
    \end{definition}


    \begin{definition}[Coset and Index]
        If $H\leq G$ and $g\in G$, the left and right cosets of $H$ represented
        by $g$ are
        \[
            gH=\{gh\mid h\in H\},
            \qquad
            Hg=\{hg\mid h\in H\}.
        \]
        The number of left cosets is the index of $H$ in $G$, denoted
        $[G:H]$. For finite groups,
        \[
            [G:H]=\frac{|G|}{|H|}.
        \]
    \end{definition}

    \begin{definition}[Quotient Group]
        If $N\trianglelefteq G$, the quotient group $G/N$ is the set of
        left cosets of $N$,
        \[
            G/N=\{gN\mid g\in G\},
        \]
        with multiplication
        \[
            (gN)(hN)=ghN.
        \]
        Normality is exactly the condition that makes this multiplication
        well-defined.
    \end{definition}

    \begin{definition}[Order of a Group and of an Element]
        The order of a finite group $G$ is its cardinality $|G|$. The order of
        an element $g\in G$, denoted $\ord(g)$, is the least positive integer
        $m$ s.t. $g^m=e$, if such an integer exists; otherwise $g$ has infinite
        order. Equivalently,
        \[
            \ord(g)=|\langle g\rangle|.
        \]
    \end{definition}

    \begin{definition}[Permutation, Symmetric Group, and Alternating Group]
        A permutation of $\{1,\dots,n\}$ is a bijection from the set to itself.
        Under composition, all such permutations form the symmetric group $S_n$.
        A transposition swaps two elements and fixes the others. A permutation is
        \emph{even} if it is a product of an even number of transpositions and
        \emph{odd} otherwise. The even permutations form the alternating group
        $A_n\leq S_n$, with
        \[
            |S_n|=n!,
            \qquad
            |A_n|=\frac{n!}{2}
            \quad(n\geq2).
        \]
    \end{definition}

    \begin{fact}[Smallest Nonabelian Simple Group]
        The alternating group $A_5$ is simple and has order
        \[
            |A_5|=60.
        \]
        It is the smallest nonabelian simple group. This fact, together with the
        Sylow theorems, is a standard shortcut in questions asking which small
        integers can occur as the order of a simple group.
    \end{fact}

    \begin{theorem}[Disjoint-Cycle Decomposition of a Permutation]
        Every permutation $\sigma\in S_n$ can be written as a product of
        disjoint cycles, uniquely up to the order of the cycles and cyclic
        rotation within each cycle. If $c(\sigma)$ is the number of cycles,
        including fixed points, then the minimum number of transpositions needed
        to express $\sigma$ is
        \[
            n-c(\sigma).
        \]
    \end{theorem}

    \begin{formula}[Inversions and the Sign of a Permutation]
        The inversion number is
        \[
            \inv(\sigma)
            =
            \left|
                \{(i,j)\mid i<j,\ \sigma(i)>\sigma(j)\}
            \right|.
        \]
        The sign satisfies
        \[
            \sgn(\sigma)
            =
            (-1)^{\inv(\sigma)}
            =
            (-1)^{n-c(\sigma)}.
        \]
        Thus any transposition reverses parity.
    \end{formula}

    \begin{definition}[Group Presentation]
        A group presentation
        \[
            G=\langle S\mid R\rangle
        \]
        specifies a generating set $S$ together with defining relations $R$.
        Presentations give compact descriptions of groups whose elements may be
        difficult to list explicitly.
    \end{definition}

    \begin{fact}[Dihedral Group]
        The symmetries of a regular $n$-gon form the dihedral group
        \[
            D_{2n}
            =
            \langle r,s\mid r^n=e,\ s^2=e,\ srs=r^{-1}\rangle.
        \]
        Its elements are
        \[
            e,r,r^2,\dots,r^{n-1},
            \qquad
            s,rs,r^2s,\dots,r^{n-1}s.
        \]
        This presentation is the standard algebraic model for rotation and
        reflection counting with Burnside's Lemma.
    \end{fact}

    \begin{definition}[Cyclic Group]
        A group $G$ is cyclic if there exists an element $g\in G$ s.t. every
        element of $G$ is a power of $g$. In this case, we write
        \[
            G=\langle g\rangle.
        \]
    \end{definition}

    \begin{theorem}[Lagrange's Theorem]
        If $G$ is a finite group and $H\leq G$, then
        \[
            |H|\mid |G|.
        \]
        Equivalently, the order of every subgroup divides the order of the group.
    \end{theorem}

    \begin{corollary}[Order of an Element]
        If $G$ is a finite group and $g\in G$, then the order of $g$ divides
        $|G|$. In particular,
        \[
            g^{|G|}=e.
        \]
    \end{corollary}

    \begin{theorem}[Cauchy's Theorem]
        If $G$ is a finite group and a prime $p$ divides $|G|$, then $G$ contains
        an element of order $p$.
    \end{theorem}

    \begin{corollary}[Groups of Prime Order]
        Every group of prime order $p$ is cyclic and isomorphic to $\Z_p$.
    \end{corollary}

    \begin{proposition}[Subgroups and Generators of a Finite Cyclic Group]
        If $G=\langle g\rangle$ has order $n$, then for every divisor $d\mid n$
        there is exactly one subgroup of order $d$. Moreover, $G$ has exactly
        $\varphi(n)$ generators, and
        \[
            \ord(g^k)=\frac{n}{\gcd(n,k)}.
        \]
    \end{proposition}

    \begin{theorem}[First Isomorphism Theorem for Groups]
        If $\varphi\colon G\to H$ is a group homomorphism, then
        \[
            G/\ker\varphi\cong \image\varphi.
        \]
    \end{theorem}

    \begin{definition}[Group Action, Orbit, and Stabilizer]
        A group action of $G$ on a set $X$ is a map
        \[
            G\times X\to X,
            \qquad
            (g,x)\mapsto g\cdot x,
        \]
        satisfying $e\cdot x=x$ and $(gh)\cdot x=g\cdot(h\cdot x)$. For
        $x\in X$, the orbit and stabilizer are
        \[
            G\cdot x=\{g\cdot x\mid g\in G\},
            \qquad
            G_x=\{g\in G\mid g\cdot x=x\}.
        \]
        Group actions translate algebraic symmetry into permutations of a set.
    \end{definition}

    \begin{theorem}[Cayley's Theorem]
        Every group is isomorphic to a subgroup of a symmetric group. In
        particular, every finite group of order $n$ is isomorphic to a subgroup
        of $S_n$.
    \end{theorem}

    \begin{theorem}[Orbit--Stabilizer Theorem]
        If a finite group $G$ acts on a set $X$ and $x\in X$, then
        \[
            |G\cdot x|
            =
            [G:G_x]
            =
            \frac{|G|}{|G_x|},
        \]
        where $G\cdot x$ is the orbit of $x$ and $G_x$ is its stabilizer.
    \end{theorem}

    \begin{definition}[Conjugacy, Center, and Centralizer]
        Elements $x,y\in G$ are conjugate if $y=gxg^{-1}$ for some $g\in G$.
        The conjugacy class of $x$ is
        \[
            \operatorname{Cl}_G(x)=\{gxg^{-1}\mid g\in G\}.
        \]
        The center and centralizer are
        \[
            Z(G)=\{z\in G\mid zg=gz\text{ for all }g\in G\},
            \qquad
            C_G(x)=\{g\in G\mid gx=xg\}.
        \]
        The center consists exactly of the elements whose conjugacy classes have
        size $1$.
    \end{definition}

    \begin{theorem}[Class Equation]
        Let $G$ be a finite group, and choose one representative $x_i$ from
        each conjugacy class not contained in the center $Z(G)$. Then
        \[
            |G|
            =
            |Z(G)|
            +
            \sum_i [G:C_G(x_i)].
        \]
    \end{theorem}

    \begin{theorem}[Sylow Theorems]
        Let $G$ be a finite group and suppose
        \[
            |G|=p^k m,
            \qquad
            p\nmid m,
        \]
        where $p$ is prime. Then $G$ has a subgroup of order $p^k$, called a
        Sylow $p$-subgroup. Moreover, all Sylow $p$-subgroups are conjugate, and
        the number $n_p$ of Sylow $p$-subgroups satisfies
        \[
            n_p\equiv 1\pmod p,
            \qquad
            n_p\mid m.
        \]
    \end{theorem}

    \begin{theorem}[Basic Finite $p$-Group Facts]
        If $G$ is a finite group of order $p^n$ for a prime $p$, then its center
        $Z(G)$ is nontrivial. In particular, every group of order $p^2$ is
        abelian. These are standard quick consequences of the class equation.
    \end{theorem}

    \begin{strategy}
        For a finite group, subgroup orders and element orders must divide the group
        order. When a
        prime divides $|G|$, Cauchy's theorem guarantees an element of that prime
        order.
    \end{strategy}

    \begin{fact}[Quaternion Group]
        The quaternion group is
        \[
            Q_8=\{\pm 1,\pm i,\pm j,\pm k\},
        \]
        with
        \[
            i^2=j^2=k^2=ijk=-1.
        \]
        It is nonabelian, but every subgroup of $Q_8$ is normal.
    \end{fact}

    \begin{remark}
        The quaternion group is a standard small group with behavior that is
        less intuitive than cyclic or dihedral groups, so it is a natural
        source of quick algebra questions.
    \end{remark}

    \begin{definition}[Direct Product and Free Group]
        For groups $G$ and $H$, the direct product $G\times H$ consists of
        pairs $(g,h)$ with componentwise multiplication. A free group
        $F(S)$ on a set $S$ consists of reduced words in the symbols of $S$ and
        their formal inverses; it has no relations beyond the group axioms. A
        free group on $n$ generators is denoted $F_n$.
    \end{definition}

    \begin{theorem}[Fundamental Theorem of Finitely Generated Abelian Groups]
        Every finitely generated abelian group is isomorphic to
        \[
            G
            \cong
            \Z^r\times
            \Z_{n_1}\times\cdots\times\Z_{n_k},
            \qquad
            n_1\mid n_2\mid\cdots\mid n_k,
        \]
        for uniquely determined $r\geq0$ and invariant factors
        $n_i\geq2$, up to isomorphism. The integer $r$ is the free rank.
        In particular, a finite abelian group has $r=0$ and can equivalently be
        written uniquely, up to reordering, as a direct product of cyclic groups
        of prime-power order.
    \end{theorem}

    \begin{definition}[Ring and Field]
        A ring is a set $R$ with operations $+$ and $\cdot$ s.t.
        $(R,+)$ is an abelian group, multiplication is associative, and
        multiplication distributes over addition on both sides. Throughout this
        book, rings are assumed to have a multiplicative identity $1_R$ unless
        stated otherwise. A field is a commutative ring in which $1_R\neq0_R$
        and every nonzero element has a multiplicative inverse.
    \end{definition}


    \begin{definition}[Unit, Zero Divisor, and Characteristic]
        In a ring $R$ with identity, a \emph{unit} is an element $u$ having a
        multiplicative inverse. A nonzero element $a$ is a \emph{zero divisor}
        if $ab=0$ for some nonzero $b$. The characteristic of $R$ is the least
        positive integer $n$ s.t.
        \[
            n1_R=0,
        \]
        if such an integer exists; otherwise the characteristic is $0$.
    \end{definition}

    \begin{definition}[Ideal, Ring Homomorphism, and Quotient Ring]
        An ideal $I$ of a commutative ring $R$ is an additive subgroup s.t.
        \[
            r\in R\land a\in I
            \quad\Longrightarrow\quad
            ra\in I.
        \]
        A ring homomorphism $\varphi\colon R\to S$ preserves addition,
        multiplication, and the multiplicative identity under the conventions
        of this book. Its kernel is an ideal. If $I$ is an ideal of $R$, the
        quotient ring $R/I$ consists of additive cosets with the induced ring
        operations.
    \end{definition}

    \begin{definition}[Integral Domain]
        An integral domain is a commutative ring with identity in which there are
        no zero divisors. In other words, if $ab=0$, then $a=0$ or $b=0$.
    \end{definition}

    \begin{theorem}[Every Finite Integral Domain Is a Field]
        Every finite integral domain is a field. Indeed, for each nonzero
        $a$, multiplication by $a$ is injective and hence, on a finite set,
        surjective; therefore some element maps to $1$.
    \end{theorem}

    \begin{definition}[Irreducible and Prime Elements]
        Let $R$ be an integral domain. A nonzero nonunit $p\in R$ is
        \emph{irreducible} if $p=ab$ implies that $a$ or $b$ is a unit. It is
        \emph{prime} if
        \[
            p\mid ab
            \quad\Longrightarrow\quad
            p\mid a\ \text{or}\ p\mid b.
        \]
        Every prime element is irreducible. In a UFD, the two notions coincide.
    \end{definition}

    \begin{definition}[UFD, PID, and Euclidean Domain]
        A unique factorization domain, abbreviated UFD, is an integral domain in
        which every nonzero nonunit factors into irreducibles and this
        factorization is unique up to order and multiplication by units.

        A principal ideal domain, abbreviated PID, is an integral domain in
        which every ideal has the form
        \[
            (a)=\{ra\mid r\in R\}
        \]
        for some $a\in R$.

        A Euclidean domain is an integral domain equipped with a function
        \[
            \delta\colon R\setminus\{0\}\to\Nzero
        \]
        s.t. for all $a,b\in R$ with $b\neq0$, there exist $q,r\in R$
        satisfying
        \[
            a=bq+r,
            \qquad
            r=0\lor \delta(r)<\delta(b).
        \]
    \end{definition}

    \begin{theorem}[Factorization-Domain Hierarchy]
        The standard implications are
        \[
            \text{field}
            \Longrightarrow
            \text{Euclidean domain}
            \Longrightarrow
            \text{PID}
            \Longrightarrow
            \text{UFD}
            \Longrightarrow
            \text{integral domain}.
        \]
        None of the reverse implications holds in general.
    \end{theorem}

    \begin{definition}[Quaternion Algebra]
        The quaternions form the real algebra
        \[
            \mathbb{H}
            =
            \{a+bi+cj+dk\mid a,b,c,d\in\R\},
        \]
        where
        \[
            i^2=j^2=k^2=ijk=-1.
        \]
        Quaternion multiplication is associative but not commutative; for
        example,
        \[
            ij=k,
            \qquad
            ji=-k.
        \]
        If
        \[
            q=a+bi+cj+dk,
        \]
        then
        \[
            \overline{q}=a-bi-cj-dk,
            \qquad
            |q|^2=q\overline{q}=a^2+b^2+c^2+d^2.
        \]
        Every nonzero quaternion has inverse
        \[
            q^{-1}=\frac{\overline{q}}{|q|^2}.
        \]
        Thus $\mathbb{H}$ is a noncommutative division ring, not a field under
        the convention that fields are commutative.
    \end{definition}

    \begin{formula}[Basic Factorization Identities]
        The most frequently used polynomial factorizations include
        \[
            a^2-b^2=(a-b)(a+b),
        \]
        \[
            a^3-b^3=(a-b)(a^2+ab+b^2),
            \qquad
            a^3+b^3=(a+b)(a^2-ab+b^2).
        \]
        More generally,
        \[
            a^n-b^n
            =
            (a-b)\sum_{k=0}^{n-1}a^{n-1-k}b^k,
        \]
        and, when $n$ is odd,
        \[
            a^n+b^n
            =
            (a+b)\sum_{k=0}^{n-1}(-1)^k a^{n-1-k}b^k.
        \]
    \end{formula}

    \begin{formula}[Cubic Factorization Identity]
        For all $a,b,c$,
        \[
            a^3+b^3+c^3-3abc
            =
            (a+b+c)(a^2+b^2+c^2-ab-bc-ca),
        \]
        or equivalently,
        \[
            a^3+b^3+c^3-3abc
            =
            \frac12(a+b+c)
            \left((a-b)^2+(b-c)^2+(c-a)^2\right).
        \]
        In particular, if $a+b+c=0$, then
        \[
            a^3+b^3+c^3=3abc.
        \]
    \end{formula}

    \begin{algorithm}[Polynomial Division Algorithm]
        Let $F$ be a field. If $f(x),g(x)\in F[x]$ and $g(x)\neq 0$, then there
        exist unique polynomials $q(x),r(x)\in F[x]$ s.t.
        \[
            f(x)=q(x)g(x)+r(x),
            \qquad
            r(x)=0\lor \deg r<\deg g.
        \]
    \end{algorithm}


    \begin{algorithm}[Euclidean Algorithm for Polynomials]
        Over a field $F$, repeatedly apply polynomial division
        \[
            f=qg+r,
            \qquad
            r=0\lor\deg r<\deg g,
        \]
        and replace $(f,g)$ by $(g,r)$. The last nonzero remainder, normalized
        to be monic, is $\gcd(f,g)$. Back-substitution expresses the gcd as
        \[
            af+bg=\gcd(f,g)
        \]
        for some $a,b\in F[x]$.
    \end{algorithm}

    \begin{theorem}[Factor Theorem]
        For a polynomial $f(x)\in F[x]$ and an element $a\in F$,
        \[
            f(a)=0
            \quad\Longleftrightarrow\quad
            x-a \text{ divides } f(x).
        \]
    \end{theorem}

    \begin{proposition}[Root-Counting Principle for Polynomials]
        A nonzero polynomial of degree $n$ over a field has at most $n$ roots.
        Hence two polynomials of degree at most $n$ that agree at $n+1$
        distinct points are identical.
    \end{proposition}

    \begin{formula}[Lagrange Interpolation Formula]
        For distinct $x_0,\dots,x_n$, the unique polynomial of degree at most
        $n$ satisfying $P(x_i)=y_i$ is
        \[
            P(x)
            =
            \sum_{i=0}^{n}y_i
            \prod_{\substack{0\leq j\leq n\\j\neq i}}
            \frac{x-x_j}{x_i-x_j}.
        \]
    \end{formula}

    \begin{formula}[Vi\`ete's Formulas]
        If
        \[
            a_nx^n+a_{n-1}x^{n-1}+\cdots+a_0
            =
            a_n\prod_{i=1}^{n}(x-r_i),
        \]
        where the roots are counted with multiplicity, then
        \[
            \sum_{i=1}^{n}r_i=-\frac{a_{n-1}}{a_n},
            \qquad
            \prod_{i=1}^{n}r_i=(-1)^n\frac{a_0}{a_n}.
        \]
        More generally, the elementary symmetric sums of the roots are the
        coefficients with alternating signs, divided by $a_n$.
    \end{formula}

    \begin{formula}[Newton's Identities; Newton--Girard Formulas]
        Let $e_k$ be the elementary symmetric polynomials in roots
        $r_1,\dots,r_n$, and let
        \[
            p_k=r_1^k+\cdots+r_n^k.
        \]
        With $e_0=1$, for $1\leq k\leq n$,
        \[
            p_k-e_1p_{k-1}+e_2p_{k-2}-\cdots
            +(-1)^{k-1}e_{k-1}p_1+(-1)^k k e_k=0.
        \]
    \end{formula}

    \begin{formula}[Polynomial Discriminant]
        If
        \[
            f(x)=a_n\prod_{i=1}^{n}(x-r_i),
        \]
        then
        \[
            \Disc(f)
            =
            a_n^{2n-2}\prod_{1\leq i<j\leq n}(r_i-r_j)^2.
        \]
        It vanishes if and only if $f$ has a repeated root.
    \end{formula}

    \begin{definition}[Algebraic Number, Minimal Polynomial, and Algebraic Conjugates]
        A complex number $\alpha$ is algebraic over $\Q$ if it is a root of a
        nonzero polynomial in $\Q[x]$. Its minimal polynomial
        $m_{\alpha,\Q}(x)$ is the unique monic irreducible polynomial in
        $\Q[x]$ having $\alpha$ as a root. A complex number that is not
        algebraic over $\Q$ is transcendental over $\Q$. The other roots of the
        minimal polynomial are the algebraic conjugates of $\alpha$ over $\Q$.
    \end{definition}

    \begin{theorem}[Conjugate Root Theorems]
        The standard conjugate-root principles are as follows.
        \begin{conditionlist}
            \item If $f\in\R[x]$ and $f(\alpha)=0$, then
            $f(\overline{\alpha})=0$. Hence non-real roots of a real-coefficient
            polynomial occur in complex-conjugate pairs, with the same
            multiplicity.
            \item If $f\in\Q[x]$, $a,b\in\Q$, and $d$ is a squarefree
            integer, then
            \[
                f(a+b\sqrt d)=0
                \quad\Longrightarrow\quad
                f(a-b\sqrt d)=0.
            \]
            Thus quadratic surd roots occur with their quadratic conjugates.
            \item More generally, if $f\in\Q[x]$ and $f(\alpha)=0$, then
            $m_{\alpha,\Q}$ divides $f$ in $\Q[x]$. Therefore every algebraic
            conjugate of $\alpha$ over $\Q$ is also a root of $f$. The same
            conclusion applies to integer-coefficient polynomials because
            $\Z[x]\subseteq\Q[x]$.
        \end{conditionlist}
    \end{theorem}

    \begin{theorem}[Rational Root Theorem]
        Let
        \[
            f(x)=a_nx^n+\cdots+a_0\in\Z[x].
        \]
        If a rational number $p/q$ in lowest terms is a root of $f$, then
        \[
            p\mid a_0,
            \qquad
            q\mid a_n.
        \]
    \end{theorem}


    \begin{theorem}[Cauchy Root Bound]
        If
        \[
            f(z)=a_nz^n+a_{n-1}z^{n-1}+\cdots+a_0,
            \qquad
            a_n\neq0,
        \]
        then every complex root $z$ of $f$ satisfies
        \[
            |z|
            \leq
            1+
            \max_{0\leq k<n}\left|\frac{a_k}{a_n}\right|.
        \]
        Thus all roots lie in an explicitly computable disk before any root is
        solved for exactly.
    \end{theorem}

    \begin{theorem}[Eisenstein's Irreducibility Criterion]
        Let
        \[
            f(x)=a_nx^n+\cdots+a_0\in\Z[x].
        \]
        If there exists a prime $p$ s.t.
        \[
            p\nmid a_n,
            \qquad
            p\mid a_i\quad(0\leq i<n),
            \qquad
            p^2\nmid a_0,
        \]
        then $f$ is irreducible over $\Q$.
    \end{theorem}

    \begin{theorem}[Fundamental Theorem of Symmetric Polynomials]
        Every symmetric polynomial can be expressed uniquely as a polynomial in
        the elementary symmetric polynomials.
    \end{theorem}

    \begin{theorem}[Combinatorial Nullstellensatz, Coefficient Form]
        Let $P\in F[x_1,\dots,x_n]$ have total degree
        $t_1+\cdots+t_n$, and suppose the coefficient of
        $x_1^{t_1}\cdots x_n^{t_n}$ is nonzero. If
        \[
            |S_i|>t_i
        \]
        for every finite set $S_i\subseteq F$, then some
        $(s_1,\dots,s_n)\in S_1\times\cdots\times S_n$ satisfies
        \[
            P(s_1,\dots,s_n)\neq0.
        \]
    \end{theorem}

    \begin{strategy}[Polynomial Method]
        Encode a forbidden configuration as zeros of a polynomial, then use a
        degree bound, interpolation, factorization, or a nonzero coefficient to
        prove that the polynomial cannot vanish everywhere.
    \end{strategy}

    \begin{theorem}[Fundamental Theorem of Algebra]
        Every nonconstant complex polynomial has at least one complex root.
        Equivalently, every polynomial of degree $n\geq 1$ over $\C$ splits into
        $n$ linear factors over $\C$, counted with multiplicity.
    \end{theorem}

    \begin{definition}[Algebraically Closed Field]
        A field $F$ is algebraically closed if every nonconstant polynomial in
        $F[x]$ has a root in $F$.
    \end{definition}

    \begin{definition}[Field Extension and Galois Group]
        If $F\subseteq K$ are fields, then $K/F$ is a field extension and its
        degree is
        \[
            [K:F]=\dim_F K.
        \]
        The extension is finite if this dimension is finite. The group
        \[
            \Gal(K/F)
            =
            \{\sigma\in\Aut(K)\mid \sigma(a)=a\ \forall a\in F\}
        \]
        is the Galois group. For a finite Galois extension, intermediate fields
        and subgroups of the Galois group are related by the Galois
        correspondence.
    \end{definition}

    \begin{theorem}[Fundamental Theorem of Galois Theory; Galois Correspondence]
        For a finite Galois extension $L/K$, there is an inclusion-reversing
        correspondence between intermediate fields $K\subseteq E\subseteq L$ and
        subgroups $H\leq \Gal(L/K)$, given by
        \[
            E\longmapsto \Gal(L/E),
            \qquad
            H\longmapsto L^H.
        \]
    \end{theorem}

    \begin{theorem}[Abel--Ruffini Theorem]
        There is no formula for the roots of a general polynomial of degree $5$
        or higher using only the field operations and radicals.
    \end{theorem}

    \begin{remark}
        The Abel--Ruffini theorem does not say that every quintic equation is
        unsolvable. It says that no universal radical formula exists for the
        general quintic.
    \end{remark}

    \chapter{Topology}

    \relatedcourses{(MATH321) General Topology}

    \relatedbooks{Topology (James R. Munkres)}

    \begin{definition}[Topology]
        A topology on a set $X$ is a collection $\mathcal{T}$ of subsets of $X$
        whose elements are called open sets, satisfying:
        \begin{conditionlist}
            \item $\varnothing,X\in\mathcal{T}$;
            \item arbitrary unions of sets in $\mathcal{T}$ belong to
            $\mathcal{T}$;
            \item finite intersections of sets in $\mathcal{T}$ belong to
            $\mathcal{T}$.
        \end{conditionlist}
    \end{definition}


    \begin{definition}[Basis for a Topology]
        A collection $\mathcal{B}$ of subsets of $X$ is a basis if every
        $x\in X$ lies in some $B\in\mathcal{B}$ and, whenever
        $x\in B_1\cap B_2$ with $B_1,B_2\in\mathcal{B}$, there exists
        $B_3\in\mathcal{B}$ s.t.
        \[
            x\in B_3\subseteq B_1\cap B_2.
        \]
        The topology generated by $\mathcal{B}$ consists of arbitrary unions
        of basis elements. Open balls form the standard basis for the metric
        topology of a metric space.
    \end{definition}


    \begin{definition}[Open Sets, Closed Sets, and Neighborhoods]
        If $(X,\mathcal{T})$ is a topological space, a subset $U\subseteq X$ is
        open if $U\in\mathcal{T}$. A subset $F\subseteq X$ is closed if
        $X\setminus F\in\mathcal{T}$. A neighborhood of $x\in X$ is a set that
        contains an open set containing $x$. A local basis at $x$ is a family
        of neighborhoods with the property that every neighborhood of $x$
        contains one of them.
    \end{definition}

    \begin{definition}[First- and Second-Countability]
        A topological space is first-countable if every point has a countable
        local neighborhood basis. It is second-countable if its topology has a
        countable basis. Every second-countable space is first-countable, but
        not conversely in general.
    \end{definition}

    \begin{definition}[Compactness]
        An open cover of $X$ is a family of open sets whose union contains
        $X$. A topological space $X$ is compact if every open cover has a finite
        subcover. Equivalently, for every family $\mathcal{U}\subseteq\mathcal{T}$,
        \[
            X\subseteq\bigcup_{U\in\mathcal{U}}U
            \Rightarrow
            \exists U_1,\dots,U_n\in\mathcal{U}
            \ \text{s.t.}\
            X\subseteq U_1\cup\cdots\cup U_n.
        \]
    \end{definition}

    \begin{definition}[Connectedness]
        A topological space $X$ is connected if it cannot be written as the union
        of two disjoint nonempty open sets. Equivalently, the only subsets of
        $X$ that are both open and closed are $\varnothing$ and $X$.
    \end{definition}


    \begin{definition}[Continuity in Topological Spaces]
        A map $f\colon (X,\mathcal{T}_X)\to(Y,\mathcal{T}_Y)$ is continuous if the
        inverse image of every open set is open:
        \[
            \forall V\in\mathcal{T}_Y,
            \qquad
            f^{-1}(V)\in\mathcal{T}_X.
        \]
    \end{definition}

    \begin{definition}[Homeomorphism]
        A homeomorphism is a bijective continuous map whose inverse is also
        continuous. Equivalently, $f\colon X\xrightarrow{\sim}Y$ is a homeomorphism
        if $f$ and $f^{-1}$ are continuous. Two spaces are homeomorphic, written
        $X\cong Y$, if such a map exists.
    \end{definition}


    \begin{theorem}[Basic Topological Invariants under Continuous Maps]
        A continuous image of a compact space is compact, and a continuous
        image of a connected space is connected. Consequently, compactness and
        connectedness are preserved by homeomorphism.
    \end{theorem}

    \begin{definition}[Hausdorff Space]
        A topological space $X$ is Hausdorff if distinct points can be separated
        by disjoint open neighborhoods:
        \[
            x\neq y
            \quad\Longrightarrow\quad
            \exists U,V\text{ open s.t. }x\in U,\ y\in V,\ U\cap V=\varnothing.
        \]
        Every metric space is Hausdorff.
    \end{definition}

    \begin{theorem}[Compact Subsets of Hausdorff Spaces]
        Every compact subset of a Hausdorff space is closed. In particular, a
        continuous bijection from a compact space onto a Hausdorff space is
        automatically a homeomorphism.
    \end{theorem}

    \begin{definition}[Path, Path-Connectedness, and Homotopy]
        A path from $x_0$ to $x_1$ in $X$ is a continuous map
        \[
            \gamma\colon [0,1]\to X,
            \qquad
            \gamma(0)=x_0,
            \quad
            \gamma(1)=x_1.
        \]
        The space $X$ is path-connected if every two points can be joined by a
        path. Two maps $f,g\colon X\to Y$ are homotopic, written $f\simeq g$, if
        there exists a continuous map
        \[
            H\colon X\times[0,1]\to Y
        \]
        with $H(x,0)=f(x)$ and $H(x,1)=g(x)$. A loop is a path with the same
        initial and terminal point.
    \end{definition}

    \begin{theorem}[Path-Connectedness Implies Connectedness]
        Every path-connected space is connected. The converse is false in
        general, so connectedness and path-connectedness should not be
        interchanged without additional hypotheses.
    \end{theorem}

    \begin{theorem}[Jordan Curve Theorem]
        Every simple closed curve in the plane separates the plane into exactly
        two connected components, an interior and an exterior, and is the
        common boundary of both.
    \end{theorem}

    \begin{theorem}[Brouwer Fixed-Point Theorem]
        Every continuous map from a closed ball in $\R^n$ to itself has at least
        one fixed point.
    \end{theorem}

    \begin{theorem}[Borsuk--Ulam Theorem]
        Every continuous map
        \[
            f\colon S^n\to\R^n
        \]
        takes some pair of antipodal points to the same value.
    \end{theorem}

    \begin{corollary}[Ham Sandwich Theorem]
        Any $n$ bounded measurable subsets of $\R^n$ can be bisected
        simultaneously by one hyperplane.
    \end{corollary}

    \begin{theorem}[Hairy Ball Theorem]
        Every continuous tangent vector field on an even-dimensional sphere
        $S^{2m}$ vanishes at some point.
    \end{theorem}

    \begin{theorem}[Invariance of Domain]
        If $U\subseteq\R^n$ is open and $f\colon U\to\R^n$ is continuous and injective,
        then $f(U)$ is open and $f$ is a homeomorphism from $U$ onto $f(U)$.
    \end{theorem}

    \begin{definition}[Manifold]
        An $n$-dimensional topological manifold is a Hausdorff, second-countable
        topological space in which every point has a neighborhood homeomorphic
        to an open subset of $\R^n$.
    \end{definition}

    \begin{definition}[Fundamental Group]
        Fix a base point $x_0\in X$. The fundamental group
        $\pi_1(X,x_0)$ is the set of endpoint-preserving homotopy classes of
        loops based at $x_0$, with multiplication induced by concatenation of
        loops. It is a basic algebraic invariant of a topological space. A
        path-connected space is simply connected if
        \[
            \pi_1(X,x_0)\cong\{e\},
        \]
        equivalently, every loop can be contracted to a point.
    \end{definition}

    \begin{fact}[Fundamental Groups of Familiar Spaces]
        The following basic examples are worth recognizing:
        \[
            \pi_1(S^1)\cong\Z,
            \qquad
            \pi_1(S^n)\cong\{e\}\quad(n\geq 2),
            \qquad
            \pi_1(\mathbb{T}^2)\cong\Z^2.
        \]
    \end{fact}

    \begin{strategy}
        For introductory topology questions, the most common quick checks are:
        compactness means every open cover has a finite subcover, connectedness
        means the space cannot be separated into two disjoint nonempty open
        sets, and a homeomorphism is a continuous bijection with continuous
        inverse.
    \end{strategy}

    \begin{formula}[Euler Characteristic of a Surface]
        For a triangulated closed surface, the Euler characteristic is
        \[
            \chi=V-E+F.
        \]
        For an orientable closed surface of genus $g$,
        \[
            \chi=2-2g.
        \]
    \end{formula}

    \begin{theorem}[Classification of Closed Surfaces]
        Every connected closed surface is homeomorphic to either an orientable
        surface of genus $g\geq 0$ or a nonorientable surface obtained as a
        connected sum of $k\geq 1$ projective planes.
    \end{theorem}

    \begin{theorem}[Poincaré Conjecture]
        Every closed simply connected $3$-manifold is homeomorphic to the
        $3$-sphere $S^3$.
    \end{theorem}

    \begin{remark}
        The Poincaré conjecture was proved by Grigori Perelman using Ricci flow.
        The fundamental correspondence is
        \[
            \text{Poincaré conjecture} \longleftrightarrow
            \text{Perelman} \longleftrightarrow
            \text{Ricci flow}.
        \]
    \end{remark}

    \begin{theorem}[Thurston Geometrization Theorem]
        Every closed orientable $3$-manifold can be decomposed along embedded
        $2$-spheres and incompressible tori into pieces whose interiors admit
        geometric structures modeled on one of the eight Thurston geometries:
        \[
            S^3,\ \mathbb{E}^3,\ \mathbb{H}^3,\ S^2\times\R,\ \mathbb{H}^2\times\R,
            \ \widetilde{\SL}_2(\R),\ \mathrm{Nil},\ \mathrm{Sol}.
        \]
        The sphere decomposition is the prime decomposition, and the torus
        decomposition is the JSJ decomposition of the irreducible pieces.
    \end{theorem}

    \begin{fact}[The Eight Thurston Geometries]
        The eight model geometries are
        \[
            S^3,
            \quad \mathbb{E}^3,
            \quad \mathbb{H}^3,
            \quad S^2\times\R,
            \quad \mathbb{H}^2\times\R,
            \quad \widetilde{\SL}_2(\R),
            \quad \mathrm{Nil},
            \quad \mathrm{Sol}.
        \]
    \end{fact}

    \chapter{Miscellaneous Topics}

    \relatedcourses{}

    \relatedbooks{}

    \begin{definition}[Category]
        A category $\mathcal{C}$ consists of objects and morphisms between
        objects. For objects $A,B,C$, morphisms can be composed,
        \[
            A\xrightarrow{f}B\xrightarrow{g}C
            \quad\longmapsto\quad
            g\circ f\colon A\to C,
        \]
        composition is associative, and every object $A$ has an identity
        morphism $\id_A$. For every morphism $f\colon A\to B$, the identity laws are
        \[
            f\circ\id_A=f,
            \qquad
            \id_B\circ f=f.
        \]
        Category theory studies mathematical structures by emphasizing objects,
        structure-preserving maps, and their composition.
    \end{definition}

    \begin{definition}[Nash Equilibrium]
        In a strategic game, a strategy profile is a Nash equilibrium if no
        player can improve their payoff by changing only their own strategy
        while every other player's strategy is held fixed. Equivalently, every
        player's chosen strategy is a best response to the strategies chosen by
        the others.
    \end{definition}

    \begin{definition}[Self-Similar Dimension]
        Suppose a self-similar set is the union of $N$ geometrically similar
        copies of itself, each scaled by a factor $r$ with $0<r<1$, and the
        usual separation assumptions hold. Its similarity dimension $d$ is
        determined by
        \[
            Nr^d=1,
            \qquad
            d=\frac{\ln N}{\ln(1/r)}.
        \]
        This gives a standard noninteger dimension for many classical fractals.
    \end{definition}

    \begin{fact}[Menger Sponge]
        At each stage of the Menger sponge construction, a cube is replaced by
        $20$ copies scaled by $1/3$. Hence its similarity dimension is
        \[
            \boxed{\frac{\ln20}{\ln3}}.
        \]
    \end{fact}

    \begin{strategy}[Symmetry and Without Loss of Generality]
        Before splitting into cases, identify transformations that preserve the
        problem. If several cases differ only by a permutation, reflection,
        sign change, or cyclic relabeling, choose one representative and state
        explicitly why the remaining cases follow by symmetry. ``Without loss
        of generality'' is justified only after this invariance has been
        identified.
    \end{strategy}

    \begin{strategy}[Parity and Modular Obstructions]
        Before searching for a construction or integer solution, reduce the
        relevant quantities modulo a small modulus. Parity, residues of powers,
        and invariant congruence classes often rule out entire families of
        possibilities with almost no calculation.
    \end{strategy}

    \begin{strategy}[Telescoping and Local Cancellation]
        For a sum or product indexed consecutively, try to rewrite each term so
        that neighboring factors cancel. Typical forms are
        \[
            a_k=b_k-b_{k+1}
            \qquad\text{or}\qquad
            a_k=\frac{b_{k+1}}{b_k}.
        \]
        Then
        \[
            \sum_{k=m}^{n}a_k=b_m-b_{n+1},
            \qquad
            \prod_{k=m}^{n}a_k=\frac{b_{n+1}}{b_m}.
        \]
    \end{strategy}

    \begin{strategy}[Descent and Minimal Counterexample]
        If a statement about positive integers or finite objects is suspected to
        hold universally, assume a counterexample exists and choose one minimal
        under a natural size parameter. Use the defining relations to construct
        a strictly smaller counterexample. This contradiction is the abstract
        form of infinite descent and is useful far beyond number theory.
    \end{strategy}

    \begin{strategy}[Complex Numbers as Plane Geometry]
        Identify the Euclidean plane with $\C$ when rotations, similarities,
        roots of unity, or cyclic configurations are prominent. Multiplication
        by $e^{i\theta}$ is rotation by $\theta$, multiplication by a positive
        scalar is dilation, and complex conjugation is reflection across the
        real axis. Translate back to ordinary geometry only after the algebraic
        relation has been simplified.
    \end{strategy}

    \begin{strategy}[Polynomial Encoding]
        Encode a finite set of numbers as roots of a polynomial, or encode a
        sequence by a generating polynomial or generating function, when sums,
        products, symmetric expressions, or recurrence relations are involved.
        Root--coefficient relations, interpolation, and coefficient extraction
        can then replace case-by-case manipulation.
    \end{strategy}

    \begin{tactic}[Floor and Fractional-Part Identities]
        For every real $x$,
        \[
            x=\lfloor x\rfloor+\{x\},
            \qquad
            0\leq\{x\}<1.
        \]
        Moreover,
        \[
            \lfloor x\rfloor+\lfloor-x\rfloor
            =
            \begin{cases}
                0, & x\in\Z,\\
                -1, & x\notin\Z,
            \end{cases}
        \]
        and for $n\in\Z$,
        \[
            \lfloor x+n\rfloor=\lfloor x\rfloor+n.
        \]
        These identities convert many apparently analytic floor problems into
        integer casework.
    \end{tactic}

    \begin{strategy}[Iterated Knowledge Elimination]
        Let $\Omega_0$ be a finite set of possible states. For each agent $i$ and
        state $\omega\in\Omega$, let
        \[
            \mathcal{I}_i^\Omega(\omega)
            =
            \{\omega'\in\Omega\mid
              \omega'\text{ gives agent }i\text{ the same information as }\omega\}.
        \]
        Relative to the current state space $\Omega$, agent $i$ knows the state
        precisely when
        \[
            \bigl|\mathcal{I}_i^\Omega(\omega)\bigr|=1.
        \]
        A truthful public announcement restricts $\Omega$ to exactly the states
        at which that announcement is true. Repeating this restriction models
        successive statements about knowledge or ignorance. In particular, the
        statement ``I knew that $Q$'' means
        \[
            \forall\omega'\in\mathcal{I}_i^\Omega(\omega),
            \qquad
            Q(\omega').
        \]
    \end{strategy}

    \begin{strategy}[Transform the Given Objects]
        Replace the original configuration by an equivalent one whose structure
        is easier to see: partition a region, introduce prefix sums, encode a
        movement by vectors, pass to residues, diagonalize a matrix, or complete
        a partial object. The transformed problem should preserve the quantity or
        property being asked about.
    \end{strategy}

    \begin{strategy}[Invariant, Monovariant, and Extremal Principles]
        An invariant remains unchanged under every allowed move. A monovariant
        moves strictly in one direction and therefore prevents cycles or proves
        termination. The extremal principle chooses a largest, smallest, first,
        or last object and exploits the extra restrictions forced by extremality.
        These are among the most common contest ideas hidden behind elementary
        statements.
    \end{strategy}

    \begin{strategy}[Complete the Configuration and Work Backward]
        First imagine a completed solution, then determine which parts are forced
        and which choices remain free. This is effective in geometric
        constructions, permutation completions, lattice paths, tilings, and
        assignment problems. Fixing one object under a symmetry often removes
        redundant cases before the remaining count begins.
    \end{strategy}

    \begin{strategy}[Rotate, Reflect, and Superimpose]
        A rotation, reflection, or translated copy can convert a broken path into
        a straight one, expose equal vectors, pair terms, or turn a local pattern
        into a global invariant. When two configurations overlap, compare what
        cancels and what remains rather than recomputing each separately.
    \end{strategy}

    \begin{strategy}[Homogenization and Normalization]
        If an inequality is homogeneous, scale the variables to impose a simple
        condition such as $a+b+c=1$, $abc=1$, or $x^2+y^2=1$.
    \end{strategy}

    \begin{strategy}[Equality Cases and Smoothing]
        Identify the expected equality case before selecting an inequality.
        For a symmetric expression, replacing two variables by their average
        often reveals the extremal configuration.
    \end{strategy}

    \begin{theorem}[Triangle and Reverse Triangle Inequalities]
        For real or complex numbers $x,y$,
        \[
            \bigl||x|-|y|\bigr|
            \leq
            |x-y|
            \leq
            |x|+|y|.
        \]
        More generally, every norm satisfies
        \[
            \|\mathbf{x}+\mathbf{y}\|
            \leq
            \|\mathbf{x}\|+\|\mathbf{y}\|.
        \]
    \end{theorem}

    \begin{formula}[Basic Sum-of-Squares Identities]
        For real $a,b,c$,
        \[
            a^2+b^2+c^2-ab-bc-ca
            =
            \frac12\left((a-b)^2+(b-c)^2+(c-a)^2\right)
            \geq0.
        \]
        Consequently,
        \[
            a^2+b^2+c^2\geq ab+bc+ca
        \]
        and
        \[
            (a+b+c)^2\leq3(a^2+b^2+c^2).
        \]
    \end{formula}

    \begin{formula}[Frequently Used Two-Variable Inequalities]
        For real $x,y$,
        \[
            x^2+y^2\geq2xy.
        \]
        For positive $x,y$,
        \[
            \frac{x}{y}+\frac{y}{x}\geq2,
            \qquad
            x+\frac1x\geq2.
        \]
        Equality holds exactly when the two compared positive quantities are
        equal.
    \end{formula}

    \begin{theorem}[Lagrange's Identity]
        For real numbers $a_1,\dots,a_n$ and $b_1,\dots,b_n$,
        \[
            \left(\sum_{i=1}^{n}a_i^2\right)
            \left(\sum_{i=1}^{n}b_i^2\right)
            -
            \left(\sum_{i=1}^{n}a_ib_i\right)^2
            =
            \sum_{1\leq i<j\leq n}(a_ib_j-a_jb_i)^2.
        \]
        The nonnegativity of the right-hand side gives the
        Cauchy--Schwarz inequality immediately.
    \end{theorem}

    \begin{theorem}[QM--AM--GM--HM Inequalities]
        For positive $x_1,\dots,x_n$,
        \[
            \sqrt{\frac{x_1^2+\cdots+x_n^2}{n}}
            \geq
            \frac{x_1+\cdots+x_n}{n}
            \geq
            \sqrt[n]{x_1\cdots x_n}
            \geq
            \frac{n}{\frac1{x_1}+\cdots+\frac1{x_n}}.
        \]
        Equality holds if and only if all variables are equal.
    \end{theorem}

    \begin{theorem}[Weighted AM--GM Inequality]
        If $w_i\geq0$, $\sum_{i=1}^{n}w_i=1$, and $x_i>0$, then
        \[
            \sum_{i=1}^{n}w_ix_i
            \geq
            \prod_{i=1}^{n}x_i^{w_i}.
        \]
        Equality holds when all $x_i$ with positive weight are equal.
    \end{theorem}

    \begin{theorem}[Power Mean Inequality]
        For positive $x_1,\dots,x_n$ and real $r\neq0$, define
        \[
            M_r
            =
            \left(
                \frac{x_1^r+\cdots+x_n^r}{n}
            \right)^{1/r},
        \]
        and define $M_0=(x_1\cdots x_n)^{1/n}$. If $r<s$, then
        \[
            M_r\leq M_s.
        \]
        The harmonic, geometric, arithmetic, and quadratic means correspond to
        $r=-1,0,1,2$, respectively.
    \end{theorem}

    \begin{theorem}[Young's Inequality]
        If $x,y\geq0$, $p,q>1$, and
        \[
            \frac1p+\frac1q=1,
        \]
        then
        \[
            xy\leq\frac{x^p}{p}+\frac{y^q}{q}.
        \]
    \end{theorem}

    \begin{theorem}[Cauchy--Schwarz Inequality, Engel Form]
        For real $a_i$ and positive $b_i$,
        \[
            \sum_{i=1}^{n}\frac{a_i^2}{b_i}
            \geq
            \frac{(a_1+\cdots+a_n)^2}{b_1+\cdots+b_n}.
        \]
    \end{theorem}

    \begin{theorem}[H\"older's Inequality]
        If $p,q>1$ and $1/p+1/q=1$, then
        \[
            \sum_{i=1}^{n}|a_ib_i|
            \leq
            \left(\sum_{i=1}^{n}|a_i|^p\right)^{1/p}
            \left(\sum_{i=1}^{n}|b_i|^q\right)^{1/q}.
        \]
    \end{theorem}

    \begin{theorem}[Minkowski's Inequality]
        For $p\geq1$,
        \[
            \left(\sum_{i=1}^{n}|a_i+b_i|^p\right)^{1/p}
            \leq
            \left(\sum_{i=1}^{n}|a_i|^p\right)^{1/p}
            +
            \left(\sum_{i=1}^{n}|b_i|^p\right)^{1/p}.
        \]
    \end{theorem}

    \begin{theorem}[Rearrangement Inequality]
        If
        \[
            a_1\leq\cdots\leq a_n,
            \qquad
            b_1\leq\cdots\leq b_n,
        \]
        then for every permutation $\sigma\in S_n$,
        \[
            \sum_{i=1}^{n}a_ib_{n+1-i}
            \leq
            \sum_{i=1}^{n}a_ib_{\sigma(i)}
            \leq
            \sum_{i=1}^{n}a_ib_i.
        \]
    \end{theorem}

    \begin{theorem}[Chebyshev's Sum Inequality]
        If $a_1\leq\cdots\leq a_n$ and $b_1\leq\cdots\leq b_n$, then
        \[
            \frac1n\sum_{i=1}^{n}a_ib_i
            \geq
            \left(\frac1n\sum_{i=1}^{n}a_i\right)
            \left(\frac1n\sum_{i=1}^{n}b_i\right).
        \]
    \end{theorem}

    \begin{theorem}[Schur's Inequality]
        For $a,b,c\geq0$ and $r\geq0$,
        \[
            \sum_{\mathrm{cyc}}a^r(a-b)(a-c)\geq0.
        \]
    \end{theorem}

    \begin{formula}[Nesbitt's Inequality]
        For positive $a,b,c$,
        \[
            \frac{a}{b+c}+\frac{b}{c+a}+\frac{c}{a+b}
            \geq
            \frac32.
        \]
    \end{formula}

    \begin{theorem}[Bernoulli's Inequality]
        If $x\geq-1$ and $r\geq1$, then
        \[
            (1+x)^r\geq1+rx.
        \]
    \end{theorem}

    \begin{strategy}[Comparing Powers by Logarithms]
        To compare two positive numbers such as $a^b$ and $b^a$, take logarithms:
        \[
            a^b>b^a
            \quad\Longleftrightarrow\quad
            b\ln a>a\ln b
            \quad\Longleftrightarrow\quad
            \frac{\ln a}{a}>\frac{\ln b}{b}.
        \]
        The function $\ln x/x$ is decreasing for $x>e$. This quickly handles
        comparisons such as $11^9$ versus $9^{11}$ or $e^\pi$ versus $\pi^e$.
    \end{strategy}

    \begin{formula}[More Useful Series Values]
        The following values are common in fast calculation problems:
        \[
            \sum_{n=1}^{\infty}\frac{1}{n(n+1)}=1,
            \qquad
            \sum_{n=1}^{\infty}\frac{1}{n(n+2)}=\frac34,
        \]
        \[
            \sum_{n=1}^{\infty}\frac{1}{(2n-1)^4}
            =
            \left(1-\frac{1}{16}\right)\frac{\pi^4}{90}
            =
            \frac{\pi^4}{96},
        \]
        where the value of $\sum_{n=1}^{\infty}n^{-4}$ is recorded with the
        classical series values in the Calculus chapter.
    \end{formula}

    \begin{strategy}
        These values often appear after a simple transformation. For example,
        \[
            \frac12-\frac14+\frac16-\frac18+\cdots
            =
            \frac12
            \left(
                1-\frac12+\frac13-\frac14+\cdots
            \right)
            =
            \frac{\ln 2}{2}.
        \]
    \end{strategy}
